% Leçon 28 — Grammaire : 把字句 — Chinois 1ère A — Cours signé OnBuch+
% Programme officiel MINESEC. Compiler avec : tectonic ZHA28-grammaire.tex
\newcommand{\DOCMATIERE}{Chinois}
\newcommand{\DOCNIVEAU}{1ère A}
\newcommand{\DOCMODULE}{Module 3 : Citoyenneté et ouverture au monde}
\newcommand{\DOCLECON}{ZHA28}
\newcommand{\DOCTITRE}{Leçon 28 — Grammaire : 把字句}
% OnBuch+ — préambule « cours signés » ENRICHI (compilé avec tectonic / XeLaTeX)
% Système visuel « Pop » d'OnBuch+ : fond crème (jamais blanc), bordures encre
% épaisses, ombres « dures » décalées, titres Archivo Black en capitales.
% Version MATHÉMATIQUES : graphes (pgfplots/tikz), boîtes pédagogiques
% (définition, propriété, méthode, attention, le savais-tu, exercice résolu,
% exercice, TP, bilan), page de garde et signature OnBuch+ ancrée en bas.
%
% Macros attendues dans main.tex AVANT \input{preamble} :
%   \DOCMATIERE  \DOCNIVEAU  \DOCMODULE  \DOCLECON (numéro)  \DOCTITRE
\documentclass[11pt]{article}

\usepackage{fontspec}
\usepackage[french]{babel} % français = langue principale ; le chinois via \zh{..} / textebox (xeCJK)
\usepackage{xeCJK}
\usepackage{pifont} % \ding{51} \ding{55} utilisés par les modèles
\usepackage[a4paper,margin=2cm,top=3.4cm,bottom=2.8cm,headheight=1.4cm,footskip=1.3cm]{geometry}
\usepackage{amsmath,amssymb,amsfonts}
\usepackage{mathtools} % \xmapsto, \coloneqq... souvent utilisés par les modèles
\usepackage{cancel} % simplifications barrées (bilans de forces)
% \abs{z} (module, valeur absolue) et \norm{u} (norme), utilisés par les modèles
\providecommand{\abs}{}\let\abs\relax\DeclarePairedDelimiter{\abs}{\lvert}{\rvert}
\providecommand{\norm}{}\let\norm\relax\DeclarePairedDelimiter{\norm}{\lVert}{\rVert}
\usepackage[table]{xcolor} % [table] : fournit aussi \rowcolors (alternance de lignes)
\usepackage{pagecolor}
\usepackage{tikz}
\usetikzlibrary{arrows.meta,positioning,calc,shapes.geometric,shapes.misc,shapes.arrows,decorations.pathmorphing,decorations.pathreplacing,decorations.markings,patterns,shadows,mindmap,trees,fit,backgrounds,matrix,intersections,angles,quotes}
\usepackage{pgfplots}
\usepackage[version=4]{mhchem} % équations nucléaires \ce{^{235}_{92}U}
\usepackage[european,siunitx]{circuitikz} % schémas électriques
\pgfplotsset{compat=1.18}
\usepgfplotslibrary{fillbetween,groupplots} % aires entre courbes (calcul intégral)
\usepackage{enumitem}
\usepackage{fancyhdr}
% [most] charge toutes les bibliothèques tcolorbox (listings, minted, raster,
% documentation...) et gonfle inutilement la mémoire TeX sur un document long
% (~300 boîtes) : on ne charge que ce qui est réellement utilisé.
\usepackage[skins,breakable]{tcolorbox}
\usepackage{graphicx}
\usepackage{colortbl}
\usepackage{booktabs}
\usepackage{tabularx}
\usepackage{diagbox} % case d'en-tête diagonale des tableaux à double entrée
% Colonnes extensibles centrée / gauche / droite (C, L, R), souvent utilisées par les modèles
\newcolumntype{C}{>{\centering\arraybackslash}X}
\newcolumntype{L}{>{\raggedright\arraybackslash}X}
\newcolumntype{R}{>{\raggedleft\arraybackslash}X}
\usepackage{array}
\usepackage{multicol}
\usepackage{multirow}
\usepackage{siunitx}
% parse-numbers=false : accepte aussi \SI{2,5 \times 10^{-3}}{...}
\sisetup{locale=FR,output-decimal-marker={,},group-minimum-digits=4,parse-numbers=false}
\DeclareSIUnit\ohm{\ensuremath{\symup{\Omega}}} % U+2126 absent de la police
\DeclareSIUnit\division{div} % réglages d'oscilloscope (ms/div, V/div)
\DeclareSIUnit\tr{tr} % tours (vitesse de rotation, ex. \SI{1500}{\tr\per\minute})
\DeclareSIUnit\molar{M} % concentration molaire (ex. \SI{3}{\molar}, \SI{1}{\micro\molar})
\DeclareSIUnit\an{an} % année (le modèle confond parfois avec \year, un compteur TeX) — ex. \SI{5}{\kilo\meter\per\an}
\DeclareSIUnit\FCFA{FCFA} % franc CFA (ex. \SI{500}{\FCFA\per\kilo\gram})
\DeclareSIUnit\degreeBrix{\textdegree Brix} % teneur en sucre d'un sirop/jus (réfractométrie)
\DeclareSIUnit\month{mois} % le modèle utilise parfois \month, un compteur TeX, comme unité de durée
\usepackage{qrcode}
% \degree : commande fréquemment utilisée par les modèles pour un angle ou une
% température en dehors de \SI{}{} (non fournie par les paquets chargés).
\providecommand{\degree}{^{\circ}}
% \qed (fin de démonstration), utilisé par les modèles sans amsthm
\providecommand{\qed}{\hfill$\square$}
% \barre : double barre des valeurs interdites dans un tableau de variations (tkz-tab)
\providecommand{\barre}{\ensuremath{\|}}
% environnement proof (amsthm non chargé) utilisé par les modèles
\ifdefined\proof\else\newenvironment{proof}[1][Démonstration]{\par\noindent\textit{#1.}\ }{\qed\par}\fi
\usepackage{lastpage}
\usepackage{hyperref}
\hypersetup{hidelinks}

% ============================================================
% Palette « Pop » OnBuch+ — mêmes hex que la classe P (plus_ui.dart)
% ============================================================
\definecolor{poporange}{HTML}{F59321}
\definecolor{popdark}{HTML}{E07A0C}
\definecolor{popcream}{HTML}{FFF6E8}
\definecolor{popink}{HTML}{1C1714}
\definecolor{poppurple}{HTML}{7A5AE0}
\definecolor{popgreen}{HTML}{2E9E6A}
\definecolor{poppink}{HTML}{E0457B}
\definecolor{popblue}{HTML}{2F7DE1}
\definecolor{popgold}{HTML}{C98A12}
\definecolor{popmuted}{HTML}{8A7F76}
\definecolor{popline}{HTML}{EADFD2}
\definecolor{popgray}{HTML}{8A7F76}
\colorlet{popgrey}{popgray}
% Alias tolérants (le modèle nomme parfois les couleurs différemment)
\colorlet{popwhite}{white}
\colorlet{popblack}{popink}
\colorlet{popred}{poppink}
\colorlet{popyellow}{popgold}
% Teintes claires (fonds de tableaux/encadrés) — le modèle nomme parfois
% les couleurs avec un suffixe L (ex. popblueL) sans les avoir définies.
\colorlet{poporangeL}{poporange!20}
\colorlet{popdarkL}{popdark!20}
\colorlet{poppurpleL}{poppurple!20}
\colorlet{popgreenL}{popgreen!20}
\colorlet{poppinkL}{poppink!20}
\colorlet{popblueL}{popblue!20}
\colorlet{popgoldL}{popgold!20}
\colorlet{popredL}{popred!20}
\colorlet{popgrayL}{popgray!20}
% Teintes claires pour fonds de tableaux / zones
\colorlet{popblueL}{popblue!10!white}
\colorlet{poporangeL}{poporange!14!white}
\colorlet{popgreenL}{popgreen!12!white}
\colorlet{poppurpleL}{poppurple!12!white}
\colorlet{poppinkL}{poppink!12!white}
\colorlet{popgoldL}{popgold!14!white}

\pagecolor{popcream}

% ============================================================
% Typographie
% ============================================================
\defaultfontfeatures{Ligatures=TeX}
\setmainfont{PlusJakartaSans-Medium.ttf}[Path=../../../fonts/,BoldFont=PlusJakartaSans-Bold.ttf]
% Symboles, API et emojis absents de la police principale : repli sur DejaVu Sans (caractères actifs)
\newfontface\symfont{DejaVuSans.ttf}[Path=../../../fonts/]
\makeatletter
\newcommand\symactive[1]{\catcode#1=13 \begingroup\lccode`\~=#1 \lowercase{\endgroup\def~}{{\symfont\char#1}}}
\newcommand\symblank[1]{\catcode#1=13 \begingroup\lccode`\~=#1 \lowercase{\endgroup\def~}{}}
\newcommand\symhyph[1]{\catcode#1=13 \begingroup\lccode`\~=#1 \lowercase{\endgroup\def~}{\mbox{-}}}
\newcount\symc
\newcommand\symrange[2]{\symc=#1 \@whilenum\symc<#2 \do{\expandafter\symactive\expandafter{\the\symc}\advance\symc by 1 }}
\newcommand\symblankrange[2]{\symc=#1 \@whilenum\symc<#2 \do{\expandafter\symblank\expandafter{\the\symc}\advance\symc by 1 }}
\makeatother
\symrange{"0250}{"02FF}\symactive{"2713}\symactive{"2714}\symactive{"2717}\symactive{"2718}\symactive{"2610}\symactive{"2611}\symactive{"2612}
\symrange{"2600}{"27BF}
\catcode"2705=13 \begingroup\lccode`\~="2705 \lowercase{\endgroup\def~}{{\symfont\char"2713}}\symblank{"26BD}\symblank{"26A1}
\symrange{"2460}{"257F}\symactive{"223C}\symactive{"2248}\symactive{"2260}
\symactive{"01F8}\symactive{"01F9}\symactive{"1E3E}\symactive{"1E3F}
\symhyph{"2011}\symhyph{"2E1A}\symblankrange{"1F300}{"1FAFF}\symblank{"FE0F}
% Caractères chinois : WenQuanYi Zen Hei (sous-ensemble GB2312 + pinyin), gras/italique simulés
\setCJKmainfont{WenQuanYiZenHei-subset.ttf}[Path=../../../fonts/,AutoFakeBold=3,AutoFakeSlant=0.2]
\xeCJKsetup{CJKspace=false}
% Pinyin : voyelles à caron (ǎ ǐ ǒ ǔ ǖ ǘ ǚ ǜ) absentes de la police principale -> caractères actifs
\newfontface\pyfont{WenQuanYiZenHei-subset.ttf}[Path=../../../fonts/]
\newcommand\pyactive[1]{\catcode#1=13 \begingroup\lccode`\~=#1 \lowercase{\endgroup\def~}{{\pyfont\char#1}}}
\pyactive{"01CD}\pyactive{"01CE}\pyactive{"01CF}\pyactive{"01D0}\pyactive{"01D1}\pyactive{"01D2}\pyactive{"01D3}\pyactive{"01D4}\pyactive{"01D5}\pyactive{"01D6}\pyactive{"01D7}\pyactive{"01D8}\pyactive{"01D9}\pyactive{"01DA}\pyactive{"01DB}\pyactive{"01DC}
% Maths en sans-serif (Fira Math) avec chiffres et lettres droites de Plus
% Jakarta, pour que les formules se fondent dans le texte.
\usepackage[math-style=ISO,bold-style=ISO]{unicode-math}
\setmathfont{FiraMath-Regular.otf}[Path=../../../fonts/]
\setmathfont{PlusJakartaSans-Medium.ttf}[Path=../../../fonts/,range={up/num,up/{Latin,latin},\mathup}]
\usepackage{icomma} % virgule décimale sans espace en mode math
% Caractères mathématiques Unicode absents des polices de texte (Plus Jakarta,
% Archivo Black) : rendus par leur équivalent mathématique, en texte comme en
% formule — sinon ils s'affichent en carré vide (ex. « Arithmétique dans ℤ »).
\usepackage{newunicodechar}
\newunicodechar{ℝ}{\ensuremath{\mathbb{R}}}
\newunicodechar{ℤ}{\ensuremath{\mathbb{Z}}}
\newunicodechar{ℂ}{\ensuremath{\mathbb{C}}}
\newunicodechar{ℕ}{\ensuremath{\mathbb{N}}}
\newunicodechar{ℚ}{\ensuremath{\mathbb{Q}}}
\newunicodechar{𝔻}{\ensuremath{\mathbb{D}}}
\newunicodechar{α}{\ensuremath{\alpha}}
\newunicodechar{β}{\ensuremath{\beta}}
\newunicodechar{γ}{\ensuremath{\gamma}}
\newunicodechar{δ}{\ensuremath{\delta}}
\newunicodechar{ε}{\ensuremath{\varepsilon}}
\newunicodechar{θ}{\ensuremath{\theta}}
\newunicodechar{λ}{\ensuremath{\lambda}}
\newunicodechar{μ}{\ensuremath{\mu}}
\newunicodechar{σ}{\ensuremath{\sigma}}
\newunicodechar{τ}{\ensuremath{\tau}}
\newunicodechar{φ}{\ensuremath{\varphi}}
\newunicodechar{ψ}{\ensuremath{\psi}}
\newunicodechar{ω}{\ensuremath{\omega}}
\newunicodechar{Σ}{\ensuremath{\Sigma}}
% majuscules produites par \MakeUppercase dans les titres (α-aminés -> Α)
\newunicodechar{Α}{\ensuremath{\alpha}}
\newunicodechar{Β}{\ensuremath{\beta}}
\newunicodechar{Γ}{\ensuremath{\Gamma}}
\newunicodechar{Δ}{\ensuremath{\Delta}}
\newunicodechar{Ε}{\ensuremath{\varepsilon}}
\newunicodechar{Θ}{\ensuremath{\Theta}}
\newunicodechar{Λ}{\ensuremath{\Lambda}}
\newunicodechar{Μ}{\ensuremath{\mu}}
\newunicodechar{Τ}{\ensuremath{\tau}}
\newunicodechar{Φ}{\ensuremath{\Phi}}
\newunicodechar{Ψ}{\ensuremath{\Psi}}
\newunicodechar{Ω}{\ensuremath{\Omega}}
\newunicodechar{↦}{\ensuremath{\mapsto}}
\newunicodechar{∈}{\ensuremath{\in}}
\newunicodechar{∉}{\ensuremath{\notin}}
\newunicodechar{∧}{\ensuremath{\wedge}}
\newunicodechar{⇔}{\ensuremath{\Leftrightarrow}}
\newunicodechar{⟺}{\ensuremath{\Longleftrightarrow}}
\newunicodechar{⇒}{\ensuremath{\Rightarrow}}
\newunicodechar{⟹}{\ensuremath{\Longrightarrow}}
\newunicodechar{∘}{\ensuremath{\circ}}
\newunicodechar{∪}{\ensuremath{\cup}}
\newunicodechar{∩}{\ensuremath{\cap}}
\newunicodechar{≡}{\ensuremath{\equiv}}
\newunicodechar{⊂}{\ensuremath{\subset}}
\newunicodechar{⊥}{\ensuremath{\perp}}
\newunicodechar{∃}{\ensuremath{\exists}}
\newunicodechar{∀}{\ensuremath{\forall}}
\newunicodechar{∛}{\ensuremath{\sqrt[3]{\,}}}
\newunicodechar{∜}{\ensuremath{\sqrt[4]{\,}}}
\newunicodechar{⃗}{\ensuremath{{}^{\to}}}
\newunicodechar{ⁿ}{\ifmmode^{n}\else\textsuperscript{\ensuremath{n}}\fi}
\newunicodechar{ˣ}{\ifmmode^{x}\else\textsuperscript{\ensuremath{x}}\fi}
\newunicodechar{ᵇ}{\ifmmode^{b}\else\textsuperscript{\ensuremath{b}}\fi}
\newunicodechar{ʳ}{\ifmmode^{r}\else\textsuperscript{\ensuremath{r}}\fi}
\newunicodechar{ᵘ}{\ifmmode^{u}\else\textsuperscript{\ensuremath{u}}\fi}
\newunicodechar{ʸ}{\ifmmode^{y}\else\textsuperscript{\ensuremath{y}}\fi}
\newunicodechar{ᵃ}{\ifmmode^{a}\else\textsuperscript{\ensuremath{a}}\fi}
\newunicodechar{ᵉ}{\ifmmode^{e}\else\textsuperscript{\ensuremath{e}}\fi}
\newunicodechar{ᵏ}{\ifmmode^{k}\else\textsuperscript{\ensuremath{k}}\fi}
\newunicodechar{ᵖ}{\ifmmode^{p}\else\textsuperscript{\ensuremath{p}}\fi}
\newunicodechar{ᴺ}{\ifmmode^{N}\else\textsuperscript{\ensuremath{N}}\fi}
\newunicodechar{ᶜ}{\ifmmode^{c}\else\textsuperscript{\ensuremath{c}}\fi}
\newunicodechar{ʰ}{\ifmmode^{h}\else\textsuperscript{\ensuremath{h}}\fi}
\newunicodechar{ᵠ}{\ifmmode^{q}\else\textsuperscript{\ensuremath{q}}\fi}
\newunicodechar{ₙ}{\ifmmode_{n}\else\textsubscript{\ensuremath{n}}\fi}
\newunicodechar{ₐ}{\ifmmode_{a}\else\textsubscript{\ensuremath{a}}\fi}
\newunicodechar{ᵢ}{\ifmmode_{i}\else\textsubscript{\ensuremath{i}}\fi}
\newunicodechar{ᵣ}{\ifmmode_{r}\else\textsubscript{\ensuremath{r}}\fi}
\newunicodechar{ₖ}{\ifmmode_{k}\else\textsubscript{\ensuremath{k}}\fi}
\newunicodechar{ᵦ}{\ifmmode_{\beta}\else\textsubscript{\ensuremath{\beta}}\fi}
\newunicodechar{ₓ}{\ifmmode_{x}\else\textsubscript{\ensuremath{x}}\fi}
\newfontfamily\popdisplay{ArchivoBlack-Regular.ttf}[Path=../../../fonts/]
\newfontfamily\poplogo{SpaceGrotesk-Bold.ttf}[Path=../../../fonts/]
\newfontfamily\popsemi{PlusJakartaSans-SemiBold.ttf}[Path=../../../fonts/]
\newfontfamily\popxbold{PlusJakartaSans-ExtraBold.ttf}[Path=../../../fonts/]
\newfontfamily\popxboldls{PlusJakartaSans-ExtraBold.ttf}[Path=../../../fonts/,LetterSpace=3.0]

\setlist[enumerate]{leftmargin=1.7em,itemsep=5pt,topsep=4pt}
\setlist[itemize]{leftmargin=1.7em,itemsep=4pt,topsep=4pt,label={\color{poporange}\textbullet}}
\setlength{\parindent}{0pt}
\setlength{\parskip}{5pt}
\renewcommand{\arraystretch}{1.25}

% pgfplots : style Pop par défaut
\pgfplotsset{
  every axis/.append style={axis line style={popink,line width=1pt},tick style={popink},
    grid style={popline,line width=0.5pt},label style={font=\small},tick label style={font=\footnotesize}},
  popcurve/.style={poporange,line width=1.8pt,no markers,smooth},
}
% Même style disponible dans un \draw TikZ simple (hors pgfplots)
\tikzset{popcurve/.style={poporange,line width=1.8pt}}
\tikzset{popgrid/.style={popline,very thin}}
\colorlet{popcurve}{poporange} % le nom du style est parfois utilisé comme couleur

% ============================================================
% Pill / squiggle
% ============================================================
\newcommand{\poppill}[3]{%
  \tikz[baseline=-0.55ex]\node[draw=popink,line width=1.2pt,rounded corners=2.6pt,%
    fill=#2,inner xsep=6.5pt,inner ysep=3pt]{%
    \popxboldls\fontsize{7}{7}\selectfont\color{#3}\MakeUppercase{#1}};%
}
\newcommand{\popsquiggle}[1]{%
  \tikz[baseline=-0.4ex]{
    \draw[#1,line width=2.6pt,line cap=round]
      plot [smooth] coordinates {(0,0.09) (0.34,0.01) (0.68,0.21) (1.02,0.01) (1.36,0.10)};
  }%
}
% Mot-clé surligné (marqueur orange sous le texte)
\newcommand{\cle}[1]{{\popxbold\color{popdark}#1}}

% ============================================================
% En-tête / pied
% ============================================================
\pagestyle{fancy}
\fancyhf{}
\renewcommand{\headrulewidth}{0pt}
\renewcommand{\footrulewidth}{0pt}
\lhead{\raisebox{-0.15ex}{\poplogo\fontsize{13}{13}\selectfont\color{popink}OnBuch\color{poporange}+}}
\rhead{\poppill{\DOCMATIERE\ \textbullet\ \DOCNIVEAU\ \textbullet\ Leçon \DOCLECON}{white}{popink}}
\lfoot{\popxbold\fontsize{7.6}{7.6}\selectfont\color{popmuted}\MakeUppercase{Cours signé OnBuch+}\ \textbullet\ {\popsemi\DOCTITRE}}
\rfoot{\tikz[baseline=-0.6ex]\node[rounded corners=8.5pt,fill=popink,minimum height=17pt,minimum width=17pt,inner xsep=5pt,inner sep=0]{\popxbold\fontsize{8}{8}\selectfont\color{popcream}\thepage\,/\,\pageref*{LastPage}};}

% ============================================================
% Titres
% ============================================================
\newcommand{\chaptitre}[2]{%
  \begin{tcolorbox}[enhanced,colback=poporange,colframe=popink,boxrule=2.6pt,arc=11pt,
      drop shadow=popink,left=15pt,right=15pt,top=12pt,bottom=11pt,before skip=3pt,after skip=18pt]
    \poppill{#2}{popink}{popcream}\par\vspace{9pt}
    {\raggedright\hyphenpenalty=10000\exhyphenpenalty=10000\popdisplay\fontsize{22}{24}\selectfont\color{popcream}\MakeUppercase{#1}\par}
  \end{tcolorbox}
}
% Section numérotée : pastille orange avec le numéro + titre Archivo
\newcounter{popsec}
\newcommand{\coursec}[1]{%
  \stepcounter{popsec}\setcounter{popsub}{0}%
  \par\vspace{16pt}\noindent
  \begin{minipage}{\linewidth}
  \tikz[baseline=-0.7ex]\node[draw=popink,line width=1.6pt,fill=poporange,rounded corners=4pt,minimum size=22pt,inner sep=0]{\popdisplay\fontsize{12}{12}\selectfont\color{popcream}\thepopsec};%
  \hspace{8pt}{\popdisplay\fontsize{15}{17}\selectfont\color{popink}\MakeUppercase{#1}}\par
  \vspace{4pt}\noindent{\color{poporange}\rule{36pt}{3.6pt}}
  \end{minipage}\par\nopagebreak\vspace{8pt}
}
\newcounter{popsub}
\newcommand{\courssub}[1]{%
  \stepcounter{popsub}%
  \par\vspace{9pt}\noindent{\popxbold\fontsize{11.5}{13}\selectfont\color{popink}{\color{poporange}\thepopsec.\thepopsub}\hspace{6pt}#1}\par\nopagebreak\vspace{4pt}
}

% ============================================================
% Boîtes pédagogiques — même recette : fond blanc, bordure épaisse colorée,
% ombre dure encre, badge plein. Titre optionnel [#1].
% ============================================================
\tcbset{popbase/.style={breakable,enhanced,colback=white,boxrule=2.3pt,arc=9pt,
  shadow={3pt}{-3pt}{0pt}{popink},left=14pt,right=14pt,top=11pt,bottom=11pt,before skip=11pt,after skip=15pt,
  fontupper=\popsemi\fontsize{10.6}{14.5}\selectfont\color{popink},
  fontlower=\popsemi\fontsize{10.6}{14.5}\selectfont\color{popink}}}
% Coche dessinée (le glyphe ✓ n'existe pas dans les polices du document)
\newcommand{\popcheck}[1]{\tikz[baseline=-0.5ex]\draw[#1,line width=1.6pt,line cap=round,line join=round](-0.13,0.0)--(-0.04,-0.09)--(0.13,0.1);}
\providecommand{\coche}{\popcheck{popgreen}}
\providecommand{\resultat}[1]{\par\smallskip\noindent\fcolorbox{popgreen}{popgreenL}{\parbox{\dimexpr\linewidth-2\fboxsep-2\fboxrule\relax}{\textbf{Résultat.} #1}}\par\smallskip}
\newcommand{\popboxhead}[3]{\poppill{#1}{#2}{white}\ifx\relax#3\relax\else\hspace{6pt}{\popxbold\fontsize{10.6}{12}\selectfont\color{popink}#3}\fi\par\vspace{7pt}}

\newtcolorbox{aretenir}[1][]{popbase,colframe=popgreen,before upper={\popboxhead{À retenir}{popgreen}{#1}}}
% Texte en chinois (document de lecture, dialogue, extrait) : cadre
% d'étude. Un titre optionnel (source, auteur) s'affiche dans l'en-tête.
\newtcolorbox{textebox}[1][]{popbase,colframe=popblue,before upper={\popboxhead{文本 · Texte}{popblue}{#1}},after upper={}}
% Marques ✓ / ✗ (forme correcte / fautive) souvent utilisées par les modèles
\providecommand{\cross}{\textcolor{poppink}{\ensuremath{\times}}}
\providecommand{\xmark}{\cross}\providecommand{\crossmark}{\cross}\providecommand{\cmark}{\ensuremath{\checkmark}}
% Ligne de réponse / justification à compléter (inventée par les modèles)
\providecommand{\justification}[1]{\par\noindent\textit{Justification~:} \dotfill\par}
% \textipa (package tipa non chargé) : la police ne couvre pas l'API -> simple passage
\providecommand{\textipa}[1]{#1}
% Soulignement (ulem non chargé) : \uline, \ul
\providecommand{\uline}[1]{\underline{#1}}\providecommand{\ul}[1]{\underline{#1}}
% \glossaire{\item ...} inventée par les modèles : liste de glossaire
\providecommand{\glossaire}[1]{\begin{itemize}#1\end{itemize}}
% \alert (beamer) inventée par les modèles : texte mis en évidence
\providecommand{\alert}[1]{\textbf{\textcolor{poppink}{#1}}}
% variantes ASCII d'ordinaux écrites en macro
\providecommand{\iere}{\textsuperscript{re}}\providecommand{\ieme}{\textsuperscript{e}}\providecommand{\ere}{\textsuperscript{re}}\providecommand{\eme}{\textsuperscript{e}}\providecommand{\er}{\textsuperscript{er}}\providecommand{\ie}{\textsuperscript{e}}
% Ordinaux français écrits en macro par les modèles : 1\ière, 3\ième
\newcommand{\ière}{\textsuperscript{re}}\newcommand{\ième}{\textsuperscript{e}}
% \exemple (inventée par les modèles) : étiquette « Exemple : »
\providecommand{\exemple}{\textbf{Exemple~:}\ }
% \square utilisable aussi hors mode math (cases à cocher)
\AtBeginDocument{\let\squareMath\square\renewcommand{\square}{\ensuremath{\squareMath}}}
% Mot ou phrase en chinois à l'intérieur d'un paragraphe français
\newcommand{\zh}[1]{\ifmmode\text{#1}\else{#1}\fi}
\let\eng\zh \let\esp\zh \let\all\zh % alias tolérés
% flèches d'intonation inventées par les modèles
\providecommand{\textsearrow}{}\renewcommand{\textsearrow}{\ensuremath{\searrow}}\providecommand{\textnearrow}{}\renewcommand{\textnearrow}{\ensuremath{\nearrow}}
% \fr{..} : mot français (inventé par les modèles)
\providecommand{\fr}[1]{\textit{#1}}
% zhbox : encadré de texte chinois inventé par les modèles
\newenvironment{zhbox}{\begin{quote}}{\end{quote}}
% \courssubsub : titre de niveau 3 inventé par les modèles
\providecommand{\courssubsub}[1]{\par\medskip\noindent\textbf{#1}\par\smallskip}
% \zhbf : texte chinois en gras inventé par les modèles
\providecommand{\zhbf}[1]{\textbf{#1}}
% \vs : « versus » inventé par les modèles
\providecommand{\vs}{\textit{vs}\ }
% \blank : case à compléter inventée par les modèles
\providecommand{\blank}[1][]{\underline{\hspace{1.6em}}}
\newtcolorbox{exemplebox}[1][]{popbase,colframe=poporange,before upper={\popboxhead{Exemple}{poporange}{#1}}}
\newtcolorbox{definition}[1][]{popbase,colframe=popblue,before upper={\popboxhead{Définition}{popblue}{#1}}}
\newtcolorbox{propriete}[1][]{popbase,colframe=poppurple,before upper={\popboxhead{Propriété}{poppurple}{#1}}}
\newtcolorbox{methode}[1][]{popbase,colframe=popgold,before upper={\popboxhead{Méthode}{popgold}{#1}}}
\newtcolorbox{attention}[1][]{popbase,colframe=poppink,before upper={\popboxhead{Attention}{poppink}{#1}}}
\newtcolorbox{experience}[1][]{popbase,colframe=popblue,colback=popblueL,before upper={\popboxhead{Expérience}{popblue}{#1}}}
% « Le savais-tu ? » : carte violette pleine (contexte, culture, Cameroun)
\newtcolorbox{savaistu}[1][]{popbase,colback=poppurple,colframe=popink,
  fontupper=\popsemi\fontsize{10.4}{14.2}\selectfont\color{white},
  before upper={\poppill{Le savais-tu\,?}{white}{poppurple}\ifx\relax#1\relax\else\hspace{6pt}{\popxbold\color{white}#1}\fi\par\vspace{7pt}}}
% Exercice résolu : énoncé en haut, \tcblower puis la solution sur fond crème
\newcounter{popexo}\newcounter{popexr}
\newtcolorbox{exoresolu}[1][]{popbase,colframe=poporange,colbacklower=popcream,
  segmentation style={popink,line width=1.2pt,dashed},
  before upper={\stepcounter{popexr}\popboxhead{Exercice résolu \thepopexr}{poporange}{#1}},
  before lower={\poppill{Solution}{popink}{popcream}\par\vspace{6pt}}}
% Exercice (non résolu en place) — la correction est dans la partie Corrigés
\newtcolorbox{exercice}[1][]{popbase,colframe=popink,
  before upper={\stepcounter{popexo}\popboxhead{Exercice \thepopexo}{popink}{#1}}}
% Corrigé
\newtcolorbox{corrige}[1][]{popbase,colframe=popgreen,colback=popgreenL,
  before upper={\popboxhead{Corrigé}{popgreen}{#1}}}
% Objectifs / prérequis / bilan
\newtcolorbox{objectifs}{popbase,colframe=popink,colback=poporangeL,
  before upper={\popboxhead{Objectifs de la leçon}{popink}{}}}
\newtcolorbox{prerequis}{popbase,colframe=popink,colback=popblueL,
  before upper={\popboxhead{Prérequis}{popblue}{}}}
\newtcolorbox{bilan}{popbase,colframe=popink,colback=white,boxrule=2.8pt,
  before upper={\popboxhead{Fiche bilan}{poporange}{L'essentiel à connaître pour l'examen}}}
% Figure encadrée avec légende
\newtcolorbox{popfigure}[1][]{enhanced,breakable=false,colback=white,colframe=popline,boxrule=1.4pt,arc=8pt,
  left=10pt,right=10pt,top=10pt,bottom=8pt,before skip=10pt,after skip=12pt,halign=center,
  lower separated=false,halign lower=center,
  fontlower=\popsemi\fontsize{9}{11.5}\selectfont\color{popmuted},
  #1}
\newcommand{\legende}[1]{\par\vspace{6pt}{\popsemi\fontsize{9}{11.5}\selectfont\color{popmuted}#1}}

% ============================================================
% Signature OnBuch+
% ============================================================
\newcommand{\poplogoinline}[1]{{\poplogo\fontsize{#1}{#1}\selectfont\color{popink}OnBuch\color{poporange}+}}
\newcommand{\signatureonbuch}{%
  \begin{tcolorbox}[enhanced,colback=popink,colframe=popink,boxrule=2.2pt,arc=10pt,
      drop shadow=poporange,left=14pt,right=14pt,top=10pt,bottom=10pt,before skip=0pt,after skip=0pt]
    \tikz[baseline=-0.7ex]\node[circle,fill=popgreen,draw=popcream,line width=1.3pt,minimum size=22pt,inner sep=0]{\popcheck{white}};%
    \hspace{9pt}\begin{minipage}[c]{0.62\linewidth}
      {\popxbold\fontsize{12}{13}\selectfont\color{popcream}Signé\ \textbullet\ {\poplogo OnBuch\color{poporange}+}}\par\vspace{2pt}
      {\raggedright\popsemi\fontsize{8}{10.5}\selectfont\color{popline}\MakeUppercase{Équipe pédagogique OnBuch+}\\\MakeUppercase{Conforme au programme officiel MINESEC}\par}
    \end{minipage}\hfill
    {\poplogo\fontsize{22}{22}\selectfont\color{popcream}OnBuch\color{poporange}+}
  \end{tcolorbox}%
}

% ============================================================
% Page de garde
% #1 titre · #2 matière · #3 niveau · #4 module · #5 n° leçon · #6 durée ·
% #7 description / objectifs (texte ou itemize)
% ============================================================
\newcommand{\pagegarde}[7]{%
  \thispagestyle{empty}
  \newgeometry{a4paper,margin=1.7cm,top=1.6cm,bottom=1.4cm}%
  \begin{tikzpicture}[remember picture,overlay]
    \fill[poporange!18] ([xshift=-3.2cm,yshift=-3.6cm]current page.north east) circle (4.6cm);
    \fill[poppurple!14] ([xshift=2.4cm,yshift=7.6cm]current page.south west) circle (3.8cm);
  \end{tikzpicture}%
  {\poplogo\fontsize{30}{30}\selectfont\color{popink}OnBuch\color{poporange}+}\hfill
  \raisebox{0.6ex}{\poppill{Cours signé \textbullet\ Premium}{popink}{popcream}}\par
  \vspace{1pt}\popsquiggle{poppurple}\par
  \vspace{8pt}
  \begin{tcolorbox}[enhanced,colback=poporange,colframe=popink,boxrule=2.8pt,arc=13pt,
      drop shadow=popink,left=18pt,right=18pt,top=12pt,bottom=13pt,after skip=10pt]
    \poppill{#2 \textbullet\ #3}{popink}{popcream}\hspace{5pt}\poppill{#4}{white}{popink}\par\vspace{8pt}
    {\popdisplay\fontsize{14}{14}\selectfont\color{popink}LEÇON #5}\par\vspace{4pt}
    {\raggedright\hyphenpenalty=10000\exhyphenpenalty=10000\popdisplay\fontsize{25}{27}\selectfont\color{popcream}\MakeUppercase{#1}\par}
    \vspace{8pt}
    {\popxbold\fontsize{9.5}{9.5}\selectfont\color{popink}\MakeUppercase{Durée conseillée : #6}}
  \end{tcolorbox}
  \begin{tcolorbox}[enhanced,colback=white,colframe=popink,boxrule=2.2pt,arc=10pt,
      drop shadow=popink,left=16pt,right=16pt,top=10pt,bottom=10pt,after skip=8pt]
    \poppill{Dans cette leçon}{poppurple}{white}\par\vspace{5pt}
    \popsemi\fontsize{9.3}{12.6}\selectfont\color{popink}#7
  \end{tcolorbox}
  \noindent\begin{minipage}[t]{0.487\linewidth}
    \begin{tcolorbox}[enhanced,colback=popgreen,colframe=popink,boxrule=2.2pt,arc=10pt,
        drop shadow=popink,left=11pt,right=11pt,top=10pt,bottom=10pt,height=3.1cm,valign=center]
      \tcbox[colback=white,colframe=popink,boxrule=1.3pt,arc=5pt,boxsep=0pt,nobeforeafter,
        left=3pt,right=3pt,top=3pt,bottom=3pt,valign=center]{\qrcode[height=1.9cm]{https://play.google.com/store/apps/details?id=com.luvvix.onbuch&hl=fr}}%
      \hspace{8pt}\begin{minipage}[c]{0.5\linewidth}
        {\popdisplay\fontsize{11}{12.5}\selectfont\color{white}\MakeUppercase{Télécharge OnBuch}}\par\vspace{4pt}
        {\raggedright\popsemi\fontsize{8.4}{11}\selectfont\color{white}Gratuit \textbullet\ Google Play\\Tuteur IA, annales, concours, résultats\par}
      \end{minipage}
    \end{tcolorbox}
  \end{minipage}\hfill
  \begin{minipage}[t]{0.487\linewidth}
    \begin{tcolorbox}[enhanced,colback=poppurple,colframe=popink,boxrule=2.2pt,arc=10pt,
        drop shadow=popink,left=11pt,right=11pt,top=10pt,bottom=10pt,height=3.1cm,valign=center]
      \tikz[baseline=-0.7ex]\node[circle,fill=white,draw=popink,line width=1.3pt,minimum size=1.6cm,inner sep=0]{\popdisplay\fontsize{24}{24}\selectfont\color{poppurple}+};%
      \hspace{8pt}\begin{minipage}[c]{0.6\linewidth}
        {\popdisplay\fontsize{11}{12.5}\selectfont\color{white}\MakeUppercase{Passe à OnBuch+}}\par\vspace{4pt}
        {\raggedright\popsemi\fontsize{8.4}{11}\selectfont\color{white}Tous les cours signés, Tuteur IA illimité, fiches PDF — onglet OnBuch+ de l'app\par}
      \end{minipage}
    \end{tcolorbox}
  \end{minipage}
  \vspace{8pt}
  \popcontenu
  \vfill
  \signatureonbuch
  \restoregeometry
  \newpage
}

% Rangée « Ce cours contient » (page de garde)
\newcommand{\poptile}[3]{% #1 couleur #2 gros texte #3 légende
  \begin{tcolorbox}[enhanced,colback=#1,colframe=popink,boxrule=2pt,arc=9pt,shadow={2.5pt}{-2.5pt}{0pt}{popink},
      left=6pt,right=6pt,top=9pt,bottom=9pt,halign=center,valign=center,height=1.75cm,width=\linewidth,nobeforeafter]
    {\popdisplay\fontsize{12.5}{13}\selectfont\color{white}\MakeUppercase{#2}}\par\vspace{3pt}
    {\centering\popsemi\fontsize{7.8}{9.5}\selectfont\color{white}#3\par}
  \end{tcolorbox}}
\newcommand{\popcontenu}{%
  \par\noindent{\popxboldls\fontsize{7.5}{7.5}\selectfont\color{popmuted}CE COURS SIGNÉ CONTIENT}\par\vspace{7pt}
  \noindent\begin{minipage}[t]{0.235\linewidth}\poptile{popblue}{Cours}{complet, illustré, pas à pas}\end{minipage}\hfill
  \begin{minipage}[t]{0.235\linewidth}\poptile{poporange}{Méthodes}{et exercices résolus}\end{minipage}\hfill
  \begin{minipage}[t]{0.235\linewidth}\poptile{poppink}{Exercices}{type examen + corrigés}\end{minipage}\hfill
  \begin{minipage}[t]{0.235\linewidth}\poptile{popgreen}{Fiche bilan}{l'essentiel à retenir}\end{minipage}\par}

% Sommaire
\newtcolorbox{sommaire}{enhanced,colback=white,colframe=popink,boxrule=2.2pt,arc=10pt,
  drop shadow=popink,left=18pt,right=18pt,top=13pt,bottom=13pt,before skip=6pt,after skip=20pt,
  before upper={\poppill{Sommaire}{poppurple}{white}\par\vspace{9pt}\popsemi\fontsize{10.6}{14.5}\selectfont\color{popink}}}

% Signature de fin de cours — poussée tout en bas de la dernière page
\newcommand{\findecours}{%
  \par\vfill\nopagebreak
  \begin{center}\popsquiggle{poporange}\end{center}\vspace{4pt}
  \signatureonbuch
}
\AtBeginDocument{\def\llbracket{\mathopen{[\mkern-3.5mu[}}\def\rrbracket{\mathclose{]\mkern-3.5mu]}}} % intervalles d'entiers (glyphe ⟦ absent de la police)
% Glyphes absents de Fira Math : redéfinis à partir de glyphes disponibles
\AtBeginDocument{%
  \def\setminus{\mathbin{\backslash}}\let\smallsetminus\setminus
  \def\vdots{\mathord{\vbox{\baselineskip3.2pt\lineskiplimit0pt\kern4pt\hbox{.}\hbox{.}\hbox{.}}}}%
  \def\ddots{\mathinner{\mkern1mu\raise6.5pt\hbox{.}\mkern2mu\raise3.6pt\hbox{.}\mkern2mu\raise0.7pt\hbox{.}\mkern1mu}}%
  \def\triangle{\mathord{\scalebox{0.9}{$\symup{\Delta}$}}}%
  \def\nabla{\mathord{\raisebox{\depth}{\scalebox{1}[-1]{$\symup{\Delta}$}}}}%
  \def\checkmark{\popcheck{popdark}}%
  \def\star{\mathbin{\ast}}\def\bigstar{\mathord{\ast}}%
  \def\diamond{\mathbin{\scalebox{0.6}{$\Diamond$}}}%
  \def\bigcup{\mathop{\mathchoice{\raisebox{-0.3em}{\scalebox{1.7}{$\cup$}}}{\scalebox{1.25}{$\cup$}}{\cup}{\cup}}\displaylimits}%
  \def\bigcap{\mathop{\mathchoice{\raisebox{-0.3em}{\scalebox{1.7}{$\cap$}}}{\scalebox{1.25}{$\cap$}}{\cap}{\cap}}\displaylimits}%
  \def\lesseqqgtr{\mathrel{\lessgtr}}\def\bigoplus{\mathop{\oplus}\displaylimits}%
}
\providecommand{\ssi}{si et seulement si} % abréviation fréquente
\AtBeginDocument{\providecommand{\wideparen}{\overparen}} % arc de cercle

%% FIGFIX-BEGIN v3 — correctifs globaux des figures (docs/onbuch-generation/tools/figfix)
\usepackage{adjustbox}\usepackage{etoolbox}
% toute figure TikZ/pgfplots plus large que la ligne est réduite à la largeur disponible
\BeforeBeginEnvironment{tikzpicture}{\begin{adjustbox}{max width=\linewidth}}
\AfterEndEnvironment{tikzpicture}{\end{adjustbox}}
% légendes pgfplots lisibles par-dessus les courbes
\pgfplotsset{every axis/.append style={legend style={fill=white,fill opacity=0.92,text opacity=1,draw=black!15,inner sep=3pt}}}
% étiquettes d'arêtes (syntaxe edge["..."]) avec fond blanc
\tikzset{every edge quotes/.append style={fill=white,inner sep=1.2pt}}
% bulles de cartes mentales : texte à la ligne dans la bulle, bulle un peu plus grande
\tikzset{every concept/.append style={text width=2.0cm,align=center,minimum size=2.3cm,inner sep=2pt}}
%% FIGFIX-END
\begin{document}

\pagegarde{Leçon 28 — Grammaire : 把字句}{Chinois}{1ère A}{Module 3 : Citoyenneté et ouverture au monde}{ZHA28}{2 h}{Cette leçon de grammaire présente la phrase en 把 (把字句), structure emblématique du chinois qui met en avant la disposition d'un objet déterminé. Les élèves apprennent son schéma, ses emplois, ses contraintes et les pièges à éviter pour un francophone.
\vspace{4pt}
\begin{multicols}{2}\small
\begin{itemize}
\item À la fin de cette leçon, je sais reconnaître une phrase en 把 dans un texte
\item À la fin de cette leçon, je sais construire le schéma Sujet + 把 + Objet + Verbe + complément
\item À la fin de cette leçon, je sais expliquer pourquoi le verbe d'une phrase en 把 ne peut pas rester « nu »
\item À la fin de cette leçon, je sais employer 把 avec les compléments de direction (在/到/给), de résultat (完/好/成) et 了
\item À la fin de cette leçon, je sais placer correctement la négation (没/别) et les auxiliaires avant 把
\item À la fin de cette leçon, je sais transformer une phrase SVO classique en phrase en 把 quand c'est possible
\end{itemize}\end{multicols}}

\begin{sommaire}
\begin{multicols}{2}
\begin{enumerate}
\item Observation guidée d'exemples
\item Schéma de la structure
\item Emplois et cas d'usage
\item Exceptions et restrictions
\item Comparaison avec le français et pièges des francophones
\item Entraînement : application et transformation
\item Activité d'intégration
\item Méthodes et exercices résolus
\item Exercices
\item Corrigés des exercices
\item Fiche bilan
\end{enumerate}
\end{multicols}
\end{sommaire}

\begin{objectifs}
\begin{itemize}
\item À la fin de cette leçon, je sais reconnaître une phrase en 把 dans un texte
\item À la fin de cette leçon, je sais construire le schéma Sujet + 把 + Objet + Verbe + complément
\item À la fin de cette leçon, je sais expliquer pourquoi le verbe d'une phrase en 把 ne peut pas rester « nu »
\item À la fin de cette leçon, je sais employer 把 avec les compléments de direction (在/到/给), de résultat (完/好/成) et 了
\item À la fin de cette leçon, je sais placer correctement la négation (没/别) et les auxiliaires avant 把
\item À la fin de cette leçon, je sais transformer une phrase SVO classique en phrase en 把 quand c'est possible
\item À la fin de cette leçon, je sais identifier les cas où la structure en 把 est interdite (objet indéfini, verbes sans action sur l'objet)
\item À la fin de cette leçon, je sais comparer la structure en 把 avec la construction française et éviter les calques
\item À la fin de cette leçon, je sais rédiger un court paragraphe utilisant au moins trois phrases en 把 correctes
\end{itemize}
\end{objectifs}

\begin{prerequis}
\begin{itemize}
\item Ordre de base de la phrase chinoise Sujet-Verbe-Objet (Seconde)
\item Compléments de résultat : 完, 好, 懂, 见 (début d'année)
\item Compléments de direction avec 在, 到, 给 et la particule 了
\item Négations 不 et 没(有) et leur emploi
\item Notion d'objet déterminé vs indéterminé (un/une vs le/la/ce)
\end{itemize}
\end{prerequis}

\newpage

\coursec{Observation guidée d'exemples}

\courssub{Situation d'accroche : le rangement de la salle de classe}

Il est 17 h au lycée de Yaoundé. Les cours sont terminés, et le professeur de chinois supervise le rangement de la salle. Il s'adresse à un élève : « Mets les cahiers sur la table ! » En chinois, pour insister sur ce qu'on fait \emph{subir} à un objet précis, on utilise une structure spéciale : \zh{把本子放在桌子上} (\emph{Bǎ běnzi fàng zài zhuōzi shang.} — « Prends les cahiers et pose-les sur la table »). De même, quand un élève camerounais annonce fièrement « J'ai fini mes devoirs », le chinois préfère dire \zh{我把作业做完了} (\emph{Wǒ bǎ zuòyè zuò wán le.}).

Cette structure en \zh{把}, appelée \cle{phrase en 把} (\zh{把字句}, \emph{bǎzìjù}), est absente du français : elle n'a pas d'équivalent direct. C'est pourtant l'une des structures les plus fréquentes du chinois parlé et écrit, et l'une des plus redoutées des apprenants francophones. \textbf{Problématique :} comment le chinois déplace-t-il l'objet devant le verbe, et comment construire cette phrase sans faute ? Observons d'abord des exemples avant de dégager la règle.

\courssub{Lecture guidée : quatre phrases en 把}

\begin{textebox}[Quatre phrases à observer — lire à voix haute en respectant les tons]
\zh{我把饭吃完了。}\\
\emph{Wǒ bǎ fàn chī wán le.}\\
« J'ai fini de manger le riz. »\\[0.4em]
\zh{他把书放在桌子上。}\\
\emph{Tā bǎ shū fàng zài zhuōzi shang.}\\
« Il a posé le livre sur la table. »\\[0.4em]
\zh{请把门关上。}\\
\emph{Qǐng bǎ mén guān shang.}\\
« Ferme la porte, s'il te plaît. »\\[0.4em]
\zh{妈妈把衣服洗干净了。}\\
\emph{Māma bǎ yīfu xǐ gānjìng le.}\\
« Maman a lavé les vêtements (ils sont maintenant propres). »
\end{textebox}

Lis chaque phrase deux fois : une première fois lentement, syllabe par syllabe, une deuxième fois à vitesse naturelle. Fais attention au troisième ton de \zh{把} (\emph{bǎ}) et au deuxième ton de \zh{完} (\emph{wán}), qui ne doit pas être avalé dans la rapidité du discours.

\courssub{Questions guides d'observation}

Observe les quatre phrases et réponds aux questions suivantes :

\begin{enumerate}
\item \textbf{Où se trouve l'objet ?} Dans une phrase française (et dans une phrase chinoise simple Sujet–Verbe–Objet), l'objet se trouve \emph{après} le verbe : « manger \emph{le riz} », \zh{吃饭}. Ici, où sont passés \zh{饭}, \zh{书}, \zh{门}, \zh{衣服} ?
\item \textbf{Quel mot précède l'objet ?} Quel petit mot revient dans les quatre phrases juste avant l'objet ?
\item \textbf{Que remarque-t-on après le verbe ?} Le verbe est-il seul ? Que trouve-t-on après \zh{吃}, \zh{放}, \zh{关}, \zh{洗} ?
\end{enumerate}

\begin{exemplebox}[Réponses attendues]
\begin{itemize}
\item L'objet (\zh{饭}, \zh{书}, \zh{门}, \zh{衣服}) est passé \textbf{avant le verbe}, introduit par \zh{把}.
\item Le mot-outil \zh{把} (\emph{bǎ}) revient dans les quatre phrases : \zh{把 + objet}.
\item Le verbe n'est \textbf{jamais seul} : il est suivi d'un complément (\zh{完}, \zh{在桌子上}, \zh{上}, \zh{干净}) et souvent de \zh{了}.
\end{itemize}
\end{exemplebox}

\courssub{Comparaison phrase SVO simple / phrase en 把}

Pour bien voir le déplacement, comparons la phrase ordinaire et la phrase en \zh{把} :

\begin{center}
\begin{tabularx}{\linewidth}{|X|X|X|}
\hline
\rowcolor{popblueL} Phrase SVO simple & Phrase en 把 & Traduction \\
\hline
\zh{我吃完饭了。}\newline \emph{Wǒ chī wán fàn le.} & \zh{我把饭吃完了。}\newline \emph{Wǒ bǎ fàn chī wán le.} & J'ai fini de manger le riz. \\
\hline
\zh{他放书在桌子上。}\newline \emph{Tā fàng shū zài zhuōzi shang.} & \zh{他把书放在桌子上。}\newline \emph{Tā bǎ shū fàng zài zhuōzi shang.} & Il a posé le livre sur la table. \\
\hline
\zh{请关门。}\newline \emph{Qǐng guān mén.} & \zh{请把门关上。}\newline \emph{Qǐng bǎ mén guān shang.} & Ferme la porte, s'il te plaît. \\
\hline
\zh{妈妈洗干净衣服了。}\newline \emph{Māma xǐ gānjìng yīfu le.} & \zh{妈妈把衣服洗干净了。}\newline \emph{Māma bǎ yīfu xǐ gānjìng le.} & Maman a lavé les vêtements. \\
\hline
\end{tabularx}
\end{center}

\begin{popfigure}
\begin{tikzpicture}[font=\small]
\node[draw=popblue, fill=popblueL, rounded corners, align=center, minimum width=5.6cm, minimum height=1.5cm] (svo) at (0,0) {\textbf{Phrase SVO}\\ 我 \textcolor{popgreen}{吃完了} \textcolor{popink}{饭}。\\ \emph{S + V + compl. + O}};
\node[draw=poporange, fill=poporangeL, rounded corners, align=center, minimum width=5.6cm, minimum height=1.5cm] (ba) at (8.5,0) {\textbf{Phrase en 把}\\ 我 \textcolor{popink}{把饭} \textcolor{popblue}{吃完了}。\\ \emph{S + 把 + O + V + compl.}};
\draw[-{Stealth[length=3mm]}, very thick, popink, bend left=25] (svo.north) to node[above, align=center]{l'objet \textcolor{popink}{饭} passe\\ devant le verbe} (ba.north);
\draw[-{Stealth[length=3mm]}, very thick, popmuted] (svo.east) -- (ba.west) node[midway, below, align=center]{transformation};
\end{tikzpicture}
\legende{De la phrase SVO à la phrase en 把 : l'objet (en rouge) passe avant le verbe ; le verbe et son complément (en bleu) restent groupés après.}
\end{popfigure}

\courssub{Dégagement collectif : ce que fait 把}

\begin{aretenir}[Première conclusion d'observation]
Dans une phrase en \zh{把} :
\begin{enumerate}
\item \zh{把} fait passer \textbf{l'objet avant le verbe} : Sujet + \zh{把} + Objet + Verbe + \ldots
\item Le verbe \textbf{porte obligatoirement un complément} (résultat, direction, lieu, \zh{了}\ldots) : il ne reste jamais seul.
\item La phrase insiste sur \textbf{ce qu'on fait subir à l'objet} : on le « prend en main » et on le transforme, déplace ou traite.
\end{enumerate}
\end{aretenir}

Comparaison avec le français : le français garde toujours l'ordre Sujet–Verbe–Objet (« Maman a lavé \emph{les vêtements} ») et ne possède aucun mot-outil comme \zh{把}. Pour rendre l'idée, on peut paraphraser : « Maman \emph{a pris} les vêtements et les a lavés (résultat : ils sont propres) ». Cette paraphrase aide à comprendre, mais en chinois la structure en \zh{把} est naturelle et très courante.

\begin{exemplebox}[Deux exemples supplémentaires traduits]
\zh{我把作业做完了。}\\
\emph{Wǒ bǎ zuòyè zuò wán le.}\\
« J'ai terminé mes devoirs. » (l'objet \zh{作业} passe devant \zh{做}, complément de résultat \zh{完})\\[0.4em]
\zh{他把钱给了弟弟。}\\
\emph{Tā bǎ qián gěi le dìdi.}\\
« Il a donné l'argent à son petit frère. » (l'objet \zh{钱} passe devant \zh{给}, suivi de \zh{了} + le bénéficiaire)
\end{exemplebox}

\begin{savaistu}[D'où vient 把 ?]
À l'origine, \zh{把} était un verbe plein signifiant « prendre en main, saisir ». On le retrouve encore dans \zh{一把伞} (\emph{yì bǎ sǎn}, « un parapluie », littéralement « une poignée de parapluie ») où \zh{把} sert de classificateur des objets qu'on tient à la main. La structure en \zh{把} garde cette image : le sujet \emph{prend l'objet en main} et lui fait subir quelque chose — on le mange jusqu'au bout (\zh{吃完}), on le pose quelque part (\zh{放在}), on le ferme (\zh{关上}), on le rend propre (\zh{洗干净}).
\end{savaistu}

\begin{methode}[Repérer une phrase en 把 en quatre pas]
\begin{enumerate}
\item Cherche le mot \zh{把} dans la phrase.
\item Vérifie qu'un \textbf{nom (objet)} le suit immédiatement : \zh{把 + nom}.
\item Vérifie que ce groupe \zh{把 + objet} se trouve \textbf{avant le verbe}.
\item Vérifie que le verbe est \textbf{suivi d'un complément} (\zh{完}, \zh{干净}, \zh{在}\ldots, \zh{上}\ldots) ou au moins de \zh{了}. Si le verbe est seul, la phrase est fautive.
\end{enumerate}
\end{methode}

\begin{exoresolu}[Repérage guidé]
Énoncé : dans les phrases suivantes, souligne l'objet déplacé et entoure le complément du verbe. (a) \zh{我把水喝完了。} (b) \zh{他把电脑放在椅子上。}
\tcblower
Solution : (a) \zh{我把}\underline{\zh{水}}\zh{喝}\textbf{完}\zh{了} — l'objet est \zh{水} (\emph{shuǐ}, « l'eau »), le complément de résultat est \zh{完} (« fini ») : « J'ai fini de boire l'eau. » (b) \zh{他把}\underline{\zh{电脑}}\zh{放}\textbf{在椅子上} — l'objet est \zh{电脑} (\emph{diànnǎo}, « l'ordinateur »), le complément de lieu est \zh{在椅子上} (« sur la chaise ») : « Il a posé l'ordinateur sur la chaise. » Dans les deux cas, \zh{把} + objet précède le verbe, et le verbe n'est jamais seul.
\end{exoresolu}

\begin{attention}[Piège immédiat des francophones]
Ne traduis jamais mot à mot « J'ai fini mes devoirs » par \zh{我做完了作业} quand tu veux insister sur le résultat subi par l'objet : le chinois préfère alors \zh{我把作业做完了}. Et surtout, une phrase en \zh{把} avec un verbe seul est \textbf{impossible} : on ne dit jamais \zh{我把饭吃}. Le verbe doit toujours être « habillé » d'un complément ou de \zh{了} : \zh{我把饭吃了} ou \zh{我把饭吃完了}.
\end{attention}

\begin{exercice}[À toi d'observer]
Dans chaque phrase, indique si la structure en \zh{把} est correcte (C) ou incorrecte (I), et justifie : (1) \zh{我把窗户打开了。} (2) \zh{他把歌唱。} (3) \zh{请把灯关上。} (4) \zh{我把信寄。}
\end{exercice}

\begin{corrige}[Exercice « À toi d'observer »]
(1) C : \zh{把} + objet \zh{窗户} avant le verbe \zh{开}, complément \zh{打开} + \zh{了}. (2) I : le verbe \zh{唱} est seul, sans complément ; il faudrait par exemple \zh{他把歌唱完了}. (3) C : phrase impérative correcte, complément directionnel \zh{上}. (4) I : \zh{寄} est seul ; il faut un complément ou \zh{了} : \zh{我把信寄了} ou \zh{我把信寄出去了}.
\end{corrige}

\coursec{Schéma de la structure}

\courssub{La formule générale}

\begin{propriete}[Schéma de la phrase en 把]
\textbf{Sujet + 把 + Objet + Verbe + Complément (ou 了)}

Le groupe \zh{把 + objet} forme un bloc inséparable placé \emph{avant} le verbe. Tout ce qui accompagne le verbe (complément de résultat, de direction, de lieu, \zh{了}, \zh{过}) reste \emph{après} le verbe.
\end{propriete}

\begin{center}
\begin{tabularx}{\linewidth}{|X|X|X|}
\hline
\rowcolor{popblueL} Structure & Exemple en caractères + pinyin & Traduction \\
\hline
S + 把 + O + V + résultat & \zh{我把饭吃完了。}\newline \emph{Wǒ bǎ fàn chī wán le.} & J'ai fini de manger le riz. \\
\hline
S + 把 + O + V + 在 + lieu & \zh{他把书放在桌子上。}\newline \emph{Tā bǎ shū fàng zài zhuōzi shang.} & Il a posé le livre sur la table. \\
\hline
S + 把 + O + V + direction & \zh{请把门关上。}\newline \emph{Qǐng bǎ mén guān shang.} & Ferme la porte, s'il te plaît. \\
\hline
S + 把 + O + V + 了 & \zh{我把信寄了。}\newline \emph{Wǒ bǎ xìn jì le.} & J'ai envoyé la lettre. \\
\hline
S + 把 + O + V + 给 + bénéficiaire & \zh{他把钱给了弟弟。}\newline \emph{Tā bǎ qián gěi le dìdi.} & Il a donné l'argent à son petit frère. \\
\hline
\end{tabularx}
\end{center}

\courssub{Vocabulaire utile de la leçon}

\begin{center}
\begin{tabularx}{\linewidth}{|X|X|X|X|}
\hline
\rowcolor{popblueL} Caractères & Pinyin & Nature & Traduction \\
\hline
\zh{把} & bǎ & mot-outil / classificateur & (marque de l'objet déplacé) ; poignée \\
\hline
\zh{放} & fàng & verbe & poser, placer \\
\hline
\zh{关} & guān & verbe & fermer \\
\hline
\zh{洗} & xǐ & verbe & laver \\
\hline
\zh{干净} & gānjìng & adjectif & propre \\
\hline
\zh{完} & wán & complément de résultat & fini, terminé \\
\hline
\zh{寄} & jì & verbe & envoyer (par la poste) \\
\hline
\zh{窗户} & chuānghu & nom & fenêtre \\
\hline
\zh{灯} & dēng & nom & lampe, lumière \\
\hline
\zh{作业} & zuòyè & nom & devoirs \\
\hline
\zh{本子} & běnzi & nom & cahier \\
\hline
\zh{伞} & sǎn & nom & parapluie \\
\hline
\end{tabularx}
\end{center}

\coursec{Emplois et cas d'usage}

\courssub{Quand utilise-t-on la phrase en 把 ?}

\begin{definition}[Emplois principaux de la phrase en 把]
On emploie la phrase en \zh{把} quand on veut présenter l'action du point de vue de \textbf{ce que subit l'objet} :
\begin{itemize}
\item \textbf{Déplacement} de l'objet : \zh{把书放在桌子上} (poser le livre sur la table), \zh{把椅子搬到外面} (\emph{bǎ yǐzi bān dào wàimiàn}, déplacer la chaise dehors).
\item \textbf{Transformation ou résultat} affectant l'objet : \zh{把衣服洗干净} (laver les vêtements jusqu'à ce qu'ils soient propres), \zh{把饭吃完} (finir le riz).
\item \textbf{Traitement complet} de l'objet : \zh{把信寄出去} (envoyer la lettre), \zh{把作业交给老师} (\emph{bǎ zuòyè jiāo gěi lǎoshī}, remettre les devoirs au professeur).
\end{itemize}
\end{definition}

\begin{exemplebox}[Trois situations camerounaises]
\zh{妈妈把鱼做好了。}\\
\emph{Māma bǎ yú zuò hǎo le.}\\
« Maman a préparé le poisson (il est prêt). »\\[0.4em]
\zh{他把自行车放在学校门口。}\\
\emph{Tā bǎ zìxíngchē fàng zài xuéxiào ménkǒu.}\\
« Il a laissé son vélo devant le portail de l'école. »\\[0.4em]
\zh{请把窗户打开，教室里很热。}\\
\emph{Qǐng bǎ chuānghu dǎkāi, jiàoshì lǐ hěn rè.}\\
« Ouvre la fenêtre, s'il te plaît, il fait très chaud dans la salle. »
\end{exemplebox}

\coursec{Exceptions et restrictions}

\courssub{Restriction 1 : l'objet doit être défini}

\begin{propriete}[L'objet de 把 est connu et précis]
L'objet introduit par \zh{把} doit être \textbf{défini} : le locuteur et l'auditeur savent de quoi il s'agit (\emph{le} livre, \emph{ces} cahiers, \emph{mes} devoirs). On ne peut pas introduire un objet indéfini du type « un livre quelconque » : on ne dit pas \zh{我把一本书放在桌子上} pour « j'ai posé un livre (quelconque) sur la table » ; on dira simplement \zh{我放了一本书在桌子上} ou on précisera \zh{我把那本书放在桌子上} (« j'ai posé ce livre-là sur la table »).
\end{propriete}

\courssub{Restriction 2 : le verbe ne reste jamais seul}

C'est la règle déjà dégagée en observation : le verbe doit être suivi d'un complément (\zh{完}, \zh{干净}, \zh{在} + lieu, \zh{上}, \zh{给} + bénéficiaire) ou au minimum de \zh{了}. \zh{我把门关} est impossible ; \zh{我把门关了} ou \zh{我把门关上} sont corrects.

\courssub{Restriction 3 : les verbes incompatibles avec 把}

\begin{attention}[Verbes qui ne prennent pas 把]
La phrase en \zh{把} exige un verbe d'\textbf{action} qui affecte concrètement l'objet. Les verbes d'état, de sentiment ou de connaissance sont \textbf{impossibles} avec \zh{把} :
\begin{itemize}
\item \zh{是} (être), \zh{有} (avoir), \zh{在} (se trouver) : jamais avec \zh{把}.
\item \zh{喜欢} (aimer), \zh{爱} (aimer), \zh{知道} (savoir), \zh{认识} (connaître), \zh{想} (penser) : on ne dit pas \zh{我把这首歌喜欢}, mais \zh{我喜欢这首歌} (\emph{Wǒ xǐhuan zhè shǒu gē}, « j'aime cette chanson »).
\end{itemize}
\end{attention}

\courssub{Restriction 4 : place de la négation et des adverbes}

\begin{propriete}[La négation se place AVANT 把]
Dans une phrase en \zh{把}, la négation \zh{没(有)} (passé) ou \zh{不} (présent/futur, volonté) et les adverbes modaux (\zh{已经}, \zh{别}, \zh{要}) se placent \textbf{avant} \zh{把}, jamais entre \zh{把} et le verbe :
\begin{itemize}
\item \zh{我没有把作业做完。} \emph{Wǒ méiyǒu bǎ zuòyè zuò wán.} « Je n'ai pas fini mes devoirs. »
\item \zh{别把门关上！} \emph{Bié bǎ mén guān shang!} « Ne ferme pas la porte ! »
\item \zh{他已经把信寄出去了。} \emph{Tā yǐjing bǎ xìn jì chūqù le.} « Il a déjà envoyé la lettre. »
\end{itemize}
Faute typique : \zh{我把作业没有做完} est incorrect.
\end{propriete}

\coursec{Comparaison avec le français et pièges des francophones}

\courssub{Ce qui change par rapport au français}

\begin{center}
\begin{tabularx}{\linewidth}{|X|X|X|}
\hline
\rowcolor{popblueL} Point de comparaison & Français & Chinois (phrase en 把) \\
\hline
Ordre des mots & Sujet + Verbe + Objet : « Il a posé \emph{le livre} sur la table. » & Sujet + 把 + Objet + Verbe + complément : \zh{他把书放在桌子上。} \\
\hline
Marque de l'objet & aucune particule spéciale & le mot-outil \zh{把} introduit l'objet \\
\hline
Verbe seul & possible : « Il a fermé la porte. » & impossible : il faut \zh{关上} ou \zh{了} : \zh{他把门关上了。} \\
\hline
Négation & « Il n'a \emph{pas} fermé la porte. » & \zh{他没有把门关上。} (négation avant \zh{把}) \\
\hline
\end{tabularx}
\end{center}

\begin{attention}[Les quatre erreurs les plus fréquentes des francophones]
\begin{enumerate}
\item \textbf{Laisser l'objet après le verbe} : \zh{我把吃完了饭} est un désordre total ; l'objet suit \zh{把} et précède le verbe : \zh{我把饭吃完了}.
\item \textbf{Oublier le complément} : \zh{我把衣服洗} est fautif ; dire \zh{我把衣服洗了} ou \zh{我把衣服洗干净了}.
\item \textbf{Placer la négation après 把} : \zh{我把作业没做完} est fautif ; dire \zh{我没把作业做完}.
\item \textbf{Utiliser 把 avec un verbe d'état} : \zh{我把汉语喜欢} est impossible ; dire \zh{我喜欢汉语} (\emph{Wǒ xǐhuan Hànyǔ}, « j'aime le chinois »).
\end{enumerate}
\end{attention}

\begin{methode}[Construire une phrase en 把 en cinq pas]
\begin{enumerate}
\item Identifie le sujet, l'objet \emph{défini} et l'action : « j'ai fini mes devoirs » → sujet \zh{我}, objet \zh{作业}, action \zh{做}.
\item Vérifie que le verbe est un verbe d'action (sinon, pas de \zh{把}).
\item Place l'objet après \zh{把} : \zh{我把作业}.
\item Ajoute le verbe suivi d'un complément ou de \zh{了} : \zh{做完了}.
\item Assemble : \zh{我把作业做完了。} — puis vérifie : négation éventuelle \emph{avant} \zh{把}.
\end{enumerate}
\end{methode}

\coursec{Entraînement : application et transformation}

\begin{exoresolu}[Transformation SVO → 把]
Énoncé : transforme en phrase en \zh{把} : « Il a fini de boire l'eau. » (phrase SVO : \zh{他喝完水了。})
\tcblower
Solution : l'objet défini est \zh{水} (l'eau), le verbe d'action est \zh{喝}, le complément de résultat est \zh{完}. On place \zh{把水} avant le verbe : \zh{他把水喝完了。} \emph{Tā bǎ shuǐ hē wán le.} « Il a fini de boire l'eau. » Vérification : objet défini, verbe d'action, complément présent, \zh{了} final : la phrase est correcte.
\end{exoresolu}

\begin{exercice}[Transforme en phrase en 把]
Transforme les phrases suivantes : (1) \zh{我做完作业了。} (2) \zh{他关窗户了。} (3) \zh{妈妈洗衣服了。} (4) \zh{请开门。}
\end{exercice}

\begin{corrige}[Exercice « Transforme en phrase en 把 »]
(1) \zh{我把作业做完了。} \emph{Wǒ bǎ zuòyè zuò wán le.} (2) \zh{他把窗户关上了。} \emph{Tā bǎ chuānghu guān shang le.} (3) \zh{妈妈把衣服洗干净了。} \emph{Māma bǎ yīfu xǐ gānjìng le.} (4) \zh{请把门打开。} \emph{Qǐng bǎ mén dǎkāi.}
\end{corrige}

\begin{exercice}[Corrige les phrases fautives]
Chaque phrase contient une erreur ; corrige-la : (1) \zh{我把饭吃。} (2) \zh{他把作业没做完。} (3) \zh{我把这首歌喜欢。} (4) \zh{她把放书在桌子上。}
\end{exercice}

\begin{corrige}[Exercice « Corrige les phrases fautives »]
(1) Verbe seul interdit : \zh{我把饭吃了。} ou \zh{我把饭吃完了。} (2) La négation précède \zh{把} : \zh{他没把作业做完。} (3) \zh{喜欢} est un verbe de sentiment, incompatible avec \zh{把} : \zh{我喜欢这首歌。} (4) L'objet suit \zh{把} et précède le verbe : \zh{她把书放在桌子上。}
\end{corrige}

\begin{exercice}[Traduis en chinois avec 把]
(1) J'ai envoyé la lettre. (2) Ne ferme pas la fenêtre ! (3) Il a déjà fini ses devoirs. (4) Pose le cahier sur la table, s'il te plaît.
\end{exercice}

\begin{corrige}[Exercice « Traduis en chinois avec 把 »]
(1) \zh{我把信寄了。} \emph{Wǒ bǎ xìn jì le.} (2) \zh{别把窗户关上！} \emph{Bié bǎ chuānghu guān shang!} (3) \zh{他已经把作业做完了。} \emph{Tā yǐjing bǎ zuòyè zuò wán le.} (4) \zh{请把本子放在桌子上。} \emph{Qǐng bǎ běnzi fàng zài zhuōzi shang.}
\end{corrige}

\begin{experience}[Jeu de rôle : le rangement de la salle]
Par groupes de trois : un élève joue le professeur et donne cinq ordres avec \zh{把} (\zh{请把书放在桌子上}, \zh{把灯关上}, \zh{把窗户打开}\ldots), un deuxième élève mime les actions, un troisième décrit ensuite ce qui a été fait au passé avec \zh{把} + \zh{了} (\zh{他把书放在桌子上了}). Changez ensuite les rôles. Objectif : cinq phrases correctes chacun, avec un verbe toujours accompagné d'un complément.
\end{experience}

\begin{aretenir}[Bilan de la leçon]
\begin{itemize}
\item Schéma : \textbf{Sujet + 把 + Objet + Verbe + Complément / 了}.
\item \zh{把} déplace l'objet \emph{défini} devant le verbe pour insister sur ce qu'il subit (déplacement, transformation, résultat).
\item Le verbe n'est jamais seul ; la négation et les adverbes se placent \emph{avant} \zh{把}.
\item Pas de \zh{把} avec les verbes d'état ou de sentiment (\zh{是}, \zh{有}, \zh{喜欢}, \zh{知道}\ldots).
\item Le français n'a pas d'équivalent : il faut penser « prendre l'objet en main et lui faire subir l'action ».
\end{itemize}
\end{aretenir}

\coursec{Schéma de la structure}

\noindent Après avoir \textbf{observé} dans la section précédente comment la structure \zh{把} permet de déplacer l'objet devant le verbe pour insister sur son sort, nous allons maintenant \textbf{formaliser} son architecture interne. Comprendre ce « squelette » est la clé pour ne plus jamais produire de phrases agrammaticales comme \zh{*我把书放} ou \zh{*我把一本书买了}.

\courssub{Le schéma canonique : cinq blocs indispensables}

\noindent Toute phrase en \zh{把} correcte et complète respecte l'ordre immuable suivant. Mémorisez ce schéma comme une formule : si un élément manque, la phrase ne fonctionne pas.

\begin{center}
\begin{tabularx}{\linewidth}{|X|X|X|X|}
\hline
\rowcolor{popblueL}
\textbf{Bloc} & \textbf{Contenu} & \textbf{Fonction} & \textbf{Exemple} \\
\hline
1. Sujet & Nom / Pronom & Agent de l'action & \zh{我} (Wǒ) \\
\hline
2. \cle{把} + Objet \textbf{déterminé} & \zh{把} + Nom précis (défini, possessif, démonstratif, contexte connu) & Patient, « ce qui est manipulé » & \zh{把那本书} (bǎ nà běn shū) \\
\hline
3. Verbe & Verbe d'action (souvent transitif) & Action effectuée sur l'objet & \zh{放} (fàng) \\
\hline
4. \textbf{Complément obligatoire (X)} & \zh{了} / Compl. résultat / Compl. direction / Redoublement / \zh{一}+V & \textbf{Indispensable :} montre l'aboutissement & \zh{在桌子上} (zài zhuōzi shàng) \\
\hline
5. (Autres compléments) & Temps, lieu (si pas dans X), manière & Circonstances & \zh{昨天} (zuótiān) \\
\hline
\end{tabularx}
\end{center}

\noindent \textbf{Phrase type :} \zh{我把那本书放在桌子上了。} (Wǒ bǎ nà běn shū fàng zài zhuōzi shàng le. — J'ai posé ce livre sur la table.)

\begin{popfigure}
\begin{tikzpicture}[
    node distance=0.8cm and 1.2cm,
    block/.style={draw=popblue, thick, rounded corners, fill=popblueL, minimum height=1.2cm, text width=2.8cm, align=center, font=\small},
    arrow/.style={-Stealth, thick, popdark},
    stop/.style={draw=poporange, thick, rounded corners, fill=poporangeL, minimum height=1cm, text width=2.5cm, align=center, font=\small\bfseries},
]
\node[block, align=center] (sujet){1. Sujet\\ \zh{我} (Wǒ)};
\node[block, right=of sujet, align=center] (baobj){2. \cle{把} + Objet déterminé\\ \zh{把那本书} (bǎ nà běn shū)};
\node[block, right=of baobj, align=center] (verbe){3. Verbe\\ \zh{放} (fàng)};
\node[block, right=of verbe, align=center] (comp){4. Complément obligatoire (X)\\ \zh{在桌子上了} (zài zhuōzi shàng le)};
\node[stop, below=of verbe, align=center] (stop){INTERDIT : verbe nu\\ \zh{*我把书放}};

\draw[arrow] (sujet) -- (baobj);
\draw[arrow] (baobj) -- (verbe);
\draw[arrow] (verbe) -- (comp);
\draw[arrow, poporange, dashed] (verbe.south) -- (stop.north);
\end{tikzpicture}
\legende{Schéma canonique de la phrase en \zh{把} : l'objet déterminé glisse entre le sujet et le verbe ; le verbe \textbf{doit} être suivi d'un complément (X). Un verbe nu (flèche pointillée) rend la phrase incorrecte.}
\end{popfigure}

\begin{propriete}[Règle d'or : le verbe est TOUJOURS accompagné]
Dans une phrase en \zh{把}, le verbe \textbf{ne peut jamais être nu} (seul). Il doit obligatoirement être suivi d'un élément (noté X) qui indique le \textbf{résultat}, la \textbf{direction}, l'\textbf{achèvement} ou l'\textbf{étendue} de l'action. Sans X, la phrase est agrammaticale : \zh{*我把门关}, \zh{*他把饭吃}.
\end{propriete}

\courssub{Le complément obligatoire « X » : les cinq formes autorisées}

\noindent Ce « X » n'est pas un complément libre : il appartient à une liste de cinq catégories. Si votre verbe n'est suivi d'aucun de ces cinq éléments, la structure \zh{把} est impossible (il faut alors revenir à l'ordre Sujet–Verbe–Objet simple : \zh{我关门了}).

\begin{center}
\begin{tabularx}{\linewidth}{|X|X|X|}
\hline
\rowcolor{popgreenL}
\textbf{Type de X} & \textbf{Forme chinoise + Pinyin} & \textbf{Exemple complet (Caractères + Pinyin + Traduction)} \\
\hline
1. \cle{了} (achèvement) & Verbe + \zh{了} (le) & \zh{我把作业做了。} (Wǒ bǎ zuòyè zuò le. — J'ai fait les devoirs.) \\
\hline
2. Complément de \textbf{résultat} & V + \zh{完/好/成/干净/错/坏…} & \zh{请把门关好。} (Qǐng bǎ mén guān hǎo. — Ferme bien la porte.) \\
\hline
3. Complément de \textbf{direction / destination} & V + \zh{在/到/给/回/上/下/进/出…} + Lieu/Personne & \zh{他把钥匙给我了。} (Tā bǎ yàoshi gěi wǒ le. — Il m'a donné la clé.) \\
\hline
4. \textbf{Redoublement} du verbe & V + V (souvent + \zh{了}) & \zh{把屋子扫扫。} (Bǎ wūzi sǎosao. — Balaye un peu la pièce.) \\
\hline
5. \textbf{\zh{一} + Verbe} (action brève) & \zh{一} + V (souvent + \zh{了}) & \zh{让我把这道题想一想。} (Ràng wǒ bǎ zhè dào tí xiǎng yi xiǎng. — Laisse-moi réfléchir un peu à cet exercice.) \\
\hline
\end{tabularx}
\end{center}

\begin{exemplebox}[Variations sur « manger la pomme »]
Observez comment le verbe \zh{吃} (manger) change de complément X :
\begin{itemize}
    \item \zh{我把苹果\cle{吃了}。} (Wǒ bǎ píngguǒ \cle{chī le}. — J'ai mangé la pomme / action accomplie.)
    \item \zh{我把苹果\cle{吃完了}。} (Wǒ bǎ píngguǒ \cle{chī wán le}. — J'ai fini de manger la pomme / résultat : terminée.)
    \item \zh{我把苹果\cle{吃掉了}。} (Wǒ bǎ píngguǒ \cle{chī diào le}. — J'ai mangé toute la pomme / résultat : disparue.)
    \item \zh{我把苹果\cle{吃了吃}。} (Wǒ bǎ píngguǒ \cle{chī le chī}. — J'ai goûté un peu la pomme / action brève.)
\end{itemize}
\emph{Remarque :} \zh{*我把苹果吃} est \textbf{strictement impossible}.
\end{exemplebox}

\courssub{L'objet après \cle{把} : la détermination obligatoire}

\noindent C'est le piège n°1 pour les francophones. En français, « J'ai acheté \textbf{un} livre » (indéfini) est la norme pour présenter un événement nouveau. En chinois, \zh{把} \textbf{exige} que l'objet soit \textbf{identifiable, connu, spécifique} (défini).

\begin{definition}[Objet déterminé (宾语 \zh{bīnyǔ})]
L'objet placé après \zh{把} doit désigner une entité \textbf{connue} du locuteur et de l'interlocuteur (ou présente dans le contexte immédiat). Il est marqué par :
\begin{enumerate}
    \item Un \textbf{démonstratif} : \zh{这} (ce), \zh{那} (ce…-là) → \zh{把这本书} (ce livre-ci).
    \item Un \textbf{possessif} : \zh{我的}, \zh{你的}, \zh{他的} → \zh{把我的书} (mon livre).
    \item Un \textbf{nom propre} ou une personne identifiée : \zh{把小李送到车站} (conduire Xiao Li à la gare).
    \item Un \textbf{nom défini par le contexte} (déjà mentionné) : \zh{书呢？把书放这儿。} (Shū ne? Bǎ shū fàng zhèr. — Et le livre ? Pose le livre ici.)
    \item Un \textbf{nombre + classificateur} seulement \textbf{si} l'objet est identifié : \zh{把那三本书} (ces trois livres-là), mais \textbf{pas} \zh{*把三本书} (trois livres quelconques).
\end{enumerate}
\end{definition}

\begin{attention}[Erreur fréquente des francophones : \zh{*把一本书}]
\cle{Jamais} \zh{*把一本书}, \zh{*把一个苹果}, \zh{*把一些人} au sens de « un livre quelconque / une pomme quelconque / des gens ».
\begin{itemize}
    \item \textbf{Incorrect :} \zh{*我把一本书买了。} (pour dire : « J'ai acheté un livre. »)
    \item \textbf{Correct (ordre SVO) :} \zh{我\cle{买了一本书}。} (Wǒ \cle{mǎi le yī běn shū}. — J'ai acheté un livre.)
    \item \textbf{Correct (avec \zh{把}) :} \zh{我把\cle{那本书}买了。} (Wǒ bǎ \cle{nà běn shū} mǎi le. — J'ai acheté \textbf{ce} livre / \textbf{le} livre dont on parle.)
\end{itemize}
\emph{Astuce :} si vous pouvez remplacer l'objet par « \textbf{ce} livre / \textbf{mon} livre / \textbf{le} livre », \zh{把} est possible. Si c'est « \textbf{un} livre », utilisez l'ordre SVO.
\end{attention}

\begin{popfigure}
\begin{tikzpicture}[
    node distance=1.2cm and 1.5cm,
    decision/.style={diamond, draw=poppurple, thick, aspect=2, fill=poppurpleL, text width=3cm, align=center, font=\small},
    process/.style={rectangle, draw=popblue, thick, rounded corners, fill=popblueL, text width=3.5cm, align=center, font=\small},
    startend/.style={ellipse, draw=popgreen, thick, fill=popgreenL, text width=2.5cm, align=center, font=\small},
    arrow/.style={-Stealth, thick, popdark},
    yesno/.style={font=\small, pos=0.5, sloped, above}
]
\node[startend] (start) {Début : je veux utiliser \zh{把}};
\node[decision, below=of start, align=center] (det){L'objet est-il \textbf{déterminé} ?\\(connu, précis, démonstratif, possessif)};
\node[process, left=of det, align=center] (no_det){NON\\(indéfini : un, une, des)};
\node[process, right=of det] (yes_det) {OUI};
\node[decision, below=of yes_det, align=center] (comp){Le verbe a-t-il un \textbf{complément X} ?\\(\zh{了}, résultat, direction, redoublement, \zh{一}+V)};
\node[process, left=of comp, align=center] (no_comp){NON\\(verbe nu)};
\node[process, right=of comp] (yes_comp) {OUI};
\node[startend, below=of yes_comp] (ok) {\textbf{PHRASE EN \zh{把} CORRECTE}};
\node[process, below=of no_det, align=center] (err1){\zh{*把一本书…}\\ $\rightarrow$ Ordre SVO : \zh{我买了一本书}};
\node[process, below=of no_comp, align=center] (err2){\zh{*我把书放}\\ $\rightarrow$ Ajoutez X : \zh{我把书放好了}};

\draw[arrow] (start) -- (det);
\draw[arrow] (det) -- node[yesno] {NON} (no_det);
\draw[arrow] (det) -- node[yesno] {OUI} (yes_det);
\draw[arrow] (yes_det) -- (comp);
\draw[arrow] (comp) -- node[yesno] {NON} (no_comp);
\draw[arrow] (comp) -- node[yesno] {OUI} (yes_comp);
\draw[arrow] (yes_comp) -- (ok);
\draw[arrow, poporange] (no_det) -- (err1);
\draw[arrow, poporange] (no_comp) -- (err2);
\end{tikzpicture}
\legende{Organigramme de décision : deux conditions \textbf{sine qua non} pour la phrase en \zh{把} : 1. objet déterminé ; 2. verbe accompagné (X).}
\end{popfigure}

\courssub{Place de la négation et des auxiliaires : AVANT \cle{把}}

\noindent En français, la négation entoure le verbe (\emph{je n'ai \textbf{pas} mis}). En chinois, la négation (et les auxiliaires modaux) se place \textbf{systématiquement AVANT \zh{把}}, jamais entre \zh{把} et le verbe. C'est une règle d'ordre des mots \textbf{absolue}.

\begin{propriete}[Place de la négation (\zh{没} / \zh{不} / \zh{别}) et des auxiliaires (\zh{要} / \zh{想} / \zh{会} / \zh{应该} / \zh{能} / \zh{可以})]
\textbf{Structure :} Sujet + (Négation / Auxiliaire) + \zh{把} + Objet + Verbe + X.
\begin{itemize}
    \item \textbf{Négation du passé (\zh{没}) :} \zh{我\cle{没}把作业做完。} (Wǒ \cle{méi} bǎ zuòyè zuò wán. — Je n'ai pas fini les devoirs.)
    \item \textbf{Négation présente / future (\zh{不}) :} \zh{他\cle{不}把实话告诉我们。} (Tā \cle{bù} bǎ shíhuà gàosu wǒmen. — Il ne nous dit pas la vérité.)
    \item \textbf{Impératif négatif (\zh{别}) :} \zh{\cle{别}把门关上！} (\cle{Bié} bǎ mén guān shàng! — Ne ferme pas la porte !)
    \item \textbf{Auxiliaires :} \zh{你\cle{应该}把房间打扫干净。} (Nǐ \cle{yīnggāi} bǎ fángjiān dǎsǎo gānjìng. — Tu devrais nettoyer la chambre à fond.)
\end{itemize}
\cle{Interdiction formelle :} \zh{*把不…}, \zh{*把没…}, \zh{*把别…}, \zh{*把应该…} + Verbe.
\end{propriete}

\begin{exemplebox}[Négation et auxiliaires en contexte]
\begin{itemize}
    \item \zh{小张\cle{没}把钥匙带来。} (Xiǎo Zhāng \cle{méi} bǎ yàoshi dài lái. — Xiao Zhang n'a pas apporté la clé.)
    \item \zh{我们\cle{要}把这个问题解决掉。} (Wǒmen \cle{yào} bǎ zhège wèntí jiějué diào. — Nous devons résoudre ce problème.)
    \item \zh{请\cle{别}把手机放在桌子上。} (Qǐng \cle{bié} bǎ shǒujī fàng zài zhuōzi shàng. — Merci de ne pas poser le téléphone sur la table.)
\end{itemize}
\end{exemplebox}

\courssub{Place des adverbes de manière : avant ou après \cle{把} ?}

\noindent La position de l'adverbe de manière (par exemple \zh{好好地}, \zh{认真地}, \zh{仔细地}) change la portée de la manière exprimée.

\begin{center}
\begin{tabularx}{\linewidth}{|X|X|X|}
\hline
\rowcolor{poppinkL}
\textbf{Position} & \textbf{Sens / Portée} & \textbf{Exemple} \\
\hline
\textbf{Avant \zh{把}} & La manière qualifie \textbf{l'attitude globale} du sujet face à l'action. & \zh{他\cle{好好地}把房间收拾了。} (Tā \cle{hǎohāo de} bǎ fángjiān shōushi le. — Il a rangé la chambre \textbf{comme il faut, sérieusement}.) \\
\hline
\textbf{Après le verbe, avec \zh{得}} & La manière qualifie \textbf{le résultat} obtenu sur l'objet. & \zh{他把房间\cle{收拾得很好}。} (Tā bǎ fángjiān \cle{shōushi de hěn hǎo}. — Il a rangé la chambre \textbf{très bien} / le résultat est bon.) \\
\hline
\end{tabularx}
\end{center}

\begin{aretenir}[Résumé des positions]
\begin{itemize}
    \item \zh{Adv + 把} = « il a fait l'action \textbf{de telle manière} » (attitude du sujet).
    \item \zh{把 + Obj + V + 得 + Adv} = « il a manipulé l'objet \textbf{avec tel résultat} » (manière visible sur l'objet).
    \item \textbf{Conseil :} en classe de Première, privilégiez la position \textbf{avant \zh{把}} pour les adverbes en \zh{地} (\zh{认真地把…}, \zh{好好地把…}) : c'est la plus naturelle et la moins sujette à erreur.
\end{itemize}
\end{aretenir}

\courssub{Méthode : vérifier sa phrase en \cle{把} en trois questions}

\begin{methode}[Check-list « Phrase en 把 : valide ou à refaire ? »]
Avant d'écrire ou de dire une phrase en \zh{把}, posez-vous ces trois questions dans l'ordre. Si vous répondez \textbf{OUI} aux trois, votre phrase est grammaticale.
\begin{enumerate}
    \item \textbf{L'objet est-il DÉTERMINÉ ?} (Est-ce \zh{这/那/我的}, ou le livre dont on parle ?) $\rightarrow$ Si NON : \textbf{interdit}. Utilisez l'ordre SVO (\zh{我买了一本书}).
    \item \textbf{Le verbe a-t-il un COMPLÉMENT X ?} (Y a-t-il \zh{了}, \zh{完}, \zh{好}, \zh{给}, \zh{在}…, un redoublement ou \zh{一}+V ?) $\rightarrow$ Si NON : \textbf{interdit}. Ajoutez X ou utilisez l'ordre SVO.
    \item \textbf{La négation / l'auxiliaire est-il AVANT \zh{把} ?} (Ai-je placé \zh{没/不/别/应该/要} \textbf{devant} \zh{把} ?) $\rightarrow$ Si NON (par exemple \zh{*把没做完}) : \textbf{incorrect}. Déplacez-le avant \zh{把}.
\end{enumerate}
\emph{Exemple d'application :} je veux dire « Je n'ai pas apporté mon passeport. »
\begin{itemize}
    \item Objet : \zh{我的护照} (possessif $\rightarrow$ déterminé \checkmark).
    \item Verbe : \zh{带来} (verbe + direction \zh{来} $\rightarrow$ X présent \checkmark).
    \item Négation : \zh{没} placée avant \zh{把} $\rightarrow$ \zh{我\cle{没}把护照带来。} (Wǒ méi bǎ hùzhào dài lái.) \checkmark
\end{itemize}
\end{methode}

\begin{exoresolu}[Transformation et correction : de l'ordre SVO à \zh{把} (et inversement)]
\textbf{Consigne :} transformez les phrases SVO en phrases en \zh{把} si possible. Sinon, expliquez pourquoi (objet indéfini ou verbe nu) et corrigez la phrase en \zh{把} proposée.

\begin{enumerate}
    \item \zh{我关上了窗户。} (Wǒ guān shàng le chuānghu.)
    \item \zh{他买了一辆车。} (Tā mǎi le yī liàng chē.)
    \item \zh{*请把名字写。} (phrase incorrecte à corriger.)
    \item \zh{*我们把应该这件事忘记。} (phrase incorrecte à corriger : auxiliaire mal placé et verbe nu.)
    \item \zh{别把那个坏人放跑了！} (vérifiez si elle est correcte.)
\end{enumerate}

\tcblower

\textbf{Solutions détaillées :}

\begin{enumerate}
    \item \textbf{Transformation possible.}
        \begin{itemize}
            \item Objet : \zh{窗户} $\rightarrow$ dans le contexte, « la fenêtre (de la pièce) » $\rightarrow$ déterminé. Ajoutons \zh{那扇} pour être explicite : \zh{那扇窗户}.
            \item Verbe : \zh{关上} (verbe + complément de direction \zh{上} = X présent).
            \item Résultat : \zh{我把\cle{那扇窗户}关上了。} (Wǒ bǎ \cle{nà shàn chuānghu} guān shàng le. — J'ai fermé cette fenêtre.)
        \end{itemize}
    \item \textbf{Transformation IMPOSSIBLE.}
        \begin{itemize}
            \item Objet : \zh{一辆车} (une voiture) = \textbf{indéfini}. \zh{把} ne peut pas introduire un objet nouveau ou indéfini.
            \item On garde l'ordre SVO : \zh{他买了一辆车。} (Tā mǎi le yī liàng chē. — Il a acheté une voiture.) Si la voiture devient connue plus tard dans le discours, on pourra dire : \zh{他把\cle{那辆车}买了。} (Il a acheté cette voiture-là.)
        \end{itemize}
    \item \textbf{Correction : verbe nu.}
        \begin{itemize}
            \item \zh{写} (écrire) est nu : il faut un complément X.
            \item Corrections possibles : \zh{请把名字\cle{写上}。} (Qǐng bǎ míngzi \cle{xiě shàng}. — Inscris ton nom.) / \zh{请把名字\cle{写好}。} (\cle{xiě hǎo} — écris bien ton nom.) / \zh{请把名字\cle{写一遍}。} (\cle{xiě yī biàn} — écris le nom une fois.)
        \end{itemize}
    \item \textbf{Double correction : auxiliaire mal placé + verbe nu.}
        \begin{itemize}
            \item Erreur 1 : l'auxiliaire \zh{应该} doit se placer \textbf{avant} \zh{把}, jamais après.
            \item Erreur 2 : \zh{忘记} est nu ; il faut un complément de résultat, ici \zh{掉} (\zh{忘掉} = oublier complètement).
            \item Correction : \zh{我们\cle{应该把}这件事\cle{忘掉}。} (Wǒmen \cle{yīnggāi bǎ} zhè jiàn shì \cle{wàng diào}. — Nous devrions oublier complètement cette affaire.)
            \item Variante naturelle en ordre SVO : \zh{我们应该忘记这件事。} (Wǒmen yīnggāi wàngjì zhè jiàn shì.)
        \end{itemize}
    \item \textbf{Phrase CORRECTE.}
        \begin{itemize}
            \item Objet : \zh{那个坏人} (démonstratif \zh{那个} $\rightarrow$ déterminé \checkmark).
            \item Verbe : \zh{放跑} (verbe + complément de résultat \zh{跑} = X \checkmark).
            \item Impératif négatif : \zh{别} placé avant \zh{把} \checkmark.
            \item Traduction : « Ne laisse pas ce bandit s'enfuir ! »
        \end{itemize}
\end{enumerate}
\end{exoresolu}

\coursec{Leçon 28 — Grammaire : 把字句}

\courssub{Observation guidée d'exemples}

Avant d'énoncer la règle, comparons deux façons de dire « j'ai fermé la porte » en chinois :

\begin{exemplebox}[SVO vs 把]
\zh{我关门了。} \emph{Wǒ guān mén le.} (Ordre SVO : Sujet + Verbe + Objet)\\
\zh{我把门关上了。} \emph{Wǒ bǎ mén guān shàng le.} (Ordre 把 : Sujet + 把 + Objet + Verbe + Complément)
\end{exemplebox}

\begin{aretenir}[Ce que l'observation révèle]
\begin{itemize}
\item Les deux phrases sont correctes, mais leur \cle{point de vue} diffère. \zh{我关门了} rapporte simplement l'action (« j'ai fait l'action de fermer la porte »). \zh{我把门关上了} se focalise sur \textbf{ce qui est arrivé à la porte} (elle est maintenant fermée).
\item La phrase en \zh{把} \textbf{exige un complément après le verbe} (\zh{上}, \zh{好}, \zh{在桌子上}, etc.). Un verbe nu après \zh{把} est impossible.
\item L'objet de \zh{把} (\zh{门}) est \textbf{défini} : on sait de quelle porte il s'agit.
\end{itemize}
\end{aretenir}

\courssub{Schéma de la structure}

\begin{propriete}[Structure fondamentale de la phrase en 把]
\textbf{Formule :} \cle{Sujet + 把 + Objet + Verbe + Complément (obligatoire) + (了)}\\
\textbf{Compléments obligatoires fréquents :}
\begin{itemize}
\item Complément de \textbf{résultat} : \zh{好、完、错、干净、清楚、坏…}
\item Complément de \textbf{direction/lieu} : \zh{在/到/进/上/下/出来/回去 + Lieu}
\item Complément de \textbf{transfert} : \zh{给 + Personne}
\item Complément de \textbf{transformation} : \zh{成/做 + Résultat}
\item Complément de \textbf{quantité d'action} : \zh{一遍、两次、半小时…}
\end{itemize}
\textbf{Règle d'or :} \emph{Verbe nu = pas de phrase en 把.}
\end{propriete}

\coursec{Emplois et cas d'usage}

Dans la section précédente, nous avons établi le schéma général de la phrase en \zh{把} : \textbf{Sujet + 把 + Objet + Verbe + X}, où X est un complément obligatoire (résultat, direction, lieu, quantité…). Nous savons donc \emph{comment} construire la structure. La question qui se pose maintenant est : \emph{quand} l'utiliser ? En effet, le chinois possède aussi l'ordre simple Sujet–Verbe–Objet (SVO), comme le français. Pourquoi dire \zh{我把门关上了} plutôt que \zh{我关门了} ? Cette section répond à cette question en présentant les \cle{six grands emplois} de la phrase en \zh{把}, chacun illustré par des exemples de la vie quotidienne, au Cameroun comme en Chine.

\courssub{Vue d'ensemble : les six emplois de 把}

Avant d'entrer dans le détail, observons la carte mentale suivante, qui résume toute la section. Chaque branche correspond à un emploi, avec un mini-exemple.

\begin{popfigure}
\begin{tikzpicture}[mindmap, grow cyclic, every node/.style={concept, align=center, font=\small},
root concept/.append style={concept color=popblue, font=\large\bfseries},
level 1/.append style={level distance=4.6cm, sibling angle=60},
level 1 concept/.append style={concept color=poporangeL}]
\node[root concept, align=center]{把字句\\bǎ zì jù};
\node[concept, concept color=popblueL, align=center]{Déplacement\\把+O+V+在/到\\我把书放在桌上};
\node[concept, concept color=popgreenL, align=center]{Transfert\\把+O+给+personne\\老师把书发给我们};
\node[concept, concept color=poppinkL, align=center]{Résultat\\把+O+V+错/完/好\\我把名字写错了};
\node[concept, concept color=popgoldL, align=center]{Transformation\\把+O+V+成/做\\把句子翻译成法语};
\node[concept, concept color=poppurpleL, align=center]{Ordre poli\\请把…\\请把窗户关上};
\node[concept, concept color=poporangeL, align=center]{Traitement subi\\风把树吹倒了};
\end{tikzpicture}
\legende{Carte mentale des six emplois de la phrase en 把}
\end{popfigure}

\begin{aretenir}[L'idée directrice]
La phrase en \zh{把} répond toujours à la question : \textbf{« qu'est-il arrivé à l'objet ? »}. L'objet est \cle{défini} (on sait de quoi on parle) et il \cle{subit} une action dont on précise le résultat, la destination ou la manière. Si l'on veut simplement raconter une action, on garde l'ordre SVO.
\end{aretenir}

\courssub{Emploi 1 : déplacement et localisation — 把 + objet + V + 在/到 + lieu}

\textbf{Observation.} Comparez les deux phrases suivantes :

\zh{我把自行车放在学校门口。} \emph{Wǒ bǎ zìxíngchē fàng zài xuéxiào ménkǒu.} « J'ai laissé mon vélo devant le portail de l'école. »

\zh{他把水倒进杯子里。} \emph{Tā bǎ shuǐ dào jìn bēizi lǐ.} « Il a versé l'eau dans le verre. »

Dans les deux cas, l'objet (\zh{自行车} le vélo, \zh{水} l'eau) \emph{change de place} : il aboutit quelque part. C'est exactement le premier grand emploi de \zh{把}.

\begin{propriete}[Emploi 1 : déplacement de l'objet]
\textbf{Structure :} Sujet + 把 + Objet + Verbe de déplacement + 在/到/进 + Lieu\\
\textbf{Emploi :} on insiste sur le fait que l'objet aboutit dans un lieu précis. Les verbes typiques sont \zh{放} (poser), \zh{送} (envoyer, livrer), \zh{带} (apporter), \zh{搬} (déménager), \zh{倒} (verser), \zh{挂} (accrocher).\\
\textbf{Attention :} le verbe seul ne suffit jamais ; il faut \zh{在} + lieu (position) ou \zh{到}/\zh{进} + lieu (mouvement).
\end{propriete}

\begin{exemplebox}[Exemples — déplacement]
\zh{妈妈把菜放进冰箱里。} \emph{Māma bǎ cài fàng jìn bīngxiāng lǐ.} « Maman a mis les légumes dans le réfrigérateur. »\\
\zh{请把行李放到车上。} \emph{Qǐng bǎ xíngli fàng dào chē shàng.} « Mets les bagages dans la voiture, s'il te plaît. »\\
\zh{他把信寄到中国去了。} \emph{Tā bǎ xìn jì dào Zhōngguó qù le.} « Il a envoyé la lettre en Chine. »\\
\zh{我们把桌子搬到教室里。} \emph{Wǒmen bǎ zhuōzi bān dào jiàoshì lǐ.} « Nous avons déplacé la table dans la salle de classe. »
\end{exemplebox}

\begin{attention}[Piège des francophones]
En français, on dit « j'ai posé le livre » sans préciser où ; en chinois, \zh{我把书放了} est \textbf{incorrect} : la phrase en \zh{把} exige un complément. Il faut dire \zh{我把书放在桌子上} (je l'ai posé sur la table) ou, si l'on ne veut rien préciser, revenir à l'ordre SVO : \zh{我放了书} reste d'ailleurs peu naturel ; mieux vaut \zh{我放了一本书} avec un objet indéfini.
\end{attention}

\courssub{Emploi 2 : transfert à une personne — 把 + objet + 给 + personne}

\textbf{Observation.} \zh{老师把新书发给我们。} \emph{Lǎoshī bǎ xīn shū fā gěi wǒmen.} « Le professeur nous a distribué les nouveaux livres. » Ici, l'objet (\zh{新书}) ne change pas seulement de place : il \emph{change de possesseur}. Il passe des mains du professeur dans celles des élèves.

\begin{propriete}[Emploi 2 : transfert]
\textbf{Structure :} Sujet + 把 + Objet + Verbe + 给 + Personne\\
\textbf{Emploi :} l'objet est transmis, offert, vendu, prêté à quelqu'un. Verbes typiques : \zh{给} (donner), \zh{送} (offrir), \zh{发} (distribuer), \zh{卖} (vendre), \zh{借} (prêter), \zh{告诉} (dire, transmettre une information), \zh{介绍} (présenter).\\
\textbf{Cas particulier :} avec \zh{告诉} et \zh{介绍}, le transfert est abstrait : on « transmet » une information.
\end{propriete}

\begin{exemplebox}[Exemples — transfert]
\zh{我把这个消息告诉了妈妈。} \emph{Wǒ bǎ zhège xiāoxi gàosu le māma.} « J'ai annoncé cette nouvelle à maman. »\\
\zh{他把相机借给我了。} \emph{Tā bǎ zhàoxiàngjī jiè gěi wǒ le.} « Il m'a prêté l'appareil photo. »\\
\zh{请把这本书交给王老师。} \emph{Qǐng bǎ zhè běn shū jiāo gěi Wáng lǎoshī.} « Remets ce livre au professeur Wang, s'il te plaît. »
\end{exemplebox}

\courssub{Emploi 3 : résultat obtenu sur l'objet — 把 + objet + V + complément de résultat}

\textbf{Observation.} \zh{我把名字写错了。} \emph{Wǒ bǎ míngzi xiě cuò le.} « J'ai mal écrit le nom. » Le verbe \zh{写} (écrire) est suivi du complément de résultat \zh{错} (faux) : l'action est achevée et l'objet porte la marque du résultat. C'est l'emploi le plus fréquent de tous.

\begin{propriete}[Emploi 3 : résultat]
\textbf{Structure :} Sujet + 把 + Objet + Verbe + Complément de résultat (+ 了)\\
\textbf{Emploi :} on insiste sur le résultat concret subi par l'objet. Compléments fréquents : \zh{错} (faux), \zh{对} (juste), \zh{完} (fini), \zh{好} (bien fait), \zh{坏} (abîmé), \zh{干净} (propre), \zh{清楚} (clair).
\end{propriete}

\begin{exemplebox}[Exemples — résultat]
\zh{我把相机弄坏了。} \emph{Wǒ bǎ zhàoxiàngjī nòng huài le.} « J'ai cassé l'appareil photo. »\\
\zh{弟弟把果汁喝完了。} \emph{Dìdi bǎ guǒzhī hē wán le.} « Le petit frère a fini le jus de fruit. »\\
\zh{她把衣服洗干净了。} \emph{Tā bǎ yīfu xǐ gānjìng le.} « Elle a lavé les vêtements (ils sont propres). »\\
\zh{请你把话说清楚。} \emph{Qǐng nǐ bǎ huà shuō qīngchu.} « Exprime-toi clairement, s'il te plaît. »
\end{exemplebox}

\courssub{Emploi 4 : changement et transformation — 把 + objet + V + 成/做}

\textbf{Observation.} \zh{请把这句话翻译成法语。} \emph{Qǐng bǎ zhè jù huà fānyì chéng Fǎyǔ.} « Traduis cette phrase en français, s'il te plaît. » L'objet (\zh{这句话}) ne bouge pas, il \emph{change de nature} : de chinois, il devient français.

\begin{propriete}[Emploi 4 : transformation]
\textbf{Structure :} Sujet + 把 + Objet + Verbe + 成/做 + Résultat\\
\textbf{Emploi :} l'objet est transformé en quelque chose d'autre. Verbes typiques : \zh{翻译} (traduire), \zh{改} (changer), \zh{变} (transformer), \zh{当} (prendre pour). \zh{成} marque le devenir, \zh{做} marque la fonction ou l'usage.
\end{propriete}

\begin{exemplebox}[Exemples — transformation]
\zh{我们把教室改成了图书室。} \emph{Wǒmen bǎ jiàoshì gǎi chéng le túshūshì.} « Nous avons transformé la salle de classe en bibliothèque. »\\
\zh{他把这个词当成动词了。} \emph{Tā bǎ zhège cí dàng chéng dòngcí le.} « Il a pris ce mot pour un verbe. »\\
\zh{请把这篇文章翻译成中文。} \emph{Qǐng bǎ zhè piān wénzhāng fānyì chéng Zhōngwén.} « Traduis cet article en chinois, s'il te plaît. »
\end{exemplebox}

\courssub{Emploi 5 : ordres et demandes polies — 请把…}

\textbf{Observation.} \zh{请把窗户关上。} \emph{Qǐng bǎ chuānghu guān shàng.} « Ferme la fenêtre, s'il te plaît. » La phrase en \zh{把} est la forme la plus naturelle de l'ordre poli en chinois, parce qu'un ordre porte toujours sur un objet précis dont on attend un résultat précis.

\begin{propriete}[Emploi 5 : demande polie]
\textbf{Structure :} 请 + (Personne) + 把 + Objet + Verbe + Complément\\
\textbf{Emploi :} demander à quelqu'un d'accomplir une action sur un objet déterminé. Très fréquent en classe : le professeur dit \zh{请把书打开} (ouvrez vos livres), \zh{请把作业交上来} (rendez vos devoirs).
\end{propriete}

\begin{exemplebox}[Exemples — demandes polies]
\zh{请把这篇课文读一遍。} \emph{Qǐng bǎ zhè piān kèwén dú yí biàn.} « Lis ce texte une fois, s'il te plaît. »\\
\zh{请把门打开。} \emph{Qǐng bǎ mén dǎkāi.} « Ouvre la porte, s'il te plaît. »\\
\zh{请把你的名字写在黑板上。} \emph{Qǐng bǎ nǐ de míngzi xiě zài hēibǎn shàng.} « Écris ton nom au tableau, s'il te plaît. »
\end{exemplebox}

\begin{aretenir}[Complément de quantité = complément valide]
Remarquez le premier exemple : \zh{读一遍} contient un complément de quantité d'action (\zh{一遍}, « une fois »). Ce complément suffit à rendre la phrase en \zh{把} correcte, sans résultat ni lieu.
\end{aretenir}

\courssub{Emploi 6 : insister sur le traitement subi par l'objet}

\textbf{Observation.} \zh{风把树吹倒了。} \emph{Fēng bǎ shù chuī dǎo le.} « Le vent a renversé l'arbre. » Ici, le sujet (\zh{风}, le vent) n'est pas une personne : c'est une force. La phrase en \zh{把} met en avant ce que l'objet a \emph{subi}. Cet emploi est très fréquent dans les récits, les nouvelles et la presse.

\begin{exemplebox}[Exemples — traitement subi]
\zh{大雨把路淹了。} \emph{Dàyǔ bǎ lù yān le.} « La grosse pluie a inondé la route. »\\
\zh{孩子把杯子打破了。} \emph{Háizi bǎ bēizi dǎ pò le.} « L'enfant a cassé le verre. »\\
\zh{这个消息把大家高兴坏了。} \emph{Zhège xiāoxi bǎ dàjiā gāoxìng huài le.} « Cette nouvelle a rendu tout le monde fou de joie. »
\end{exemplebox}

\courssub{Texte d'application : la famille prépare le repas}

\begin{textebox}[晚饭以前 — Avant le dîner]
今天是星期天，我们全家都在家。妈妈把鱼做好了，爸爸把桌子擦干净了，弟弟把果汁喝完了！妈妈有点生气，她把弟弟批评了一顿。后来，弟弟把杯子洗干净了，又把椅子搬到桌子旁边。爸爸把电视关上了，说：「请把手机放下，我们吃饭吧！」我把筷子放在桌子上，姐姐把菜端上来了。我们一家人高高兴兴地吃了晚饭。\\
\emph{Jīntiān shì xīngqītiān, wǒmen quánjiā dōu zài jiā. Māma bǎ yú zuò hǎo le, bàba bǎ zhuōzi cā gānjìng le, dìdi bǎ guǒzhī hē wán le! Māma yǒudiǎn shēngqì, tā bǎ dìdi pīpíng le yí dùn. Hòulái, dìdi bǎ bēizi xǐ gānjìng le, yòu bǎ yǐzi bān dào zhuōzi pángbiān. Bàba bǎ diànshì guān shàng le, shuō : « Qǐng bǎ shǒujī fàng xià, wǒmen chīfàn ba ! » Wǒ bǎ kuàizi fàng zài zhuōzi shàng, jiějie bǎ cài duān shànglái le. Wǒmen yìjiā rén gāogāo-xìngxìng de chī le wǎnfàn.}\\
« Aujourd'hui, c'est dimanche, toute la famille est à la maison. Maman a fini de préparer le poisson, papa a essuyé la table (elle est propre), et le petit frère a fini tout le jus de fruit ! Maman est un peu fâchée, elle l'a grondé. Ensuite, le petit frère a lavé le verre et a déplacé la chaise près de la table. Papa a éteint la télévision et a dit : "Posez vos téléphones, passons à table !" J'ai posé les baguettes sur la table, ma grande sœur a apporté les plats. Toute la famille a dîné joyeusement. »
\end{textebox}

\begin{center}
\begin{tabularx}{\linewidth}{|X|X|X|X|}
\hline
\rowcolor{popblueL} Caractères & Pinyin & Nature & Traduction \\
\hline
擦 & cā & verbe & essuyer, frotter \\
\hline
批评 & pīpíng & verbe & critiquer, gronder \\
\hline
一顿 & yí dùn & complément & une bonne fois (action verbale) \\
\hline
搬 & bān & verbe & déplacer, porter (objet lourd) \\
\hline
端 & duān & verbe & apporter (un plat, à deux mains) \\
\hline
筷子 & kuàizi & nom & baguettes \\
\hline
旁边 & pángbiān & nom & côté, à côté \\
\hline
\end{tabularx}
\end{center}

\textbf{Questions de compréhension :}
\begin{enumerate}
\item Combien de phrases en \zh{把} comptez-vous dans ce texte ? Soulignez-les.
\item Pour chaque phrase en \zh{把}, identifiez l'emploi (1 à 6) et le complément qui suit le verbe.
\item Pourquoi maman est-elle fâchée ? Traduisez la phrase qui l'explique.
\end{enumerate}

\courssub{Choisir entre SVO et 把 : la règle de décision}

C'est le point le plus délicat pour un francophone. Le français n'a qu'une seule construction (« j'ai fermé la porte »), le chinois en a deux. Voici la méthode pour choisir.

\begin{methode}[SVO ou 把 ? Trois questions à se poser]
\begin{enumerate}
\item \textbf{L'objet est-il défini et connu ?} (la porte, mon vélo, ce livre…) Si non, si l'objet est indéfini (\zh{一本书}, un livre quelconque), la phrase en \zh{把} est impossible : gardez SVO.
\item \textbf{Peut-on répondre à « qu'est-il arrivé à l'objet ? »} (il a été fermé, cassé, déplacé, fini, transformé…) Si oui, la phrase en \zh{把} est naturelle.
\item \textbf{Le verbe a-t-il un complément} (résultat, lieu, \zh{给}, quantité) ? Si le verbe est nu, sans rien après, \zh{把} est impossible : gardez SVO.
\end{enumerate}
Si les trois réponses sont favorables, employez \zh{把}. Sinon, SVO.
\end{methode}

\begin{center}
\begin{tabularx}{\linewidth}{|X|X|X|}
\hline
\rowcolor{popblueL} Situation & Phrase correcte & Analyse \\
\hline
Action simple, objet indéfini & 我昨天买了一本书。Wǒ zuótiān mǎi le yì běn shū. (J'ai acheté un livre hier.) & SVO obligatoire : objet indéfini \\
\hline
Résultat sur objet défini & 我把那本书看完了。Wǒ bǎ nà běn shū kàn wán le. (J'ai fini ce livre.) & 把 : objet défini + résultat 完 \\
\hline
Verbe sans complément & 我喜欢中国菜。Wǒ xǐhuan Zhōngguó cài. (J'aime la cuisine chinoise.) & SVO : verbes d'état (aimer, savoir, avoir) jamais avec 把 \\
\hline
Demande polie & 请把窗户关上。Qǐng bǎ chuānghu guān shàng. (Ferme la fenêtre, s.t.p.) & 把 : ordre sur objet précis \\
\hline
\end{tabularx}
\end{center}

\begin{attention}[Erreurs fréquentes des francophones]
\begin{itemize}
\item \textbf{Verbe nu après 把 :} \zh{我把门关} est faux. Le français « j'ai fermé la porte » se traduit par \zh{我把门关上了} ou \zh{我把门关好了} : il faut un complément.
\item \textbf{Objet indéfini :} on ne peut pas dire \zh{我把一本书买了}. L'objet de \zh{把} doit être connu des deux interlocuteurs : \zh{我把那本书买了}.
\item \textbf{Verbes d'état :} \zh{喜欢} (aimer), \zh{知道} (savoir), \zh{有} (avoir), \zh{是} (être) ne prennent jamais \zh{把}, car rien n'« arrive » à l'objet.
\item \textbf{Place de la négation :} elle se met \emph{avant} \zh{把}, jamais avant le verbe : \zh{我没有把名字写错} (je n'ai pas mal écrit le nom), et non \zh{我把名字没有写错}.
\end{itemize}
\end{attention}

\begin{experience}[Mise en situation orale : « À la maison / Au restaurant »]
\textbf{Consigne :} En binôme, jouez un dialogue de 1 minute. L'un joue le rôle d'un parent / serveur, l'autre celui d'un enfant / client. Utilisez au moins \textbf{4 phrases en 把} couvrant 3 emplois différents (ordre, déplacement, résultat, transfert).

\textbf{Exemple de début :}
A: \zh{请把筷子放在桌子上。} (Emploi 1 : déplacement)
B: \zh{好。我把汤喝完了。} (Emploi 3 : résultat)
A: \zh{把碗给我吧。} (Emploi 2 : transfert — \zh{给} omis mais sous-entendu, structure \zh{把碗给我} = \zh{把碗给我} / \zh{把碗递给我})
B: \zh{给。请把椅子搬过来。} (Emploi 1 : déplacement)

\textbf{Objectif :} Automatiser l'ordre Sujet-\zh{把}-Objet-Verbe-Complément et le choix du complément adapté.
\end{experience}

\begin{exoresolu}[Transformation SVO → 把]
Énoncé : transformez la phrase SVO suivante en phrase en \zh{把}, en ajoutant le complément qui convient : \zh{他喝了那杯果汁。} \emph{Tā hē le nà bēi guǒzhī.} « Il a bu ce verre de jus de fruit. »
\tcblower
Solution détaillée :
\begin{enumerate}
\item L'objet \zh{那杯果汁} (ce verre de jus) est \textbf{défini} : la transformation est possible.
\item On se demande « qu'est-il arrivé au jus ? » → il a été \emph{fini} : on ajoute le complément de résultat \zh{完}.
\item On place l'objet après \zh{把}, avant le verbe, et on ajoute \zh{了} en fin de phrase.
\end{enumerate}
Résultat : \zh{他把那杯果汁喝完了。} \emph{Tā bǎ nà bēi guǒzhī hē wán le.} « Il a fini ce verre de jus de fruit. »\\
Vérification : objet défini ✓ ; complément de résultat \zh{完} ✓ ; \zh{了} final ✓. La phrase est correcte.
\end{exoresolu}

\begin{exercice}[À vous de jouer]
\begin{enumerate}
\item Traduisez en chinois avec \zh{把} :
      \begin{enumerate}
      \item « J'ai mis les bagages dans la voiture. »
      \item « Traduis cette phrase en français. »
      \item « Le vent a renversé l'arbre. »
      \item « Lis ce texte une fois, s'il te plaît. »
      \end{enumerate}
\item Transformez en phrase en \zh{把} : \zh{我写了名字。} (ajoutez un complément de résultat).
\item Corrigez la faute : \zh{他把窗户关。}
\item Choisissez : pour dire « j'ai acheté deux pommes », faut-il SVO ou \zh{把} ? Justifiez.
\end{enumerate}
\end{exercice}

\begin{corrige}[Exercice « À vous de jouer »]
\begin{enumerate}
\item a) \zh{我把行李放到车上。} \emph{Wǒ bǎ xíngli fàng dào chē shàng.} — b) \zh{请把这句话翻译成法语。} \emph{Qǐng bǎ zhè jù huà fānyì chéng Fǎyǔ.} — c) \zh{风把树吹倒了。} \emph{Fēng bǎ shù chuī dǎo le.} — d) \zh{请把这篇课文读一遍。} \emph{Qǐng bǎ zhè piān kèwén dú yí biàn.}
\item \zh{我把名字写错了。} \emph{Wǒ bǎ míngzi xiě cuò le.} (ou \zh{写好} / \zh{写完} selon le sens voulu).
\item Il manque un complément : \zh{他把窗户关上了。} \emph{Tā bǎ chuānghu guān shàng le.}
\item SVO : \zh{我买了两个苹果。} \emph{Wǒ mǎi le liǎng ge píngguǒ.} L'objet est indéfini (deux pommes quelconques) et rien de particulier ne leur « arrive » : \zh{把} est impossible.
\end{enumerate}
\end{corrige}

\begin{savaistu}[把 dans la presse et au Bac]
Dans la presse et la conversation chinoises, la phrase en \zh{把} est si fréquente qu'on la retrouve en moyenne plusieurs fois par paragraphe : les Chinois aiment préciser ce qui est arrivé aux choses. Les titres de journaux l'emploient constamment, par exemple pour rapporter un accident ou une décision. Au Bac blanc comme au Bac, les épreuves de thème contiennent presque toujours au moins une phrase qui appelle naturellement la construction en \zh{把} (« il a envoyé la lettre à son ami », « ferme la porte », « elle a fini ses devoirs »). La maîtriser, c'est sécuriser des points précieux.
\end{savaistu}

\coursec{Exceptions et restrictions}

Dans la section précédente, nous avons vu les \cle{emplois} de la phrase en \zh{把} : elle sert à insister sur ce qu'on \emph{fait subir} à un objet précis (le déplacer, le transformer, le terminer). Mais la structure en \zh{把} n'est pas libre : elle obéit à des \cle{restrictions} strictes. Un francophone qui connaît le schéma \zh{把} + objet + verbe + complément peut encore produire des phrases impossibles en chinois, simplement parce qu'il calque le français. Cette section recense les quatre grandes restrictions, un cas particulier célèbre, puis une méthode simple pour décider, avant de parler, si \zh{把} est permis ou non.

\courssub{Restriction 1 : l'objet de 把 doit être défini}

\begin{exemplebox}[Observation : deux phrases, une seule est correcte]
Correct : \zh{我把那封信寄出去了。}\\
\emph{Wǒ bǎ nà fēng xìn jì chūqù le.} — « J'ai posté cette lettre. »\\[0.3em]
Incorrect : *\zh{我把一封信寄出去了。}\\
\emph{*Wǒ bǎ yī fēng xìn jì chūqù le.} — phrase impossible si l'on veut dire « j'ai posté une lettre (quelconque) ».
\end{exemplebox}

\begin{propriete}[Objet défini obligatoire]
L'objet introduit par \zh{把} doit être \cle{défini} ou \cle{identifié} : le locuteur et l'auditeur savent de quoi il s'agit. Il est donc précédé de \zh{这} (ce), \zh{那} (ce...-là), d'un possessif (\zh{我的}, \zh{他的}), d'une relative, ou il est connu du contexte.\\[0.3em]
Interdit : \zh{把} + \zh{一} + classificateur + nom au sens de « un(e) quelconque ».\\[0.3em]
Si l'objet est indéfini, on revient à l'ordre SVO classique :\\
\zh{我吃了一个苹果。} \emph{Wǒ chī le yī gè píngguǒ.} — « J'ai mangé une pomme. » (et non *\zh{我把一个苹果吃了。})
\end{propriete}

Comparaison avec le français : le français dit indifféremment « j'ai mangé la pomme » et « j'ai mangé une pomme ». Le chinois, lui, réserve \zh{把} au premier cas : \zh{把} « attrape » un objet précis pour lui faire subir l'action ; on ne peut pas attraper quelque chose d'indéterminé.

\begin{exemplebox}[Objets définis acceptables après 把]
\zh{请把门关上。} \emph{Qǐng bǎ mén guānshang.} — « Ferme la porte, s'il te plaît. » (la porte, connue du contexte)\\
\zh{他把我的书拿走了。} \emph{Tā bǎ wǒ de shū ná zǒu le.} — « Il a emporté mon livre. » (possessif)\\
\zh{我把昨天买的词典借给同学了。} \emph{Wǒ bǎ zuótiān mǎi de cídiǎn jiè gěi tóngxué le.} — « J'ai prêté à un camarade le dictionnaire acheté hier. » (relative)
\end{exemplebox}

\courssub{Restriction 2 : les verbes d'état et de sentiment refusent 把}

\begin{attention}[Erreur fréquente des francophones : calquer « j'aime ce film »]
Ne traduisez jamais « j'aime ce film » par *\zh{我把这个电影喜欢。} ! Les verbes d'\emph{état}, de \emph{sentiment} et de \emph{connaissance} — \zh{喜欢} (aimer), \zh{爱} (aimer), \zh{知道} (savoir), \zh{认识} (connaître), \zh{想} (penser), \zh{觉得} (trouver que), \zh{看见} au sens simple de « voir » — ne font rien \emph{subir} à l'objet : ils ne le déplacent pas, ne le transforment pas, ne le consomment pas. Ils refusent donc la structure en \zh{把}.\\[0.3em]
Correct : \zh{我喜欢这个电影。} \emph{Wǒ xǐhuan zhè ge diànyǐng.} — « J'aime ce film. »\\
Correct : \zh{我认识他。} \emph{Wǒ rènshi tā.} — « Je le connais. » (et non *\zh{我把他认识。})
\end{attention}

\begin{propriete}[Condition sémantique de 把]
Le verbe d'une phrase en \zh{把} doit exprimer une \cle{action effective} qui produit un \cle{résultat concret} sur l'objet : déplacement (\zh{放} poser, \zh{送} envoyer), transformation (\zh{洗} laver, \zh{翻译} traduire), achèvement (\zh{写} écrire, \zh{做} faire), suppression (\zh{关} fermer, \zh{吃} manger). Un verbe qui ne « touche » pas l'objet ne peut pas prendre \zh{把}.
\end{propriete}

\courssub{Restriction 3 : les verbes de déplacement intransitifs ne prennent pas 把}

Les verbes \zh{来} (venir), \zh{去} (aller), \zh{走} (partir, marcher), \zh{回} (rentrer) sont \cle{intransitifs} : ils n'ont pas d'objet à « manipuler ». On ne peut donc pas les mettre dans une phrase en \zh{把}.

\begin{exemplebox}[Intransitifs : pas de 把]
\zh{他回雅温得了。} \emph{Tā huí Yǎwēndé le.} — « Il est rentré à Yaoundé. » (impossible : *\zh{他把雅温得回了。})\\
\zh{我们明天去北京。} \emph{Wǒmen míngtiān qù Běijīng.} — « Nous allons à Pékin demain. »\\
\zh{老师来了。} \emph{Lǎoshī lái le.} — « Le professeur est arrivé. »
\end{exemplebox}

Attention à ne pas confondre : \zh{我把书拿来了。} \emph{Wǒ bǎ shū ná lái le.} (« j'ai apporté le livre ») est correct, car le verbe est \zh{拿} (prendre, transitif) et \zh{来} n'est qu'un complément de direction. C'est le verbe principal qui compte.

\courssub{Restriction 4 : la forme potentielle est incompatible avec 把}

\begin{propriete}[Pas de complément de possibilité avec 把]
Le complément de possibilité (verbe + \zh{得}/\zh{不} + résultat : \zh{做得完} « pouvoir finir », \zh{做不完} « ne pas pouvoir finir ») exprime une \emph{possibilité}. Or la phrase en \zh{把} décrit un \cle{résultat réel}, effectivement obtenu sur l'objet. Les deux sont incompatibles.\\[0.3em]
Incorrect : *\zh{我把作业做得完。} / *\zh{我把作业做不完。}\\
Correct (forme potentielle, ordre SVO) : \zh{我做不完作业。} \emph{Wǒ zuò bu wán zuòyè.} — « Je ne peux pas finir les devoirs. »\\
Correct (phrase en 把, résultat réel) : \zh{我把作业做完了。} \emph{Wǒ bǎ zuòyè zuò wán le.} — « J'ai fini les devoirs. »
\end{propriete}

Logique à retenir : \zh{把} présente l'action comme \emph{accomplie} sur l'objet ; la forme potentielle présente l'action comme \emph{possible ou impossible}. On ne peut pas affirmer en même temps qu'un résultat est réel et qu'il est seulement possible.

\courssub{Cas particulier : 把 + objet + verbe + 了 (verbes de disparition)}

\begin{definition}[Verbes de disparition]
Certains verbes signifient que l'objet \emph{disparaît}, est \emph{perdu} ou \emph{oublié} : \zh{丢} (perdre), \zh{忘} (oublier), \zh{死} (mourir), \zh{坏} (se gâter). Avec eux, la structure minimale \zh{把} + objet + verbe + \zh{了} suffit, sans autre complément : le \zh{了} final marque à lui seul le résultat (la disparition est accomplie).\\[0.3em]
\zh{我把钱包丢了。} \emph{Wǒ bǎ qiánbāo diū le.} — « J'ai perdu mon portefeuille. »\\
\zh{我把他的名字忘了。} \emph{Wǒ bǎ tā de míngzi wàng le.} — « J'ai oublié son nom. »\\
\zh{他把手机弄坏了。} \emph{Tā bǎ shǒujī nòng huài le.} — « Il a abîmé son téléphone. »
\end{definition}

Notez que le français dit simplement « j'ai perdu mon portefeuille » : le chinois, avec \zh{把}, insiste sur le sort subi par l'objet (souvent avec une nuance de regret ou de reproche).

\courssub{Récapitulatif : quels verbes avec 把 ?}

\begin{center}
\begin{tabularx}{\linewidth}{|X|X|X|X|}
\hline
\rowcolor{popblueL} Verbe & Pinyin & Traduction & Avec 把 ? \\
\hline
放 & fàng & poser & Oui : \zh{我把书放在桌子上。} \\
\hline
洗 & xǐ & laver & Oui : \zh{我把衣服洗了。} \\
\hline
写 & xiě & écrire & Oui : \zh{我把字写错了。} \\
\hline
关 & guān & fermer & Oui : \zh{请把窗户关上。} \\
\hline
送 & sòng & offrir, envoyer & Oui : \zh{我把礼物送给他了。} \\
\hline
翻译 & fānyì & traduire & Oui : \zh{我把课文翻译成法语了。} \\
\hline
喜欢 & xǐhuan & aimer & Non : \zh{我喜欢他。} \\
\hline
知道 & zhīdao & savoir & Non : \zh{我知道这件事。} \\
\hline
是 & shì & être & Non : \zh{他是我的老师。} \\
\hline
在 & zài & être (lieu) & Non : \zh{他在杜阿拉。} \\
\hline
有 & yǒu & avoir & Non : \zh{我有两个弟弟。} \\
\hline
来 / 去 / 走 & lái / qù / zǒu & venir / aller / partir & Non (intransitifs) \\
\hline
\end{tabularx}
\end{center}

\begin{popfigure}
\begin{tikzpicture}
\node[draw=popgreen, fill=popgreenL, rounded corners, align=center, font=\bfseries] (oui) at (0,0) {Verbes qui aiment \zh{把}\\ action + résultat sur l'objet};
\node[draw=poppink, fill=poppinkL, rounded corners, align=center, font=\bfseries] (non) at (9.5,0) {Verbes qui refusent \zh{把}\\ état, sentiment, possession};
\node[draw=popgreen!60, fill=white, rounded corners, align=center] at (0,-2.2) {\zh{放} poser\\ \zh{洗} laver\\ \zh{写} écrire};
\node[draw=popgreen!60, fill=white, rounded corners, align=center] at (0,-4.6) {\zh{关} fermer\\ \zh{送} envoyer\\ \zh{翻译} traduire};
\node[draw=poppink!60, fill=white, rounded corners, align=center] at (9.5,-2.2) {\zh{喜欢} aimer\\ \zh{知道} savoir\\ \zh{是} être};
\node[draw=poppink!60, fill=white, rounded corners, align=center] at (9.5,-4.6) {\zh{在} être (lieu)\\ \zh{有} avoir\\ \zh{来}/\zh{去} venir/aller};
\draw[-{Stealth}, popgreen, thick] (oui) -- (0,-1.3);
\draw[-{Stealth}, poppink, thick] (non) -- (9.5,-1.3);
\end{tikzpicture}
\legende{Carte mentale : les verbes compatibles et incompatibles avec la structure en \zh{把}.}
\end{popfigure}

\begin{methode}[Le test en une question avant d'utiliser 把]
Avant de construire une phrase en \zh{把}, posez-vous la question : \emph{« le verbe fait-il subir quelque chose de concret à un objet précis ? »}
\begin{enumerate}
\item \textbf{L'objet est-il précis et connu ?} Si non (un, une quelconque) $\rightarrow$ ordre SVO : \zh{我吃了一个苹果。}
\item \textbf{Le verbe agit-il physiquement sur l'objet ?} Si c'est un verbe de sentiment, d'état, de possession (\zh{喜欢}, \zh{是}, \zh{有}) $\rightarrow$ ordre SVO : \zh{我喜欢他。}
\item \textbf{Le verbe est-il transitif ?} Si c'est \zh{来}, \zh{去}, \zh{走} seul $\rightarrow$ pas de \zh{把}.
\item \textbf{Exprimez-vous une possibilité ?} Si vous voulez dire « pouvoir / ne pas pouvoir faire » $\rightarrow$ forme potentielle sans \zh{把} : \zh{我做不完作业。}
\item Si les quatre réponses permettent \zh{把}, construisez : Sujet + \zh{把} + objet défini + verbe + complément de résultat : \zh{我把作业做完了。}
\end{enumerate}
\end{methode}

\begin{exoresolu}[Corriger les phrases fautives]
Énoncé : les phrases suivantes, produites par des élèves francophones, sont incorrectes. Identifiez la restriction violée et corrigez.\\
1. *\zh{我把一个芒果吃了。} \quad 2. *\zh{我把汉语喜欢。} \quad 3. *\zh{我把学校去了。} \quad 4. *\zh{我把这篇文章翻译得完。}
\tcblower
Solution détaillée :
\begin{enumerate}
\item Restriction 1 (objet indéfini). Correction : \zh{我吃了一个芒果。} \emph{Wǒ chī le yī gè mángguǒ.} — « J'ai mangé une mangue. » Ou bien, si la mangue est précise : \zh{我把那个芒果吃了。} \emph{Wǒ bǎ nà ge mángguǒ chī le.} — « J'ai mangé cette mangue. »
\item Restriction 2 (verbe de sentiment). Correction : \zh{我喜欢汉语。} \emph{Wǒ xǐhuan Hànyǔ.} — « J'aime le chinois. »
\item Restriction 3 (verbe intransitif). Correction : \zh{我去学校了。} \emph{Wǒ qù xuéxiào le.} — « Je suis allé à l'école. »
\item Restriction 4 (forme potentielle). Correction : \zh{我翻译得完这篇文章。} \emph{Wǒ fānyì de wán zhè piān wénzhāng.} — « Je peux finir de traduire cet article. » Ou, résultat réel : \zh{我把这篇文章翻译完了。} \emph{Wǒ bǎ zhè piān wénzhāng fānyì wán le.} — « J'ai fini de traduire cet article. »
\end{enumerate}
\end{exoresolu}

\begin{exercice}[À vous de juger]
Dites si \zh{把} est possible ; si oui, construisez la phrase, si non, donnez la phrase SVO correcte : 1. finir le riz (\zh{吃完}) ; 2. aimer le ndolé ; 3. aller au marché ; 4. fermer la porte ; 5. ne pas pouvoir finir les devoirs.
\end{exercice}

\begin{corrige}[Exercice : À vous de juger]
1. \zh{我把饭吃完了。} \emph{Wǒ bǎ fàn chī wán le.} — « J'ai fini le riz. » (action + résultat, objet connu)\\
2. \zh{我喜欢恩多雷。} \emph{Wǒ xǐhuan Ēnduōléi.} — « J'aime le ndolé. » (sentiment : pas de \zh{把})\\
3. \zh{我去市场了。} \emph{Wǒ qù shìchǎng le.} — « Je suis allé au marché. » (intransitif)\\
4. \zh{请把门关上。} \emph{Qǐng bǎ mén guānshang.} — « Ferme la porte, s'il te plaît. »\\
5. \zh{我做不完作业。} \emph{Wǒ zuò bu wán zuòyè.} — « Je ne peux pas finir les devoirs. » (potentiel : pas de \zh{把})
\end{corrige}

\begin{aretenir}[Les quatre interdits de 把]
1. Objet indéfini interdit : *\zh{把一个苹果} $\rightarrow$ SVO.\\
2. Verbes d'état et de sentiment interdits : *\zh{把他喜欢} $\rightarrow$ \zh{我喜欢他。}\\
3. Verbes intransitifs interdits : *\zh{把学校去} $\rightarrow$ \zh{我去学校。}\\
4. Forme potentielle interdite : *\zh{把作业做得完} $\rightarrow$ \zh{我做得完作业。}\\
Cas particulier : verbes de disparition + \zh{了} : \zh{我把钱包丢了。}
\end{aretenir}

\coursec{Comparaison avec le français et pièges des francophones}

Nous avons vu dans la section précédente que la phrase en \zh{把} n'accepte ni les verbes d'état, ni un verbe nu sans complément, ni une négation mal placée. Ces restrictions deviennent beaucoup plus faciles à retenir quand on comprend \emph{pourquoi} le français et le chinois organisent différemment la même information. C'est l'objet de cette section : partir de ce que vous savez déjà — le français — pour construire des réflexes solides et éviter les quatre pièges classiques des apprenants francophones.

\courssub{Le français n'a pas d'équivalent de la phrase en 把}

\begin{definition}[Une structure sans équivalent français]
Le français ne possède \cle{aucune structure} correspondant à la phrase en \zh{把}. En français, l'objet reste après le verbe : \emph{Il a posé le livre sur la table}. En chinois, la phrase en \zh{把} \textbf{déplace l'objet avant le verbe} et place le résultat ou la destination après : \zh{他把书放在桌子上。} (\emph{Tā bǎ shū fàng zài zhuōzi shàng.} — Il a posé le livre sur la table).
\end{definition}

Observons d'abord trois phrases françaises très ordinaires et leur traduction chinoise :

\begin{exemplebox}[Le français dit « normalement », le chinois déplace l'objet]
\begin{itemize}
\item \zh{我把作业做完了。} — \emph{Wǒ bǎ zuòyè zuò wán le.} — « J'ai fini mes devoirs. »
\item \zh{他把书放在桌子上。} — \emph{Tā bǎ shū fàng zài zhuōzi shàng.} — « Il a posé le livre sur la table. »
\item \zh{请把门关上。} — \emph{Qǐng bǎ mén guān shàng.} — « Ferme la porte, s'il te plaît. »
\end{itemize}
\end{exemplebox}

Remarquez que le français ne signale \emph{rien} de spécial : « j'ai fini mes devoirs » a la même forme que « j'ai regardé mes devoirs ». Le chinois, lui, fait un choix : si l'objet est précis (\emph{mes} devoirs, \emph{le} livre, \emph{la} porte) et si le verbe produit un résultat (fini, posé, fermé), alors l'objet monte avant le verbe grâce à \zh{把}.

\begin{savaistu}[L'erreur inverse des Chinois qui apprennent le français]
Les Chinois qui apprennent le français font l'erreur \textbf{exactement inverse} de la vôtre : ils cherchent à « déplacer » l'objet en français et produisent des phrases comme *\emph{Il le livre a posé sur la table}. Leur professeur doit leur expliquer que le français garde l'objet après le verbe. Cette erreur miroir prouve une chose : les deux langues \textbf{organisent différemment la même information}. Le français suit l'ordre Sujet–Verbe–Objet ; le chinois, avec \zh{把}, suit l'ordre Sujet–Objet–Verbe–Résultat. Apprendre la phrase en \zh{把}, c'est donc apprendre à \emph{penser autrement}, pas seulement à ajouter un mot.
\end{savaistu}

\courssub{Comment le français rend-il une phrase en 把 ?}

Puisque le français n'a pas de structure équivalente, la traduction passe par plusieurs procédés. Il est essentiel de les connaître pour ne pas chercher un mot-à-mot impossible.

\begin{propriete}[Trois procédés de traduction du chinois vers le français]
\begin{enumerate}
\item \textbf{Phrase transitive ordinaire} (le plus fréquent) : \zh{我把作业做完了。} → « J'ai fini mes devoirs. » Le complément de résultat (\zh{完}) se fond dans le verbe français.
\item \textbf{Pronom ou reprise} : \zh{我把衣服洗干净了。} (\emph{Wǒ bǎ yīfu xǐ gānjìng le.}) → « Les voilà propres, les vêtements » ou « J'ai lavé les vêtements, ils sont propres maintenant. »
\item \textbf{Tournure causative ou périphrastique} : \zh{他把我的电脑用坏了。} → « Il a abîmé mon ordinateur en l'utilisant », parfois « il a fait en sorte que mon ordinateur soit hors d'usage ».
\end{enumerate}
\end{propriete}

\begin{center}
\begin{tabularx}{\linewidth}{|X|X|X|}
\hline
\rowcolor{popblueL} Phrase en 把 (caractères + pinyin) & Traduction naturelle & Procédé français \\
\hline
\zh{我把作业做完了。} \emph{Wǒ bǎ zuòyè zuò wán le.} & J'ai fini mes devoirs. & verbe simple \\
\hline
\zh{他把书放在桌子上。} \emph{Tā bǎ shū fàng zài zhuōzi shàng.} & Il a posé le livre sur la table. & verbe + complément de lieu \\
\hline
\zh{请把门关上。} \emph{Qǐng bǎ mén guān shàng.} & Ferme la porte, s'il te plaît. & impératif simple \\
\hline
\zh{我把衣服洗干净了。} \emph{Wǒ bǎ yīfu xǐ gānjìng le.} & Les voilà propres, les vêtements. & reprise / résultat exprimé à part \\
\hline
\zh{他把我的电脑用坏了。} \emph{Tā bǎ wǒ de diànnǎo yòng huài le.} & Il a abîmé mon ordinateur en l'utilisant. & périphrase causative \\
\hline
\zh{我没把名字写对。} \emph{Wǒ méi bǎ míngzi xiě duì.} & Je n'ai pas écrit le nom correctement. & négation + adverbe \\
\hline
\end{tabularx}
\end{center}

\courssub{Schéma de la transformation : du français au chinois}

La figure suivante montre le « voyage » de l'objet : en français, \emph{le livre} reste derrière le verbe ; en chinois, \zh{把} le fait passer devant, et la destination \zh{在桌子上} se place après le verbe.

\begin{popfigure}
\begin{tikzpicture}[
  mot/.style={rectangle, rounded corners=2pt, draw=popblue, fill=popblueL, align=center, font=\small, inner sep=3pt},
  motc/.style={rectangle, rounded corners=2pt, draw=poporange, fill=poporangeL, align=center, font=\small, inner sep=3pt},
  fl/.style={-{Stealth[length=2.5mm]}, thick, popdark}
]
% Phrase française
\node[mot] (f1) at (0,0) {Il};
\node[mot, right=0.25cm of f1] (f2) {a posé};
\node[mot, right=0.25cm of f2] (f3) {le livre};
\node[mot, right=0.25cm of f3] (f4) {sur la table};
\node[font=\small\bfseries, text=popblue, left=0.15cm of f1] {FR};
% Phrase chinoise
\node[motc] (c1) at (0,-2.2) {他};
\node[motc, right=0.25cm of c1] (c2) {把书};
\node[motc, right=0.25cm of c2] (c3) {放在};
\node[motc, right=0.25cm of c3] (c4) {桌子上};
\node[font=\small\bfseries, text=poporange, left=0.15cm of c1] {ZH};
% Flèches de correspondance
\draw[fl] (f1.south) -- (c1.north);
\draw[fl, poporange] (f3.south) to[bend right=15] node[midway, below, font=\scriptsize, align=center]{l'objet monte\\avant le verbe} (c2.north);
\draw[fl] (f2.south) -- (c3.north);
\draw[fl] (f4.south) -- (c4.north);
\end{tikzpicture}
\legende{De « Il a posé le livre sur la table » à \zh{他把书放在桌子上。} : l'objet \zh{书} est déplacé avant le verbe grâce à \zh{把}.}
\end{popfigure}

\courssub{Les quatre pièges des francophones}

\textbf{Piège 1 : garder l'ordre français Sujet–Verbe–Objet.}\par
Le réflexe français pousse à dire le verbe d'abord, puis à caser \zh{把} n'importe où. Résultat : *\zh{我放把书在桌子上。} Cette phrase est impossible : \zh{把} doit \textbf{précéder immédiatement l'objet}, et le bloc \zh{把} + objet doit \textbf{précéder le verbe}. Correct : \zh{我把书放在桌子上。}

\textbf{Piège 2 : oublier le complément après le verbe.}\par
En français, « je ferme la porte » est complet. En chinois, \zh{把} + objet + verbe nu est \emph{incomplet} : le verbe doit être suivi d'un résultat, d'une destination, de \zh{了} ou d'un autre complément.

\begin{attention}[*我把门关 : la phrase inachevée]
En français, « ferme la porte » suffit. En chinois, *\zh{我把门关。} est une phrase \textbf{suspendue, inachevée} : l'auditeur attend la suite (« tu la fermes… et alors ? »). Il faut un complément :
\begin{itemize}
\item \zh{我把门关上。} — \emph{Wǒ bǎ mén guān shàng.} — « Je ferme la porte » (résultat : fermée).
\item \zh{我把门关了。} — \emph{Wǒ bǎ mén guān le.} — « J'ai fermé la porte » (accompli).
\end{itemize}
Même punition pour : *\zh{他把饭吃。} → \zh{他把饭吃完了。} (\emph{Tā bǎ fàn chī wán le.} — « Il a fini son riz »).
\end{attention}

\textbf{Piège 3 : placer la négation après 把, comme en français.}\par
En français, la négation encadre le verbe : « je n'ai \emph{pas} fini mes devoirs ». Le réflexe donne *\zh{我把作业没做完。} Faux ! En chinois, la négation \zh{没（有）} ou \zh{不} se place \textbf{avant} \zh{把}.

\begin{propriete}[Place de la négation]
Structure : \cle{Sujet + 没/不 + 把 + objet + verbe + complément}.\par
\zh{我没把作业做完。} — \emph{Wǒ méi bǎ zuòyè zuò wán.} — « Je n'ai pas fini mes devoirs. »\par
\zh{他没把名字写对。} — \emph{Tā méi bǎ míngzi xiě duì.} — « Il n'a pas écrit le nom correctement. »\par
Même règle pour les adverbes modaux et auxiliaires (\zh{应该}, \zh{要}, \zh{想}) : \zh{你应该把窗户关上。} — \emph{Nǐ yīnggāi bǎ chuānghu guān shàng.} — « Tu devrais fermer la fenêtre. »
\end{propriete}

\textbf{Piège 4 : utiliser 把 pour tout objet défini, même avec des verbes d'état.}\par
Le français dit « j'aime \emph{ce} film », « je connais \emph{cette} personne » : l'objet est bien défini, mais le verbe ne \emph{transforme} rien. Avec les verbes d'état ou de sentiment (\zh{爱}, \zh{喜欢}, \zh{知道}, \zh{认识}, \zh{想}, \zh{是}, \zh{有}), la phrase en \zh{把} est impossible : *\zh{我把这个电影喜欢。} On garde l'ordre ordinaire : \zh{我喜欢这个电影。} (\emph{Wǒ xǐhuan zhège diànyǐng.} — J'aime ce film). \zh{把} exige un verbe d'\textbf{action} qui produit un \textbf{effet} sur l'objet.

\begin{center}
\begin{tabularx}{\linewidth}{|X|X|X|}
\hline
\rowcolor{popblueL} Piège & Phrase fautive & Correction (caractères + pinyin + traduction) \\
\hline
1. Ordre français & *\zh{我放把书在桌子上。} & \zh{我把书放在桌子上。} \emph{Wǒ bǎ shū fàng zài zhuōzi shàng.} — J'ai posé le livre sur la table. \\
\hline
2. Verbe nu & *\zh{我把门关。} & \zh{我把门关上。} \emph{Wǒ bǎ mén guān shàng.} — Je ferme la porte. \\
\hline
3. Négation mal placée & *\zh{我把作业没做完。} & \zh{我没把作业做完。} \emph{Wǒ méi bǎ zuòyè zuò wán.} — Je n'ai pas fini mes devoirs. \\
\hline
4. Verbe d'état & *\zh{我把这个电影喜欢。} & \zh{我喜欢这个电影。} \emph{Wǒ xǐhuan zhège diànyǐng.} — J'aime ce film. \\
\hline
\end{tabularx}
\end{center}

\courssub{Stratégie de traduction : du français vers le chinois}

\begin{methode}[Traduire une phrase française avec 把 en trois questions]
Face à une phrase française à traduire en chinois, posez-vous trois questions :
\begin{enumerate}
\item \textbf{L'objet est-il précis ?} Cherchez \emph{le, la, les, ce, mon, ton…} : « fini \emph{mes} devoirs », « posé \emph{le} livre ». Si l'objet est indéfini (« j'ai mangé \emph{un} gâteau »), pas de \zh{把}.
\item \textbf{Y a-t-il un résultat ou une destination ?} Cherchez ce qui arrive à l'objet : « fini » (résultat : terminé), « posé \emph{sur la table} » (destination), « fermée », « abîmé ». Sans résultat, pas de \zh{把}.
\item \textbf{Oui aux deux ?} Alors construisez : \cle{Sujet + 把 + objet + verbe + résultat/destination (+ 了)}. Sinon, gardez l'ordre ordinaire Sujet + verbe + objet.
\end{enumerate}
Exemple guidé : « Il a abîmé mon ordinateur. » → objet précis ? oui (\emph{mon} ordinateur) ; résultat ? oui (abîmé = \zh{坏}). → \zh{他把我的电脑用坏了。}
\end{methode}

\begin{exoresolu}[Transformation et correction]
\textbf{Énoncé.} a) Traduisez en chinois avec \zh{把} : « Elle a mis les cahiers dans le sac. » b) Corrigez la phrase fautive : *\zh{我把窗户没关好。}
\tcblower
\textbf{Solution détaillée.}\par
a) Application de la méthode en trois questions. (1) Objet précis ? Oui : \emph{les} cahiers → \zh{本子}. (2) Résultat ou destination ? Oui : \emph{dans le sac} → destination \zh{在书包里}. (3) Construction : Sujet \zh{她} + \zh{把} + objet \zh{本子} + verbe \zh{放} + destination \zh{在书包里}.\par
Réponse : \zh{她把本子放在书包里。} — \emph{Tā bǎ běnzi fàng zài shūbāo lǐ.} — « Elle a mis les cahiers dans le sac. »\par
b) La phrase *\zh{我把窗户没关好。} tombe dans le piège 3 : la négation \zh{没} est placée après \zh{把}, sur le modèle français « je n'ai pas bien fermé la fenêtre ». En chinois, la négation précède \zh{把}.\par
Correction : \zh{我没把窗户关好。} — \emph{Wǒ méi bǎ chuānghu guān hǎo.} — « Je n'ai pas bien fermé la fenêtre. »
\end{exoresolu}

\begin{exemplebox}[Deux modèles à retenir par cœur]
\begin{itemize}
\item \zh{我没把名字写对。} — \emph{Wǒ méi bǎ míngzi xiě duì.} — « Je n'ai pas écrit le nom correctement. » (négation avant \zh{把} ; résultat \zh{对}).
\item \zh{他把我的电脑用坏了。} — \emph{Tā bǎ wǒ de diànnǎo yòng huài le.} — « Il a abîmé mon ordinateur en l'utilisant. » (résultat \zh{坏} ; le français passe par une périphrase).
\end{itemize}
\end{exemplebox}

\begin{aretenir}[L'essentiel de la section]
\begin{itemize}
\item Le français n'a \textbf{aucun équivalent} de la phrase en \zh{把} : on traduit par une phrase ordinaire, une reprise ou une périphrase.
\item Le chinois \textbf{déplace l'objet avant le verbe} ; le français le garde après. Ne calquez jamais l'ordre français.
\item Quatre pièges : ordre S–V–O français, verbe nu sans complément, négation après \zh{把}, \zh{把} avec un verbe d'état.
\item Stratégie : objet précis + résultat/destination = candidat pour \zh{把}.
\end{itemize}
\end{aretenir}

\begin{exercice}[À vous de jouer]
Traduisez en chinois avec \zh{把} quand c'est possible, sinon justifiez l'ordre ordinaire : 1) « J'ai rangé les livres dans l'armoire. » 2) « J'aime beaucoup ce plat camerounais. » 3) « Il n'a pas fermé la porte. » 4) « Elle a fini son thé. »
\end{exercice}

\begin{corrige}[Exercice « À vous de jouer »]
1) Objet précis (\emph{les} livres) + destination (\emph{dans l'armoire}) → \zh{我把书放在柜子里。} (\emph{Wǒ bǎ shū fàng zài guìzi lǐ.})\par
2) Verbe de sentiment (\emph{aimer}) → pas de \zh{把} : \zh{我很喜欢这个喀麦隆菜。} (\emph{Wǒ hěn xǐhuan zhège Kāmàilóng cài.})\par
3) Négation avant \zh{把} : \zh{他没把门关上。} (\emph{Tā méi bǎ mén guān shàng.})\par
4) Objet précis (\emph{son} thé) + résultat (\emph{fini}) : \zh{她把茶喝完了。} (\emph{Tā bǎ chá hē wán le.})
\end{corrige}

\coursec{Entraînement : application et transformation}

Après avoir comparé la phrase en \zh{把} avec le français et repéré les pièges qui guettent les francophones (ordre des mots calqué sur le français, objet indéfini, verbe nu sans complément, négation mal placée), il est temps de passer à la pratique. Cette dernière étape de la leçon vous entraîne progressivement : vous allez d'abord \emph{reconnaître} et \emph{transformer}, puis \emph{corriger}, \emph{traduire} et enfin \emph{produire librement}. C'est exactement le parcours exigé aux évaluations de classe et au baccalauréat.

\begin{popfigure}
\begin{tikzpicture}[font=\small]
\foreach \i/\txt/\col in {0/{Reconnaître}/popblue, 1/{Remettre en\\ordre}/popgreen, 2/{Transformer}/poporange, 3/{Corriger}/poppink, 4/{Traduire}/poppurple, 5/{Produire\\librement}/popgold}{
  \node[draw=\col, fill=\col!25, rounded corners, align=center, minimum width=2.1cm, minimum height=1cm] (n\i) at (\i*2.35,0) {\txt};
}
\foreach \i in {0,1,2,3,4}{
  \pgfmathtruncatemacro{\j}{\i+1}
  \draw[-{Stealth}, thick, popdark] (n\i) -- (n\j);
}
\node[popmuted, align=center] at (5.875,-1.1) {Progression de l'entraînement : du guidé vers l'autonome};
\end{tikzpicture}
\legende{Frise de progression des exercices : chaque étape prépare la suivante.}
\end{popfigure}

\courssub{La méthode de transformation : quatre étapes incontournables}

\begin{methode}[Transformer une phrase SVO en phrase en \zh{把}]
\begin{enumerate}
\item \textbf{Repérer l'objet} du verbe : il doit être \cle{défini} (déterminé par 这/那, un possessif, ou connu du contexte). S'il est indéfini (\zh{一个苹果}), la transformation est impossible.
\item \textbf{Placer l'objet après \zh{把}} : Sujet + 把 + objet + \ldots
\item \textbf{Garder le verbe et son complément ensemble} : le verbe ne reste jamais seul ; il est suivi de 了, d'un complément de résultat (完, 干净, 好), d'une localisation (放在\ldots 上) ou de manière. Verbe + complément forment un \cle{bloc inséparable}.
\item \textbf{Vérifier la place de la négation et des auxiliaires} : 没(有), 别, 想, 要, 会 se placent \cle{avant} \zh{把}, jamais entre 把 et le verbe.
\end{enumerate}
\end{methode}

\begin{exemplebox}[Corrigé type : les trois modèles à imiter]
\zh{我把钥匙放在桌子上了。} \emph{(Wǒ bǎ yàoshi fàng zài zhuōzi shang le.)} « J'ai posé les clés sur la table. »\\
\zh{妈妈把盘子洗干净了。} \emph{(Māma bǎ pánzi xǐ gānjìng le.)} « Maman a lavé les assiettes (elles sont propres). »\\
\zh{别把窗户关上！} \emph{(Bié bǎ chuānghu guān shang !)} « Ne ferme pas la fenêtre ! »
\end{exemplebox}

\courssub{Exercice 1 — Transformation SVO vers \zh{把字句}}

\begin{exoresolu}[Transformer en phrase en \zh{把}]
Énoncé. Transformez : \zh{他喝完了那杯咖啡。} \emph{(Tā hē wán le nà bēi kāfēi.)} « Il a fini cette tasse de café. »
\tcblower
Solution détaillée. Étape 1 : l'objet est \zh{那杯咖啡} « cette tasse de café » — il est défini grâce à \zh{那}, donc la transformation est possible. Étape 2 : on place l'objet après \zh{把} : 他 + 把 + 那杯咖啡. Étape 3 : le bloc verbe + complément \zh{喝完了} reste soudé et se place à la fin. Résultat : \zh{他把那杯咖啡喝完了。} \emph{(Tā bǎ nà bēi kāfēi hē wán le.)} « Il a fini cette tasse de café. » La phrase en \zh{把} insiste sur le \emph{sort réservé à l'objet} : le café n'existe plus, la tasse est vide.
\end{exoresolu}

\begin{exercice}[À vous de transformer]
Transformez les phrases suivantes en phrases en \zh{把} :
\begin{enumerate}
\item \zh{弟弟吃完了那碗饭。} \emph{(Dìdi chī wán le nà wǎn fàn.)}
\item \zh{我写完了今天的作业。} \emph{(Wǒ xiě wán le jīntiān de zuòyè.)}
\item \zh{她洗干净了那些衣服。} \emph{(Tā xǐ gānjìng le nàxiē yīfu.)}
\end{enumerate}
\end{exercice}

\begin{corrige}[Exercice « À vous de transformer »]
1. \zh{弟弟把那碗饭吃完了。} \emph{(Dìdi bǎ nà wǎn fàn chī wán le.)} « Mon petit frère a fini ce bol de riz. » — 2. \zh{我把今天的作业写完了。} \emph{(Wǒ bǎ jīntiān de zuòyè xiě wán le.)} « J'ai fini les devoirs d'aujourd'hui. » — 3. \zh{她把那些衣服洗干净了。} \emph{(Tā bǎ nàxiē yīfu xǐ gānjìng le.)} « Elle a lavé ces vêtements (ils sont propres). » Dans les trois cas, l'objet défini (那碗饭, 今天的作业, 那些衣服) passe après \zh{把}, et le bloc verbe + complément de résultat (吃完, 写完, 洗干净) reste soudé en fin de phrase.
\end{corrige}

\courssub{Exercice 2 — Remise en ordre des mots}

\begin{exoresolu}[Remettre en ordre]
Énoncé. Remettez en ordre : 把 / 我 / 作业 / 完 / 做 / 了
\tcblower
Solution détaillée. Le sujet est \zh{我}. La préposition \zh{把} introduit l'objet \zh{作业} : 我把作业. Restent 做, 完, 了 : le verbe \zh{做} et son complément de résultat \zh{完} forment le bloc inséparable \zh{做完}, suivi de \zh{了} qui marque l'accomplissement. Résultat : \zh{我把作业做完了。} \emph{(Wǒ bǎ zuòyè zuò wán le.)} « J'ai fini mes devoirs. »
\end{exoresolu}

\begin{attention}[Ne séparez jamais le verbe de son complément de résultat]
Erreur fréquente des francophones dans la remise en ordre : traiter chaque mot comme une bille isolée et écrire *\zh{我把作业完做了} ou *\zh{我把完作业做了}. En chinois, \zh{做完} « finir de faire », \zh{放好} « bien ranger », \zh{洗干净} « laver propre » sont des \cle{blocs soudés} : le complément de résultat se colle immédiatement après le verbe, et rien ne peut s'intercaler entre eux. Mémorisez-les comme des unités : 做完, 写好, 关上, 放在.
\end{attention}

\begin{exercice}[Remise en ordre]
Remettez les mots dans l'ordre pour former une phrase en \zh{把} correcte :
\begin{enumerate}
\item 把 / 他 / 书 / 在 / 放 / 桌子 / 上 / 了
\item 干净 / 把 / 妈妈 / 洗 / 了 / 盘子
\item 把 / 请 / 门 / 关上 / 你
\end{enumerate}
\end{exercice}

\begin{corrige}[Exercice « Remise en ordre »]
1. \zh{他把书放在桌子上了。} \emph{(Tā bǎ shū fàng zài zhuōzi shang le.)} « Il a posé le livre sur la table. » — 2. \zh{妈妈把盘子洗干净了。} \emph{(Māma bǎ pánzi xǐ gānjìng le.)} « Maman a lavé les assiettes. » — 3. \zh{请你把门关上。} \emph{(Qǐng nǐ bǎ mén guān shang.)} « Ferme la porte, s'il te plaît. » Remarquez que \zh{请} « s'il te plaît » se place en tête de phrase, avant le sujet.
\end{corrige}

\courssub{Exercice 3 — Correction de phrases fautives}

\begin{exoresolu}[Corriger les phrases fautives]
Énoncé. Corrigez : (a) *\zh{我把一个苹果吃了。} (b) *\zh{我把作业没做完。}
\tcblower
Solution détaillée. (a) L'objet \zh{一个苹果} « une pomme » est \cle{indéfini} : or la construction en \zh{把} exige un objet déterminé, connu des deux interlocuteurs. Deux corrections possibles : revenir à l'ordre SVO \zh{我吃了一个苹果。} \emph{(Wǒ chī le yí ge píngguǒ.)} « J'ai mangé une pomme », ou définir l'objet : \zh{我把那个苹果吃了。} \emph{(Wǒ bǎ nà ge píngguǒ chī le.)} « J'ai mangé cette pomme-là. » (b) La négation \zh{没} est mal placée : elle doit précéder \zh{把}, jamais se loger entre l'objet et le verbe. Correction : \zh{我没把作业做完。} \emph{(Wǒ méi bǎ zuòyè zuò wán.)} « Je n'ai pas fini mes devoirs. »
\end{exoresolu}

\begin{center}
\begin{tabularx}{\linewidth}{|X|X|X|}
\hline
\rowcolor{popblueL} Phrase fautive & Correction & Règle violée \\
\hline
*\zh{我把一个苹果吃了。} & \zh{我吃了一个苹果。} & Objet indéfini interdit après 把 \\
\hline
*\zh{我把作业没做完。} & \zh{我没把作业做完。} & La négation se place avant 把 \\
\hline
*\zh{他把咖啡喝。} & \zh{他把咖啡喝完了。} & Verbe nu interdit : complément obligatoire \\
\hline
*\zh{我把门想关上。} & \zh{我想把门关上。} & L'auxiliaire se place avant 把 \\
\hline
\end{tabularx}
\end{center}

\courssub{Exercice 4 — Choisir entre l'ordre SVO et la construction en \zh{把}}

\begin{methode}[SVO ou \zh{把} ? Le test en deux questions]
\begin{enumerate}
\item L'objet est-il \cle{défini} (déterminé, connu du contexte) ? Si non, imposez l'ordre SVO.
\item Le verbe porte-t-il un \cle{résultat}, une \cle{disposition} ou une \cle{localisation} qui affecte l'objet (喝完, 洗干净, 放在\ldots 上) ? Si oui, la construction en \zh{把} est préférable, car elle met en valeur le devenir de l'objet.
\end{enumerate}
Si les deux réponses sont « oui », choisissez \zh{把} ; sinon, restez en SVO.
\end{methode}

\begin{exercice}[Choix de la construction]
Pour chaque phrase, dites si l'ordre SVO ou la construction en \zh{把} convient, puis justifiez :
\begin{enumerate}
\item « J'ai acheté un livre. » (un livre quelconque)
\item « J'ai rangé ce livre sur l'étagère. » (ce livre précis)
\item « Il a fini le ndolé que maman a préparé. »
\item « Elle a vu un film hier. »
\end{enumerate}
\end{exercice}

\begin{corrige}[Exercice « Choix de la construction »]
1. SVO : \zh{我买了一本书。} \emph{(Wǒ mǎi le yì běn shū.)} — objet indéfini, 把 impossible. 2. \zh{把} : \zh{我把这本书放在书架上了。} \emph{(Wǒ bǎ zhè běn shū fàng zài shūjià shang le.)} — objet défini + localisation résultative. 3. \zh{把} : \zh{他把妈妈做的恩多莱吃完了。} \emph{(Tā bǎ māma zuò de ndolé chī wán le.)} « Il a fini le ndolé que maman a préparé » — objet défini (le plat précis préparé par maman) + résultat 完. 4. SVO : \zh{她昨天看了一个电影。} \emph{(Tā zuótiān kàn le yí ge diànyǐng.)} — objet indéfini, simple constat sans résultat sur l'objet.
\end{corrige}

\courssub{Exercice 5 — Thème : traduction du français vers le chinois}

\begin{exoresolu}[Traduire en chinois]
Énoncé. Traduisez : « J'ai mis les clés sur la table. Maman a lavé les assiettes. Ne ferme pas la fenêtre. »
\tcblower
Solution détaillée. Phrase 1 : « les clés » est défini, « mettre sur la table » exprime une disposition de l'objet → construction en \zh{把} avec 放在\ldots 上 : \zh{我把钥匙放在桌子上了。} \emph{(Wǒ bǎ yàoshi fàng zài zhuōzi shang le.)} Phrase 2 : « les assiettes » est défini, « laver » aboutit à un résultat (propre) → \zh{把} avec 洗干净 : \zh{妈妈把盘子洗干净了。} \emph{(Māma bǎ pánzi xǐ gānjìng le.)} Phrase 3 : impératif négatif → \zh{别} placé \emph{avant} \zh{把} : \zh{别把窗户关上！} \emph{(Bié bǎ chuānghu guān shang !)}
\end{exoresolu}

\begin{exercice}[Thème d'application]
Traduisez en chinois avec la construction en \zh{把} :
\begin{enumerate}
\item « J'ai fini mes devoirs. »
\item « Il a rangé les chaises dans la salle de classe. »
\item « N'oublie pas ton sac ! » (utiliser 别 + 把 + 忘了)
\end{enumerate}
\end{exercice}

\begin{corrige}[Exercice « Thème d'application »]
1. \zh{我把作业做完了。} \emph{(Wǒ bǎ zuòyè zuò wán le.)} — 2. \zh{他把椅子放在教室里了。} \emph{(Tā bǎ yǐzi fàng zài jiàoshì lǐ le.)} — 3. \zh{别把你的包忘了！} \emph{(Bié bǎ nǐ de bāo wàng le !)} Remarquez : 别 précède 把, et le verbe 忘 est accompagné de 了.
\end{corrige}

\courssub{Exercice 6 — Production guidée : le rangement de ma chambre}

\begin{textebox}[Modèle de production : le week-end, je range]
周末，我把房间打扫干净了，把衣服洗好了，还把作业做完了。妈妈很高兴。\\
\emph{(Zhōumò, wǒ bǎ fángjiān dǎsǎo gānjìng le, bǎ yīfu xǐ hǎo le, hái bǎ zuòyè zuò wán le. Māma hěn gāoxìng.)}\\
« Le week-end, j'ai bien nettoyé ma chambre, j'ai lavé mes vêtements, et j'ai aussi fini mes devoirs. Maman est très contente. »
\end{textebox}

Observez le modèle : trois phrases en \zh{把} se succèdent, chacune avec un complément différent (打扫干净, 洗好, 做完), et l'adverbe \zh{还} « en plus » placé avant \zh{把} pour enchaîner les actions. La dernière phrase, sans \zh{把}, donne la conséquence sur une personne : c'est un excellent plan de rédaction.

\begin{experience}[Production guidée : je range ma chambre]
Rédigez cinq phrases décrivant le rangement de votre chambre, avec \textbf{au moins trois phrases en \zh{把}}. Utilisez la banque de blocs verbe + complément : 打扫干净 (nettoyer), 放好 (bien ranger), 洗干净 (laver), 做完 (finir), 关上 (fermer), 放在\ldots 上 (poser sur). Variez les objets : 书, 衣服, 桌子, 窗户, 作业. Terminez par une phrase de bilan sans \zh{把} (par exemple : \zh{我的房间现在很干净。} « Ma chambre est maintenant très propre. »). Lisez ensuite votre texte à voix haute en soignant les tons.
\end{experience}

\begin{exemplebox}[Production possible (corrigé type)]
\zh{星期六早上，我把书放在书架上了。然后，我把衣服洗干净了。我把窗户关上了，因为外面有风。我还把作业做完了。我的房间现在很干净，我很高兴。}\\
\emph{(Xīngqīliù zǎoshang, wǒ bǎ shū fàng zài shūjià shang le. Ránhòu, wǒ bǎ yīfu xǐ gānjìng le. Wǒ bǎ chuānghu guān shang le, yīnwèi wàimiàn yǒu fēng. Wǒ hái bǎ zuòyè zuò wán le. Wǒ de fángjiān xiànzài hěn gānjìng, wǒ hěn gāoxìng.)}\\
« Samedi matin, j'ai posé les livres sur l'étagère. Ensuite, j'ai lavé mes vêtements. J'ai fermé la fenêtre parce qu'il y avait du vent dehors. J'ai aussi fini mes devoirs. Ma chambre est maintenant très propre, je suis content. »
\end{exemplebox}

\courssub{Barème indicatif type baccalauréat}

\begin{center}
\begin{tabularx}{\linewidth}{|X|X|}
\hline
\rowcolor{popblueL} Critère évalué & Points indicatifs \\
\hline
Exactitude de la structure : Sujet + 把 + objet \cle{déterminé} + verbe \cle{accompagné} (了, complément de résultat ou de localisation) & 4 points \\
\hline
Place correcte de la négation (没, 别) et des auxiliaires (想, 要, 会) \cle{avant} \zh{把} & 2 points \\
\hline
Choix pertinent du complément (完, 干净, 好, 在\ldots 上) et bloc verbe + complément soudé & 2 points \\
\hline
Choix justifié entre SVO et \zh{把} selon le contexte (objet défini + résultat) & 1 point \\
\hline
Exactitude des caractères et du pinyin (tons) & 1 point \\
\hline
\end{tabularx}
\end{center}

\begin{aretenir}[L'essentiel de l'entraînement]
Pour réussir tout exercice sur la phrase en \zh{把}, appliquez le réflexe en quatre points : (1) l'objet est-il défini ? (2) il passe après \zh{把} ; (3) le verbe n'est jamais nu — il garde son complément soudé à lui ; (4) négation et auxiliaires se placent avant \zh{把}. Avec ces quatre vérifications, la transformation, la correction, la traduction et la production libre deviennent des exercices sûrs, au contrôle continu comme au baccalauréat.
\end{aretenir}

\newpage

\coursec{Activité d'intégration}

\coursec{Leçon 28 — Grammaire : La structure \zh{把字句} (bǎ zìjù)}

\courssub{Objectifs d'apprentissage}

À l'issue de cette leçon, l'élève sera capable de :
\begin{itemize}
\item identifier la structure \zh{把字句} dans un énoncé oral ou écrit ;
\item expliquer le schéma syntaxique : \textbf{Sujet + 把 + Objet déterminé + Verbe + Complément obligatoire} ;
\item construire des phrases correctes pour exprimer une action de manipulation, de déplacement ou de transformation aboutie ;
\item placer correctement la négation (\zh{没、别、不要}) et les auxiliaires (\zh{要、应该、能}) \textbf{avant} \zh{把} ;
\item éviter les erreurs classiques des francophones (verbe nu, objet indéfini, ordre des mots) ;
\item utiliser la structure dans un contexte scolaire camerounais authentique (consignes de rangement, rapport d'activité).
\end{itemize}

\courssub{Situation déclenchante : Le grand ménage au lycée bilingue de Yaoundé}

\begin{textebox}[Annonce de M. Zhang au tableau — Vendredi après-midi]
同学们，星期一有检查。\\
Tóngxuémen, xīngqīyī yǒu jiǎnchá.\\
« Élèves, lundi il y a une inspection. »\\
\\
今天我们要打扫教室。请大家用“把”字句！\\
Jīntiān wǒmen yào dǎsǎo jiàoshì. Qǐng dàjiā yòng ``bǎ'' zìjù !\\
« Aujourd'hui nous devons nettoyer la salle. Utilisez la structure en 把, s'il vous plaît ! »\\
\\
例子：请把桌子排好。\\
Lìzi : Qǐng bǎ zhuōzi pái hǎo.\\
« Exemple : Rangez bien les tables (alignez-les). »
\end{textebox}

\courssub{Étape 1 : Observation guidée (10 min)}

Lisez l'annonce ci-dessus et répondez en groupe :
\begin{enumerate}
\item Soulignez la phrase contenant \zh{把}.
\item Identifiez ses trois composants obligatoires :
      \begin{itemize}
      \item L'objet concerné : \zh{桌子} (les tables — objet connu, déterminé).
      \item Le verbe d'action : \zh{排} (aligner, ranger).
      \item Le complément de résultat : \zh{好} (bien, correctement, jusqu'à achèvement).
      \end{itemize}
\item Pourquoi \zh{请排桌子} (Qǐng pái zhuōzi) est-il insuffisant ici ? \textit{Indice : que manque-t-il au verbe ?}
\item En français, on dit « Rangez les tables ». Quel renversement d'ordre opère le chinois avec \zh{把} ?
\end{enumerate}

\begin{aretenir}[Réponse attendue]
La structure \zh{把} place l'objet \textbf{avant} le verbe et \textbf{exige} un complément après le verbe (résultat, direction, \zh{了}, etc.) pour marquer que l'action est menée à son terme sur cet objet précis. \zh{请排桌子} est une phrase SVO simple ; elle ne met pas l'accent sur la manipulation aboutie de \emph{ces} tables-là.
\end{aretenir}

\courssub{Étape 2 : Schéma de la structure \zh{把字句}}

\begin{propriete}[Structure canonique et trois règles d'or]
\textbf{Schéma syntaxique :}
\begin{center}
\begin{tikzpicture}[node distance=0.8cm and 1.2cm, every node/.style={font=\small, align=center}]
\node (S) [draw, rounded corners, fill=popblueL] {Sujet};
\node (Ba) [right=of S, draw, rounded corners, fill=poporangeL] {\zh{把}};
\node (O) [right=of Ba, draw, rounded corners, fill=popgreenL] {Objet \textbf{déterminé}};
\node (V) [right=of O, draw, rounded corners, fill=poppinkL] {Verbe d'action};
\node (C) [right=of V, draw, rounded corners, fill=popgoldL, align=center]{Complément \textbf{obligatoire}\\ (\zh{了} / Résultat / Direction / Localisation)};
\draw[-Stealth, thick, popdark] (S) -- (Ba);
\draw[-Stealth, thick, popdark] (Ba) -- (O);
\draw[-Stealth, thick, popdark] (O) -- (V);
\draw[-Stealth, thick, popdark] (V) -- (C);
\end{tikzpicture}
\end{center}

\textbf{Règle 1 — Objet déterminé.} L'objet après \zh{把} est connu, spécifique, identifiable dans le contexte (défini par un démonstratif, un possessif, le contexte visuel, ou un quantificateur comme \zh{都、全}). \textit{Exclu :} \zh{一本书、个人} (indéfinis).

\textbf{Règle 2 — Verbe jamais nu.} Le verbe doit être suivi d'un élément qui « referme » l'action : \zh{了} (accompli), complément de résultat (\zh{好、干净、完、掉、错、坏}), complément de direction (\zh{上、下、进、出、回、过、起、开}), ou localisation (\zh{在/到/放进…里}).

\textbf{Règle 3 — Négation et auxiliaires AVANT \zh{把}.}
\begin{center}
\begin{tabularx}{\linewidth}{|X|X|X|}
\hline
\rowcolor{popblueL} Type & Position & Exemple \\
\hline
Négation \zh{没(有)} & \textbf{Sujet + 没 + 把} + O + V + Compl. & \zh{我没把门关上。} (Je n'ai pas fermé la porte.) \\
\hline
Négation impérative \zh{别 / 不要} & \textbf{别 / 不要 + 把} + O + V + Compl. & \zh{别把水倒掉。} (Ne verse pas l'eau / Ne la jette pas.) \\
\hline
Auxiliaire \zh{要、应该、能、想} & \textbf{Sujet + Aux + 把} + O + V + Compl. & \zh{你应该把作业交上来。} (Tu dois rendre le devoir.) \\
\hline
\end{tabularx}
\end{center}
\end{propriete}

\courssub{Étape 3 : Emplois et cas d'usage (avec exemples contextualisés)}

La \zh{把字句} sert à décrire \textbf{ce qu'on fait subir à un objet précis} (manipulation, déplacement, transformation, élimination). Le sujet est l'agent volontaire.

\begin{exemplebox}[Quatre grands types d'emploi]
\textbf{1. Displacement / Placement (移动/安置) — Complément de direction / localisation}\\
\zh{请把书放进书包里。} \emph{Qǐng bǎ shū fàng jìn shūbāo lǐ.} « Mets le livre dans le cartable. »\\
\zh{我们把椅子搬上去了。} \emph{Wǒmen bǎ yǐzi bān shàngqù le.} « Nous avons monté les chaises. »\\[4pt]
\textbf{2. Transformation / Traitement (处理/改变) — Complément de résultat}\\
\zh{请把黑板擦干净。} \emph{Qǐng bǎ hēibǎn cā gānjìng.} « Essuie le tableau (jusqu'à ce qu'il soit propre). »\\
\zh{他把衣服洗好了。} \emph{Tā bǎ yīfu xǐ hǎo le.} « Il a fini de laver les vêtements. »\\[4pt]
\textbf{3. Élimination / Consommation totale (消耗/处理掉) — \zh{掉、完、光}}\\
\zh{请把垃圾倒掉。} \emph{Qǐng bǎ lājī dào diào.} « Jette les ordures. »\\
\zh{弟弟把蛋糕吃完了。} \emph{Dìdi bǎ dàngāo chī wán le.} « Le petit frère a mangé tout le gâteau. »\\[4pt]
\textbf{4. Action accomplie avec \zh{了} (focus sur la fin de l'action sur l'objet)}\\
\zh{我把门关上了。} \emph{Wǒ bǎ mén guān shang le.} « J'ai fermé la porte (c'est fait). »
\end{exemplebox}

\courssub{Étape 4 : Exceptions et restrictions (verbes incompatibles)}

\begin{attention}[Verbes qui NE s'utilisent PAS (ou rarement) dans la structure \zh{把}]
\begin{itemize}
\item \textbf{Verbes d'état / sentiment / perception :} \zh{喜欢、爱、恨、知道、认识、明白、懂、觉得、认为、希望、想 (vouloir)}. \textit{On ne « manipule » pas un objet avec ces verbes.}
\item \textbf{Verbes statifs / adjectifs verbaux purs :} \zh{是、在、有、很、好、大}. \textit{Exclu :} \zh{*我把他很好。}
\item \textbf{Verbes d'action intransitifs sans objet patient :} \zh{去、来、走、跑、睡觉、起床}. \textit{Exclu :} \zh{*我把公园去了。} (Mais \zh{我把路走完了} « J'ai fini de parcourir le chemin » est possible car \zh{路} devient patient).
\item \textbf{Verbes à double objet (donner) :} \zh{给、送、还、借、教} préfèrent la structure \zh{给字句} (Sujet + 给 + Destinataire + Objet + Verbe). \textit{Ex :} \zh{我给了他一本书} (et non \zh{*我把书给他了} — bien que toléré à l'oral, la norme HSK préfère \zh{给} prépositionnel).
\end{itemize}
\end{attention}

\courssub{Étape 5 : Comparaison français / chinois — Pièges des francophones}

\begin{attention}[Erreurs fréquentes et corrections]
\begin{center}
\begin{tabularx}{\linewidth}{|X|X|X|}
\hline
\rowcolor{popblueL} Erreur typique (influence français) & Analyse & Correction \zh{把字句} \\
\hline
\zh{* 请排桌子。} (SVO simple) & Français : « Range les tables. » Verbe nu, pas de focus sur l'objet patient. & \zh{请把桌子排好。} (Verbe + \zh{好}) \\
\hline
\zh{* 我把书放好了没？} & Négation/interrogation mal placée après \zh{把}. & \zh{我没把书放好。} / \zh{你把书放好了吗？} \\
\hline
\zh{* 把一本书放好。} & Objet indéfini (\zh{一本}). \zh{把} exige la détermination. & \zh{把那本书放好。} / \zh{把你的书放好。} \\
\hline
\zh{* 我把作业写。} & Verbe nu (\zh{写}). Manque \zh{了/完/好}. & \zh{我把作业写完了。} / \zh{我把作业写好了。} \\
\hline
\zh{* 老师把我们批评。} & Verbe \zh{批评} prend souvent un complément humain, mais structure \zh{被} (passif) plus naturelle si focus sur nous. & \zh{我们被老师批评了。} (Passif) ou \zh{老师批评了我们。} (Actif SVO) \\
\hline
\end{tabularx}
\end{center}

\textbf{Astuce mnémotechnique :} \zh{把} = « Prends [l'objet] ET fais-en qqch jusqu'au bout ». L'objet est « pris en main » (\zh{把} = main \zh{手} + \zh{巴}), d'où l'obligation de résultat.
\end{attention}

\courssub{Étape 6 : Entraînement — Activité « Mission rangement » (50 min)}

\courssub{Ressources lexicales mobilisables}

\begin{center}
\begin{tabularx}{\linewidth}{|X|X|X|X|}
\hline
\rowcolor{popblueL} Caractères & Pinyin & Nature & Traduction \\
\hline
打扫 & dǎsǎo & verbe & nettoyer, faire le ménage \\
\hline
教室 & jiàoshì & nom & salle de classe \\
\hline
桌子 & zhuōzi & nom & table \\
\hline
椅子 & yǐzi & nom & chaise \\
\hline
黑板 & hēibǎn & nom & tableau noir \\
\hline
擦 & cā & verbe & essuyer \\
\hline
干净 & gānjìng & adjectif & propre \\
\hline
垃圾 & lājī & nom & ordures, déchets \\
\hline
倒 & dào & verbe & verser, jeter (liquides/ordures) \\
\hline
排 & pái & verbe & aligner, ranger en rangs \\
\hline
关 & guān & verbe & fermer \\
\hline
窗户 & chuānghu & nom & fenêtre \\
\hline
检查 & jiǎnchá & nom/verbe & inspection ; inspecter \\
\hline
纸 & zhǐ & nom & papier \\
\hline
地上 & dìshang & locatif & par terre, sur le sol \\
\hline
扔 & rēng & verbe & jeter (objets solides) \\
\hline
笔记本 & bǐjìběn & nom & cahier \\
\hline
书包 & shūbāo & nom & cartable, sac d'école \\
\hline
里 & lǐ & locatif & dedans, à l'intérieur \\
\hline
上 & shang & locatif & dessus, haut (compl. direction) \\
\hline
掉 & diào & compl. rés. & tomber, se débarrasser de, partir \\
\hline
好 & hǎo & compl. rés. & bien, correctement, achevé \\
\hline
完 & wán & compl. rés. & finir, épuiser \\
\hline
\end{tabularx}
\end{center}

\courssub{Consignes des tâches}

\textbf{Tâche 1 — Compréhension et repérage (10 min).} Reprenez l'annonce de M. Zhang. Soulignez la phrase en \zh{把}. Identifiez : Objet (\zh{桌子}), Verbe (\zh{排}), Complément (\zh{好}). Expliquez en une phrase française pourquoi \zh{请排桌子} ne convient pas pour une consigne de rangement abouti.

\textbf{Tâche 2 — Association objet / action (10 min).} Associez chaque objet à une action \textbf{avec son complément} (colonne de droite). Vérifiez : verbe + complément obligatoire.
\begin{center}
\begin{tabular}{|l|l|}
\hline
\textbf{Objets} & \textbf{Actions possibles (Verbe + Complément)} \\
\hline
\zh{桌子} & \zh{排好} \\
\zh{椅子} & \zh{放好} / \zh{排好} \\
\zh{黑板} & \zh{擦干净} \\
\zh{书 / 笔记本} & \zh{放进书包里} / \zh{收好} \\
\zh{垃圾 / 纸} & \zh{倒掉} / \zh{扔掉} / \zh{捡起来} \\
\zh{门 / 窗户} & \zh{关上} / \zh{打开} \\
\hline
\end{tabular}
\end{center}

\textbf{Tâche 3 — Rédaction collaborative : Liste de 6 instructions (15 min).} Par groupe de 4. Rédigez 6 phrases commençant par \zh{请把}.
\textbf{Contraintes obligatoires :}
\begin{itemize}
\item Au moins 2 compléments de résultat différents (\zh{好、干净、完、掉、上}).
\item Au moins 1 localisation (\zh{放进…里、放在…上}).
\item 1 phrase négative correcte (\zh{请不要把… / 请别把…}).
\item Objets déterminés (contexte : \emph{nos} tables, \emph{le} tableau, \emph{vos} cahiers).
\end{itemize}

\textbf{Tâche 4 — Jeu de rôle : Exécution et prononciation (15 min).}
\begin{itemize}
\item Chef de groupe lit les 6 ordres. Les 3 autres miment l'action.
\item Changez de chef. Répétez.
\item \textbf{Focus tons :} \zh{bǎ (3)}, \zh{hǎo (3)}, \zh{gān (1) jìng (4)}, \zh{dào (4) diào (4)}, \zh{guān (1) shang (4/neutral)}, \zh{pái (2) hǎo (3)}.
\end{itemize}

\textbf{Tâche 5 — Rapport oral et écrit au professeur (10 min).}
\begin{itemize}
\item Oral : Chaque groupe annonce : \zh{老师，我们把桌子排好了，把黑板擦干净了……} (Structure : \zh{我们把} + O + V + \zh{了/Complément}).
\item Écrit : Rédigez 4 phrases de bilan dont \textbf{1 au négatif} (\zh{没/还没 + 把}).
\end{itemize}

\textbf{Tâche 6 — Transformation guidée (10 min).} Traduisez en chinois \textbf{avec \zh{把}} :
\begin{enumerate}
\item « J'ai mis les cahiers dans le cartable. »
\item « Ne jette pas les papiers par terre ! »
\item « Elle a essuyé le tableau (jusqu'à ce qu'il soit propre). »
\item « Nous n'avons pas encore fermé les fenêtres. »
\item « Rangez bien vos chaises, s'il vous plaît. »
\end{enumerate}

\courssub{Production attendue et corrigé commenté}

\begin{exoresolu}[Modèle de production — Groupe 3 (Tâches 3, 5, 6)]
Énoncé : Liste de 6 instructions + Rapport 4 phrases + 5 phrases de transformation.
\tcblower
\textbf{Liste d'instructions (Tâche 3) :}\\
1. \zh{请把桌子排好。} \emph{Qǐng bǎ zhuōzi pái hǎo.} « Alignez les tables, s'il vous plaît. »\\
2. \zh{请把椅子放好。} \emph{Qǐng bǎ yǐzi fàng hǎo.} « Rangez bien les chaises. »\\
3. \zh{请把黑板擦干净。} \emph{Qǐng bǎ hēibǎn cā gānjìng.} « Essuyez le tableau jusqu'à ce qu'il soit propre. »\\
4. \zh{请把笔记本放进书包里。} \emph{Qǐng bǎ bǐjìběn fàng jìn shūbāo lǐ.} « Mettez les cahiers dans les cartables. »\\
5. \zh{请把垃圾倒掉。} \emph{Qǐng bǎ lājī dào diào.} « Jetez les ordures. »\\
6. \zh{请不要把纸扔在地上。} \emph{Qǐng bùyào bǎ zhǐ rēng zài dìshang.} « Ne jetez pas les papiers par terre. »\\[6pt]
\textbf{Rapport au professeur (Tâche 5) :}\\
1. \zh{老师，我们把桌子排好了，把椅子也放好了。} \emph{Lǎoshī, wǒmen bǎ zhuōzi pái hǎo le, bǎ yǐzi yě fàng hǎo le.} « Monsieur, nous avons aligné les tables et bien rangé les chaises aussi. »\\
2. \zh{我们把黑板擦干净了。} \emph{Wǒmen bǎ hēibǎn cā gānjìng le.} « Nous avons essuyé le tableau (propre). »\\
3. \zh{我们把笔记本都放进书包里了。} \emph{Wǒmen bǎ bǐjìběn dōu fàng jìn shūbāo lǐ le.} « Nous avons mis tous les cahiers dans les cartables. »\\
4. \zh{我们还没把窗户关上。} \emph{Wǒmen hái méi bǎ chuānghu guān shang.} « Nous n'avons pas encore fermé les fenêtres. »\\[6pt]
\textbf{Transformation (Tâche 6) :}\\
(a) \zh{我把笔记本放进书包里了。} \emph{Wǒ bǎ bǐjìběn fàng jìn shūbāo lǐ le.} \\
(b) \zh{别把纸扔在地上！} \emph{Bié bǎ zhǐ rēng zài dìshang!} (ou \zh{请不要把纸扔在地上。}) \\
(c) \zh{她把黑板擦干净了。} \emph{Tā bǎ hēibǎn cā gānjìng le.} \\
(d) \zh{我们还没把窗户关上。} \emph{Wǒmen hái méi bǎ chuānghu guān shang.} \\
(e) \zh{请把椅子排好。} \emph{Qǐng bǎ yǐzi pái hǎo.} \\[6pt]
\textbf{Commentaire linguistique :}
\begin{itemize}
\item \textbf{Compléments obligatoires :} \zh{好} (1, 2, 6), \zh{干净} (3), \zh{进…里} (4), \zh{掉} (5), \zh{上} (d, e). Aucun verbe nu.
\item \textbf{Objets déterminés :} Contextuels (\zh{桌子、黑板、笔记本、垃圾、纸、窗户, 椅子}) — pas d'indéfini \zh{一本/一个}.
\item \textbf{Négation :} \zh{没} (d) et \zh{别/不要} (6, b) placés \textbf{avant} \zh{把}.
\item \textbf{Adverbe \zh{都} / \zh{也} :} Se place \textbf{après l'objet introduit par \zh{把}}, avant le verbe (phrase 3 rapport : \zh{把笔记本都放…}).
\item \textbf{Tons :} Attention au \zh{3e ton} \zh{bǎ} + \zh{3e ton} \zh{hǎo} $\rightarrow$ \zh{bá hǎo} (sandhi) à l'oral, s'écrire \zh{bǎ hǎo}.
\end{itemize}
\end{exoresolu}

\courssub{Exercices de consolidation (Devoirs / Travail personnel)}

\begin{exercice}[Entraînement écrit — Niveau HSK 3]
\begin{enumerate}
\item \textbf{Transformez en phrase \zh{把字句} (ajoutez le complément nécessaire) :}
      \begin{enumerate}
      \item 我关门了。 $\rightarrow$ \zh{我把门关上了。}
      \item 他吃苹果。 $\rightarrow$ \zh{他把苹果吃完了/吃掉了。}
      \item 请开窗户。 $\rightarrow$ \zh{请把窗户打开。}
      \item 她洗衣服。 $\rightarrow$ \zh{她把衣服洗好了/洗干净了。}
      \item 我们扫地。 $\rightarrow$ \zh{我们把地扫干净了。}
      \end{enumerate}
\item \textbf{Corrigez les erreurs dans les phrases suivantes :}
      \begin{enumerate}
      \item \zh{* 请把书放。} $\rightarrow$ \zh{请把书放好 / 放进书包里。}
      \item \zh{* 我把一张纸扔了。} $\rightarrow$ \zh{我把那张纸扔掉了。} (objet déterminé)
      \item \zh{* 他没把作业做完。} $\rightarrow$ \textbf{Correct.} (Négation \zh{没} avant \zh{把}).
      \item \zh{* 别把灯关。} $\rightarrow$ \zh{别把灯关上 / 关掉。}
      \item \zh{* 老师把我们表扬了。} $\rightarrow$ \zh{老师表扬了我们。} (Verbe \zh{表扬} préfère SVO ou \zh{被} ; \zh{把} possible mais moins naturel pour action sociale).
      \end{enumerate}
\item \textbf{Traduisez en chinois (utilisez \zh{把}) :}
      \begin{enumerate}
      \item Ferme la porte, s'il te plaît. (\zh{请把门关上。})
      \item J'ai déjà fini les devoirs. (\zh{我已经把作业做完了。})
      \item Ne casse pas la tasse ! (\zh{别把杯子打碎了 / 打破了。})
      \item Ils ont nettoyé la salle de classe. (\zh{他们把教室打扫干净了。})
      \item Mettez les dictionnaires sur l'étagère. (\zh{请把词典放在书架上。})
      \end{enumerate}
\end{enumerate}
\end{exercice}

\begin{corrige}[Exercices de consolidation]
\begin{enumerate}
\item Voir flèches dans l'énoncé. Points clés : \zh{关上、吃完/掉、打开、洗好/干净、扫干净}.
\item Explications :
      \begin{enumerate}
      \item Verbe nu \zh{放} $\rightarrow$ ajouter \zh{好} ou locatif.
      \item \zh{一张纸} indéfini $\rightarrow$ \zh{那张纸} ou \zh{这张纸} ; \zh{扔} nu $\rightarrow$ \zh{扔掉}.
      \item Phrase correcte : ordre \zh{没 + 把} respecté.
      \item \zh{关} nu $\rightarrow$ \zh{关上/掉}.
      \item \zh{表扬} (louer) : action sociale, pas manipulation physique d'un patient. Préférer SVO : \zh{老师表扬了我们} ou passif \zh{我们被老师表扬了}.
      \end{enumerate}
\item Traductions :
      \begin{enumerate}
      \item \zh{请把门关上。}
      \item \zh{我已经把作业做完了。} (\zh{已经} avant \zh{把} ou après sujet).
      \item \zh{别把杯子打碎了 / 打破了。} (Complément de résultat \zh{碎/破}).
      \item \zh{他们把教室打扫干净了。}
      \item \zh{请把词典放在书架上。} (Localisation \zh{在…上}).
      \end{enumerate}
\end{enumerate}
\end{corrige}

\courssub{Bilan de la leçon}

La structure \zh{把字句} est un pivot majeur du chinois intermédiaire (HSK 3-4). Elle permet de passer de la description d'action (SVO) à la \textbf{narration de manipulation aboutie} sur un objet précis. Les trois piliers — \textbf{Objet déterminé}, \textbf{Verbe complété}, \textbf{Négation/Auxiliaires avant \zh{把}} — sont non négociables.

\begin{savaistu}[Le « grand ménage » (大扫除) : Chine vs Cameroun]
En Chine, le \zh{大扫除} (dà sǎochú) est une institution hebdomadaire (souvent le vendredi) dans toutes les écoles, de la primaire à l'université. Les élèves balayent, lavent les sols, essuient les fenêtres, rangent les bureaux \textbf{eux-mêmes}, sans personnel d'entretien. C'est un moment d'éducation civique (\zh{劳动教育} láodòng jiàoyù) : responsabilité collective, humilité, soin du bien commun. Au Cameroun, le « grand ménage » du vendredi au lycée (souvent avant l'inspection ou la rentrée) partage cet esprit : la classe est \emph{notre} espace, nous en prenons soin. La différence ? En Chine, c'est quotidiennement ritualisé (15-20 min chaque jour + grand ménage hebdo) ; au Cameroun, c'est souvent ponctuel (veille d'inspection, fin de trimestre). La structure \zh{把} est l'outil linguistique parfait pour cet univers : \zh{把地扫干净、把窗户擦亮、把桌子排整齐} — chaque verbe porte la trace de l'effort accompli sur l'objet.
\end{savaistu}

\courssub{Prolongement personnel (Évaluation formative)}

Rédigez un court paragraphe (5-6 phrases) en chinois pour décrire le rangement de votre chambre ou de votre espace de travail à la maison hier soir. Utilisez \zh{把} au moins 3 fois. Intégrez : un complément de résultat (\zh{好/干净/完}), un complément de direction/localisation (\zh{进/上/里}), une phrase au négatif (\zh{没把…/还没把…}). Exemple de début : \zh{昨天晚上，我把衣服叠好了，把书放回书架上……}

\newpage

\coursec{Méthodes et exercices résolus}

\courssub{Observation guidée : découvrir la phrase 把}

\begin{textebox}[Dialogue : après le cours, au lycée de Yaoundé]
A : 艾玛，请你把窗户关上，要下雨了。\\
\emph{Àimǎ, qǐng nǐ bǎ chuānghu guān shàng, yào xià yǔ le.}\\
« Emma, ferme la fenêtre, s'il te plaît, il va pleuvoir. »\\[4pt]
B : 好的。我已经把黑板擦干净了。\\
\emph{Hǎo de. Wǒ yǐjīng bǎ hēibǎn cā gānjìng le.}\\
« D'accord. J'ai déjà essuyé le tableau (il est propre). »\\[4pt]
A : 太好了！别忘了把作业本交给老师。\\
\emph{Tài hǎo le! Bié wàng le bǎ zuòyèběn jiāo gěi lǎoshī.}\\
« Parfait ! N'oublie pas de remettre les cahiers au professeur. »\\[4pt]
B : 放心，我没把书落在教室里。\\
\emph{Fàngxīn, wǒ méi bǎ shū là zài jiàoshì li.}\\
« Ne t'inquiète pas, je n'ai pas laissé mes livres dans la salle. »
\end{textebox}

\textbf{Questions d'observation.}
\begin{enumerate}
\item Dans chaque phrase soulignée par la présence de \zh{把}, quel est l'\emph{objet} concerné ? (la fenêtre, le tableau, les cahiers, les livres)
\item Où se place cet objet : avant ou après le verbe ?
\item Que trouve-t-on \emph{après} chaque verbe ? (\zh{关上}, \zh{擦干净}, \zh{交给老师}, \zh{落在教室里})
\item Où sont placés \zh{已经} et \zh{没} par rapport à \zh{把} ?
\end{enumerate}

\textbf{Constat.} Dans toutes ces phrases, l'objet est \emph{précis} et \emph{connu}, il est placé \emph{avant} le verbe grâce à \zh{把}, et le verbe n'est jamais seul : il est suivi d'un élément qui dit ce que l'objet est devenu. C'est la \cle{phrase 把} (\zh{把字句} \emph{bǎzìjù}).

\begin{definition}[La phrase 把 (把字句)]
La \cle{phrase 把} est une construction dans laquelle la préposition \zh{把} (\emph{bǎ}) introduit l'objet \emph{déterminé} du verbe et le place \emph{avant} celui-ci. Elle met l'accent sur ce que le sujet \emph{fait subir} à l'objet : déplacement, transformation, achèvement, changement d'état.

\textbf{Schéma général :} Sujet + (négation / adverbe / modal) + 把 + Objet déterminé + Verbe + complément obligatoire.
\end{definition}

\begin{center}
\begin{tabularx}{\linewidth}{|X|X|X|X|}
\hline
\rowcolor{popblueL} Caractères & Pinyin & Nature & Traduction \\
\hline
把 & bǎ & préposition & (marque de l'objet déplacé) \\
\hline
窗户 & chuānghu & nom & fenêtre \\
\hline
关 & guān & verbe & fermer \\
\hline
黑板 & hēibǎn & nom & tableau noir \\
\hline
擦 & cā & verbe & essuyer, frotter \\
\hline
干净 & gānjìng & adjectif & propre \\
\hline
作业本 & zuòyèběn & nom & cahier de devoirs \\
\hline
交给 & jiāo gěi & verbe & remettre à \\
\hline
落 & là & verbe & laisser, oublier (qqch. quelque part) \\
\hline
教室 & jiàoshì & nom & salle de classe \\
\hline
\end{tabularx}
\end{center}

\begin{popfigure}
\begin{tikzpicture}[
  bloc/.style={rectangle, rounded corners=3pt, draw=popblue, fill=popblueL, align=center, minimum height=0.9cm, font=\small},
  fleche/.style={-{Stealth[length=2.5mm]}, thick, popdark}
]
\node[bloc, align=center] (s) at (0,0){Sujet\\我};
\node[bloc, fill=poporangeL, draw=poporange, align=center] (ba) at (2.6,0){把\\bǎ};
\node[bloc, align=center] (o) at (5.2,0){Objet déterminé\\那本书};
\node[bloc, align=center] (v) at (8.2,0){Verbe\\放};
\node[bloc, fill=popgreenL, draw=popgreen, align=center] (c) at (11.2,0){Complément\\在桌子上};
\draw[fleche] (s) -- (ba);
\draw[fleche] (ba) -- (o);
\draw[fleche] (o) -- (v);
\draw[fleche] (v) -- (c);
\node[align=center, font=\small\itshape, popmuted] at (5.6,-1.3) {我把那本书放在桌子上。 — « J'ai posé ce livre sur la table. »};
\end{tikzpicture}
\legende{Schéma de la phrase 把 : l'objet passe devant le verbe, le complément est obligatoire.}
\end{popfigure}

\courssub{Fiche méthode 1 : Construire une phrase 把 correcte}

\begin{methode}[Construire une phrase 把 : la démarche en 5 étapes]
La phrase 把 (\emph{bǎzìjù}) sert à dire ce qu'on \emph{fait subir} à un objet précis : on le déplace, on le transforme, on le termine. Pour la construire sans faute :

\begin{enumerate}
\item \textbf{Identifier l'objet déterminé.} Le nom après 把 doit être connu ou précis (avec 这/那 + classificateur, un possessif, ou déjà mentionné dans le contexte). On ne dit pas \zh{把一个苹果吃了} si l'on parle d'une pomme quelconque.
\item \textbf{Placer la structure de base :} Sujet + 把 + Objet + Verbe + complément obligatoire.
\item \textbf{Ne jamais laisser le verbe seul.} Après le verbe, il faut toujours un élément : 了, un complément directionnel (来/去), un complément de résultat (完/好/错), un complément de lieu (在/到/给) ou un redoublement du verbe.
\item \textbf{Placer la négation et les adverbes AVANT 把 :} 没(有) + 把, 不 + 把, 已经 + 把. Jamais entre 把 et le verbe.
\item \textbf{Vérifier le sens :} la phrase 把 répond à la question « Qu'a-t-on fait de l'objet ? ». Si le verbe n'affecte pas l'objet (aimer, voir, savoir), la phrase 把 est impossible.
\end{enumerate}

\textbf{Schéma à retenir :} Sujet + (négation / adverbe) + 把 + Objet déterminé + Verbe + 了 / complément.
\end{methode}

\begin{exoresolu}[Construction : ranger la classe]
\textbf{Énoncé (type examen).} Mets les mots dans l'ordre pour construire une phrase 把 correcte, puis traduis :

a) 把 / 桌子 / 老师 / 干净 / 擦 / 了

b) 书 / 把 / 我 / 在 / 放 / 桌子上

c) 把 / 他 / 没 / 作业 / 完 / 做
\tcblower
\textbf{Solution détaillée.}

\textbf{a)} Étape 1 : le sujet est 老师 (le professeur). Étape 2 : l'objet déterminé est 桌子 (la table). Étape 3 : le verbe 擦 (essuyer) ne peut pas rester seul ; le complément de résultat 干净 (propre) + 了 convient.

\zh{老师把桌子擦干净了。} \emph{Lǎoshī bǎ zhuōzi cā gānjìng le.} « Le professeur a essuyé la table (elle est propre). »

\textbf{b)} Sujet : 我 ; objet : 书 ; verbe 放 + complément de lieu 在桌子上.

\zh{我把书放在桌子上。} \emph{Wǒ bǎ shū fàng zài zhuōzi shang.} « J'ai posé le livre sur la table. »

\textbf{c)} La négation 没 se place AVANT 把, jamais après. Le complément 完 (finir) accompagne 做.

\zh{他没把作业做完。} \emph{Tā méi bǎ zuòyè zuò wán.} « Il n'a pas fini ses devoirs. »

\textbf{Conclusion.} Dans les trois cas, l'objet est déterminé, le verbe est suivi d'un complément, et la négation précède 把. C'est le réflexe à acquérir.
\end{exoresolu}

\courssub{Fiche méthode 2 : Choisir le bon complément après le verbe}

\begin{methode}[Quel complément après le verbe dans une phrase 把 ?]
Le verbe d'une phrase 把 ne peut jamais être « nu ». Choisis le complément selon le sens :

\begin{enumerate}
\item \textbf{Résultat de l'action} (l'objet change d'état) : verbe + 完 (finir), 好 (bien fait), 错 (faux), 干净 (propre), 坏 (abîmé). Ex. : \zh{把饭吃完} « finir de manger le riz ».
\item \textbf{Déplacement} : verbe + 到 + lieu (amener à), 在 + lieu (placer à), 给 + personne (donner à). Ex. : \zh{把信寄到中国} « envoyer la lettre en Chine ».
\item \textbf{Direction} : verbe + 来 / 去 / 进来 / 出去. Ex. : \zh{把书拿出来} « sortir le livre ».
\item \textbf{Transformation en...} : 把 + A + 变成/写成/翻译成 + B. Ex. : \zh{把句子翻译成法语} « traduire la phrase en français ».
\item \textbf{Action brève ou tentative} : verbe redoublé (看看, 想一想) ou verbe + 一下. Ex. : \zh{把房间看一看} « jeter un coup d'œil à la chambre ».
\end{enumerate}

\textbf{Réflexe :} demande-toi « que devient l'objet après l'action ? » — la réponse est ton complément.
\end{methode}

\begin{center}
\begin{tabularx}{\linewidth}{|X|X|X|}
\hline
\rowcolor{popblueL} Type de complément & Exemple (caractères + pinyin) & Traduction \\
\hline
Résultat (完/好/干净) & 把饭吃完 \emph{bǎ fàn chī wán} & finir de manger le riz \\
\hline
Lieu (在/到/给) & 把信寄到中国 \emph{bǎ xìn jì dào Zhōngguó} & envoyer la lettre en Chine \\
\hline
Direction (出来/进去) & 把书拿出来 \emph{bǎ shū ná chūlái} & sortir le livre \\
\hline
Transformation (成) & 把句子翻译成法语 \emph{bǎ jùzi fānyì chéng Fǎyǔ} & traduire la phrase en français \\
\hline
Action brève (一下) & 把问题想一下 \emph{bǎ wèntí xiǎng yíxià} & réfléchir un peu à la question \\
\hline
\end{tabularx}
\end{center}

\begin{exoresolu}[Choix du complément : à la maison à Yaoundé]
\textbf{Énoncé.} Complète chaque phrase 把 avec le complément qui convient parmi : 完, 到学校里, 成法语, 出来, 一下.

a) 妈妈把饭做\_\_\_\_了，我们可以吃了。

b) 请你把这句话翻译\_\_\_\_。

c) 司机把学生送\_\_\_\_。

d) 他把字典拿\_\_\_\_，放在桌子上。

e) 你把这个问题想\_\_\_\_，再回答。
\tcblower
\textbf{Solution détaillée.}

\textbf{a)} Le riz est prêt : résultat de l'action → 完. \zh{妈妈把饭做完了，我们可以吃了。} \emph{Māma bǎ fàn zuò wán le, wǒmen kěyǐ chī le.} « Maman a fini de préparer le repas, on peut manger. »

\textbf{b)} Transformation d'une langue dans une autre → 成法语. \zh{请你把这句话翻译成法语。} \emph{Qǐng nǐ bǎ zhè jù huà fānyì chéng Fǎyǔ.} « Traduis cette phrase en français, s'il te plaît. »

\textbf{c)} Déplacement vers un lieu → 到学校里. \zh{司机把学生送到学校里。} \emph{Sījī bǎ xuésheng sòng dào xuéxiào li.} « Le chauffeur a déposé les élèves à l'école. »

\textbf{d)} Mouvement vers l'extérieur → 出来. \zh{他把字典拿出来，放在桌子上。} \emph{Tā bǎ zìdiǎn ná chūlái, fàng zài zhuōzi shang.} « Il a sorti le dictionnaire et l'a posé sur la table. »

\textbf{e)} Action brève, « réfléchis un peu » → 一下. \zh{你把这个问题想一下，再回答。} \emph{Nǐ bǎ zhège wèntí xiǎng yíxià, zài huídá.} « Réfléchis un peu à cette question, puis réponds. »

\textbf{Conclusion.} Le complément exprime toujours le devenir de l'objet : fini, transformé, déplacé, sorti ou simplement « un peu traité ».
\end{exoresolu}

\courssub{Fiche méthode 3 : Transformer une phrase SVO en phrase 把}

\begin{methode}[Transformation SVO → 把 : la recette]
Pour transformer une phrase classique Sujet + Verbe + Objet en phrase 把 :

\begin{enumerate}
\item \textbf{Vérifier que la transformation est possible :} le verbe doit pouvoir « affecter » l'objet (manger, laver, fermer, envoyer, oublier...). Verbes impossibles avec 把 : 有, 是, 喜欢, 知道, 看见.
\item \textbf{Rendre l'objet déterminé :} ajoute 这/那 + classificateur ou un possessif si nécessaire (\zh{苹果} → \zh{那个苹果}).
\item \textbf{Déplacer l'objet} entre 把 et le verbe : Sujet + 把 + Objet + Verbe.
\item \textbf{Ajouter un complément} après le verbe (了, 完, 干净, 在...) : c'est obligatoire.
\item \textbf{Replacer négation et adverbes avant 把.}
\end{enumerate}

\textbf{Formule :} S + V + O → S + 把 + O (déterminé) + V + complément.
\end{methode}

\begin{exoresolu}[Transformation : phrases du quotidien]
\textbf{Énoncé.} Transforme les phrases suivantes en phrases 把, puis traduis :

a) 我喝了那杯水。

b) 妹妹洗干净了衣服。

c) 他忘了我的电话号码。
\tcblower
\textbf{Solution détaillée.}

\textbf{a)} L'objet 那杯水 est déjà déterminé. On le déplace après 把 ; le verbe 喝 garde 了.

\zh{我把那杯水喝了。} \emph{Wǒ bǎ nà bēi shuǐ hē le.} « J'ai bu ce verre d'eau (jusqu'au bout). »

Justification : 喝 affecte l'objet (l'eau disparaît), donc la transformation est possible. Le classificateur 杯 est conservé avec l'objet.

\textbf{b)} L'objet 衣服 devient déterminé : 她的衣服 ou 这些衣服. Le complément de résultat 干净 + 了 reste après le verbe.

\zh{妹妹把她的衣服洗干净了。} \emph{Mèimei bǎ tā de yīfu xǐ gānjìng le.} « La petite sœur a lavé ses vêtements (ils sont propres). »

\textbf{c)} 忘 (oublier) peut affecter un objet précis : le numéro est déterminé (我的电话号码). On ajoute 了 après le verbe.

\zh{他把我的电话号码忘了。} \emph{Tā bǎ wǒ de diànhuà hàomǎ wàng le.} « Il a oublié mon numéro de téléphone. »

\textbf{Conclusion.} La transformation met l'accent sur le sort de l'objet : l'eau est bue, les vêtements sont propres, le numéro est oublié. Sans complément après le verbe, la phrase serait fautive.
\end{exoresolu}

\courssub{Fiche méthode 4 : Négation, adverbes et verbes modaux dans la phrase 把}

\begin{methode}[Où placer 没, 不, 已经, 要, 会 ?]
Règle d'or : \textbf{tout ce qui n'est pas le verbe principal se place AVANT 把} (ou avant le sujet pour certains adverbes de temps).

\begin{enumerate}
\item \textbf{Négation du passé :} 没(有) + 把 + O + V + complément. \zh{我没有把窗户关上。} « Je n'ai pas fermé la fenêtre. »
\item \textbf{Négation présente/future ou impératif :} 不 / 别 + 把. \zh{别把手机放在桌子上。} « Ne pose pas ton téléphone sur la table. »
\item \textbf{Adverbes :} 已经, 都, 也, 马上 se placent avant 把. \zh{他已经把作业做完了。} « Il a déjà fini ses devoirs. »
\item \textbf{Verbes modaux :} 想, 要, 会, 能, 应该 se placent avant 把. \zh{我要把这个消息告诉他。} « Je veux lui annoncer cette nouvelle. »
\item \textbf{Interdit :} placer 没, 不, 已经 entre 把 et le verbe (\zh{把作业没做完} est FAUX).
\end{enumerate}
\end{methode}

\begin{exoresolu}[Correction de phrases fautives]
\textbf{Énoncé.} Les phrases suivantes contiennent chacune une erreur de place. Corrige-les et traduis :

a) 他把作业没做完。

b) 我把这个单词已经记住了。

c) 你别把这件事告诉他，好吗？
\tcblower
\textbf{Solution détaillée.}

\textbf{a)} FAUX : la négation 没 est coincée entre 把 et le verbe. Elle doit précéder 把.

Correction : \zh{他没把作业做完。} \emph{Tā méi bǎ zuòyè zuò wán.} « Il n'a pas fini ses devoirs. »

\textbf{b)} FAUX : l'adverbe 已经 doit se placer avant 把.

Correction : \zh{我已经把这个单词记住了。} \emph{Wǒ yǐjīng bǎ zhège dāncí jì zhù le.} « J'ai déjà mémorisé ce mot. » (记住 = retenir, mémoriser ; 住 est le complément de résultat.)

\textbf{c)} CORRECT : 别 (impératif négatif) est bien placé avant 把. \zh{你别把这件事告诉他，好吗？} \emph{Nǐ bié bǎ zhè jiàn shì gàosu tā, hǎo ma?} « Ne lui raconte pas cette affaire, d'accord ? » (件 = classificateur des affaires, événements.)

\textbf{Conclusion.} Le réflexe de l'examen : repérer 把, puis vérifier que négation, adverbes et modaux sont tous à sa gauche.
\end{exoresolu}

\courssub{Fiche méthode 5 : Utiliser 把 dans un court texte rédigé}

\begin{methode}[Rédiger un paragraphe avec des phrases 把]
Pour une production écrite (récit, message, description d'actions) :

\begin{enumerate}
\item \textbf{Repère les actions qui « disposent » d'un objet} : ranger, nettoyer, préparer, envoyer, ouvrir, fermer. Ce sont tes candidates pour 把.
\item \textbf{Alterne} phrases 把 et phrases SVO simples pour un texte naturel ; n'utilise pas 把 à chaque phrase.
\item \textbf{Enchaîne avec des marqueurs temporels} : 先... 然后... 最后... (d'abord... ensuite... enfin...).
\item \textbf{Varie les compléments} : 完, 好, 干净, 在..., 给... pour éviter les répétitions.
\item \textbf{Relis avec la check-list :} objet déterminé ? verbe suivi d'un complément ? négation/adverbes avant 把 ?
\end{enumerate}

\textbf{Modèle de trame :} 今天早上，我先 + 把 + O + V + 完，然后 + 把 + O + V + 在/给...，最后 + phrase simple.
\end{methode}

\begin{exoresolu}[Production écrite : le samedi matin à Douala]
\textbf{Énoncé.} Rédige un court paragraphe (4 à 6 phrases) intitulé 星期六早上 « Samedi matin », en utilisant au moins trois phrases 把 différentes. Mots suggérés : 房间, 打扫, 干净, 衣服, 洗, 作业, 做完, 早饭, 准备, 给.
\tcblower
\textbf{Solution détaillée.}

\textbf{Raisonnement.} Les actions du matin affectent des objets précis : la chambre (la nettoyer), les vêtements (les laver), les devoirs (les finir), le petit-déjeuner (le préparer pour quelqu'un). Chacune donne une phrase 把 avec un complément différent : 干净, 完, 给.

\textbf{Texte rédigé :}

\zh{星期六早上，我七点起床。我先把房间打扫干净了，然后把衣服洗好了。妈妈把早饭准备好了，我们一起吃饭。吃完饭，我把作业做完了。下午，我把借的书还给同学了。}

\emph{Xīngqīliù zǎoshang, wǒ qī diǎn qǐchuáng. Wǒ xiān bǎ fángjiān dǎsǎo gānjìng le, ránhòu bǎ yīfu xǐ hǎo le. Māma bǎ zǎofàn zhǔnbèi hǎo le, wǒmen yìqǐ chī fàn. Chī wán fàn, wǒ bǎ zuòyè zuò wán le. Xiàwǔ, wǒ bǎ jiè de shū huán gěi tóngxué le.}

« Samedi matin, je me lève à sept heures. D'abord, je nettoie ma chambre (elle est propre), ensuite je lave mes vêtements. Maman a préparé le petit-déjeuner, nous mangeons ensemble. Après le repas, je finis mes devoirs. L'après-midi, je rends à mon camarade le livre que j'avais emprunté. »

\textbf{Justifications grammaticales.} Quatre phrases 把, quatre compléments différents : 干净 (résultat), 好 (résultat), 完 (achèvement), 还给 (rendre à). Les objets sont tous déterminés par le contexte (ma chambre, mes vêtements, mes devoirs, le livre emprunté). Les marqueurs 先, 然后 structurent le récit.

\textbf{Conclusion.} Un texte naturel mélange phrases 把 (actions sur un objet) et phrases simples (se lever, manger ensemble) : c'est ce que l'examinateur attend.
\end{exoresolu}

\begin{attention}[Les 5 erreurs qui coûtent des points à l'examen]
\begin{enumerate}
\item \textbf{Le verbe laissé seul.} \zh{我把饭吃。} est FAUX : après le verbe, il faut toujours 了 ou un complément. Écris \zh{我把饭吃了。} ou \zh{我把饭吃完了。}
\item \textbf{La négation mal placée.} \zh{我把作业没做完。} est FAUX : 没/不/别 se placent AVANT 把 → \zh{我没把作业做完。}
\item \textbf{L'objet indéterminé.} \zh{我把一个苹果吃了。} est maladroit : l'objet de 把 doit être précis → \zh{我把那个苹果吃了。} Pense à ajouter 这/那 + classificateur.
\item \textbf{Le verbe incompatible.} Les verbes d'état ou de perception (是, 有, 喜欢, 知道, 看见) ne prennent pas 把 : \zh{我把这首歌喜欢。} est impossible. Garde la structure SVO : \zh{我喜欢这首歌。}
\item \textbf{Le calque du français.} En français on dit « je ferme la porte » (SVO) ; en chinois, dès que l'action change l'état de l'objet, la phrase 把 est souvent plus naturelle : \zh{请把门关上。} Ne traduis pas mot à mot, et n'oublie pas le complément 上 qui correspond à « fermer ».
\end{enumerate}
\end{attention}

\begin{aretenir}[L'essentiel de la leçon]
\begin{itemize}
\item La phrase 把 déplace l'objet \emph{déterminé} avant le verbe : Sujet + 把 + Objet + Verbe + complément.
\item Le verbe n'est \emph{jamais seul} : il est suivi de 了, d'un complément de résultat, de lieu, de direction ou d'un redoublement.
\item Négation (没/不/别), adverbes (已经, 都) et modaux (要, 会, 想) se placent \emph{avant} 把.
\item La phrase 把 exprime la \emph{disposition} de l'objet : ce qu'on lui fait subir, ce qu'il devient.
\item Verbes interdits avec 把 : 是, 有, 喜欢, 知道, 看见 (verbes d'état ou de perception).
\end{itemize}
\end{aretenir}

\begin{savaistu}[D'où vient 把 ?]
À l'origine, \zh{把} était un verbe signifiant « prendre dans la main » (on le retrouve encore dans \zh{把手} « poignée » et \zh{一把伞} « un parapluie », où 把 est le classificateur des objets qu'on tient en main). Au fil des siècles, ce verbe s'est grammaticalisé : « prendre le livre et le poser sur la table » est devenu \zh{把书放在桌子上}. La phrase 把 est l'une des constructions les plus typiques du chinois moderne : elle n'a aucun équivalent direct en français, ce qui en fait un point délicat mais très valorisé à l'examen. Au Cameroun comme en Chine, on l'utilise sans cesse dans la vie quotidienne : \zh{把门关上！} « Ferme la porte ! », \zh{把手机给我！} « Donne-moi le téléphone ! ».
\end{savaistu}

\begin{exercice}[À toi de jouer : entraînement type examen]
\textbf{Exercice 1.} Mets les mots dans l'ordre et traduis :

a) 把 / 我 / 自行车 / 好 / 修 / 了

b) 把 / 请 / 你 / 门 / 开 / 上 (attention au sens !)

c) 把 / 我们 / 都 / 课文 / 完 / 读 / 了

\textbf{Exercice 2.} Transforme en phrase 把 :

a) 弟弟吃完了米饭。

b) 我寄出了那封信。

\textbf{Exercice 3.} Traduis en chinois avec 把 :

a) « J'ai déjà rendu le dictionnaire à mon camarade. »

b) « Ne mets pas les chaussures sous le lit. » (sous = 下面)
\end{exercice}

\begin{corrige}[Exercices 1 à 3]
\textbf{Exercice 1.}

a) \zh{我把自行车修好了。} \emph{Wǒ bǎ zìxíngchē xiū hǎo le.} « J'ai réparé le vélo (il marche). » Complément de résultat 好.

b) \zh{请你把门关上。} \emph{Qǐng nǐ bǎ mén guān shàng.} « Ferme la porte, s'il te plaît. » Piège : 开上 n'existe pas pour une porte ; « fermer » se dit 关上.

c) \zh{我们都把课文读完了。} \emph{Wǒmen dōu bǎ kèwén dú wán le.} « Nous avons tous fini de lire le texte. » L'adverbe 都 précède 把.

\textbf{Exercice 2.}

a) \zh{弟弟把米饭吃完了。} \emph{Dìdi bǎ mǐfàn chī wán le.} « Le petit frère a fini le riz. »

b) \zh{我把那封信寄出去了。} \emph{Wǒ bǎ nà fēng xìn jì chūqù le.} « J'ai envoyé cette lettre. » (封 = classificateur des lettres ; complément directionnel 出去.)

\textbf{Exercice 3.}

a) \zh{我已经把字典还给同学了。} \emph{Wǒ yǐjīng bǎ zìdiǎn huán gěi tóngxué le.} L'adverbe 已经 avant 把 ; complément 还给 « rendre à ».

b) \zh{别把鞋放在床下面。} \emph{Bié bǎ xié fàng zài chuáng xiàmian.} Impératif négatif 别 avant 把 ; complément de lieu 在床下面.
\end{corrige}

\begin{experience}[Jeu de rôle : la veille de l'examen]
Par groupes de deux, jouez la scène suivante : un élève prépare sa chambre et ses affaires la veille d'une évaluation de chinois. Chaque élève doit produire au moins quatre phrases 把 avec des compléments différents (résultat, lieu, direction, 给).

\textbf{Exemples de départs :} \zh{我把课本放进书包里了。} « J'ai rangé le manuel dans le cartable. » — \zh{请把闹钟调到六点。} « Règle le réveil sur six heures, s'il te plaît. »

Le reste de la classe vérifie à voix haute la check-list : objet déterminé ? complément après le verbe ? adverbes avant 把 ?
\end{experience}

\newpage

\coursec{Exercices}

\courssub{Je vérifie mes connaissances}

\begin{exercice}[QCM : la bonne structure]
Pour chaque question, choisis la seule phrase correcte.
\begin{enumerate}
\item \begin{enumerate}
\item 我把饭吃在食堂里。
\item 我把饭吃完了。
\item 我吃完了把饭。
\end{enumerate}
\item \begin{enumerate}
\item 他把门没关好。
\item 他没把门关好。
\item 他把没门关好。
\end{enumerate}
\item \begin{enumerate}
\item 我把书放在桌子上。
\item 我放在桌子上把书。
\item 我把放书在桌子上。
\end{enumerate}
\item \begin{enumerate}
\item 请把一杯水喝。
\item 请把这杯水喝了。
\item 请喝把这杯水。
\end{enumerate}
\end{enumerate}
\end{exercice}

\begin{exercice}[Vrai ou faux ? Justifie]
Dis si les affirmations suivantes sont vraies ou fausses, puis justifie ta réponse en une phrase.
\begin{enumerate}
\item Dans une phrase en \zh{把}, le complément d'objet se place avant le verbe.
\item On peut dire : \zh{我把书看。}
\item La négation \zh{没（有）} se place avant \zh{把}.
\item Toutes les phrases Sujet + Verbe + Objet du français peuvent se transformer en phrase en \zh{把}.
\item \zh{把} s'emploie quand l'objet est déterminé (connu, précisé).
\end{enumerate}
\end{exercice}

\begin{exercice}[Vocabulaire : complète le tableau]
Complète le tableau suivant avec le pinyin et la traduction française de chaque mot (les caractères sont donnés).
\begin{center}
\begin{tabularx}{\linewidth}{|X|X|X|X|}
\hline
\rowcolor{popblueL} Caractères & Pinyin & Nature & Traduction \\
\hline
关 & \ldots & \ldots & \ldots \\
\hline
洗 & \ldots & \ldots & \ldots \\
\hline
放 & \ldots & \ldots & \ldots \\
\hline
锁 & \ldots & \ldots & \ldots \\
\hline
修 & \ldots & \ldots & \ldots \\
\hline
行李 & \ldots & \ldots & \ldots \\
\hline
\end{tabularx}
\end{center}
\end{exercice}

\begin{exercice}[Associe caractères et pinyin]
Relie chaque caractère ou mot (colonne de gauche) à son pinyin (colonne de droite). Attention : deux pinyin sont des intrus.
\begin{enumerate}
\item \zh{把} \hspace{1cm} a. fàng
\item \zh{完} \hspace{1cm} b. bǎ
\item \zh{放} \hspace{1cm} c. wán
\item \zh{干净} \hspace{1cm} d. gǎnjìng \hspace{1cm} e. guān \hspace{1cm} f. gānjìng \hspace{1cm} g. wǎn
\end{enumerate}
\end{exercice}

\courssub{Je m'entraîne}

\begin{exercice}[Traduis en français]
Traduis les phrases suivantes en français correct.
\begin{enumerate}
\item \zh{妈妈把饭做好了。}\\ \emph{Māma bǎ fàn zuò hǎo le.}
\item \zh{弟弟把果汁喝完了。}\\ \emph{Dìdi bǎ guǒzhī hē wán le.}
\item \zh{老师把书放在桌子上了。}\\ \emph{Lǎoshī bǎ shū fàng zài zhuōzi shang le.}
\item \zh{我把作业交给老师了。}\\ \emph{Wǒ bǎ zuòyè jiāo gěi lǎoshī le.}
\item \zh{请把窗户打开。}\\ \emph{Qǐng bǎ chuānghu dǎkāi.}
\end{enumerate}
\end{exercice}

\begin{exercice}[Transforme en phrase en 把]
Transforme les phrases suivantes en phrases en \zh{把}. Si la transformation est impossible, écris « impossible ».
\begin{enumerate}
\item \zh{他洗干净了衣服。}\\ \emph{Tā xǐ gānjìng le yīfu.}
\item \zh{我看完了这本书。}\\ \emph{Wǒ kàn wán le zhè běn shū.}
\item \zh{她关了门。}\\ \emph{Tā guān le mén.}
\item \zh{我们吃完了饭。}\\ \emph{Wǒmen chī wán le fàn.}
\item \zh{他修好了自行车。}\\ \emph{Tā xiū hǎo le zìxíngchē.}
\end{enumerate}
\end{exercice}

\begin{exercice}[Texte à trous]
Complète le texte suivant avec les mots de la liste : \zh{把}、\zh{在}、\zh{给}、\zh{完}、\zh{好}、\zh{了} (chaque mot peut être utilisé une ou plusieurs fois).

\begin{textebox}[Le matin de Xiao Ming]
早上，小明\ldots\ldots{}书包收拾\ldots\ldots{}，然后\ldots\ldots{}早饭吃\ldots\ldots{}\ldots\ldots{}。他\ldots\ldots{}钥匙放\ldots\ldots{}口袋里，\ldots\ldots{}自行车骑到学校。到了教室，他\ldots\ldots{}作业交\ldots\ldots{}老师\ldots\ldots{}。\\
\emph{Zǎoshang, Xiǎo Míng \ldots{} shūbāo shōushi \ldots{}, ránhòu \ldots{} zǎofàn chī \ldots{} \ldots{}. Tā \ldots{} yàoshi fàng \ldots{} kǒudài li, \ldots{} zìxíngchē qí dào xuéxiào. Dào le jiàoshì, tā \ldots{} zuòyè jiāo \ldots{} lǎoshī \ldots{}.}\\
« Le matin, Xiao Ming range son cartable, puis finit de manger son petit-déjeuner. Il met ses clés dans sa poche et va à l'école à vélo. Arrivé en classe, il remet ses devoirs au professeur. »
\end{textebox}
\end{exercice}

\begin{exercice}[Remets les mots dans l'ordre]
Remets les mots suivants dans le bon ordre pour former une phrase en \zh{把} correcte.
\begin{enumerate}
\item \zh{把 / 我 / 了 / 喝 / 水 / 完}
\item \zh{门 / 把 / 关 / 请 / 好 / 你}
\item \zh{把 / 桌子 / 书 / 在 / 他 / 放 / 上}
\item \zh{朋友 / 把 / 给 / 我 / 词典 / 了 / 借}
\item \zh{没 / 把 / 窗户 / 她 / 打开}
\end{enumerate}
\end{exercice}

\courssub{Je me prépare au Bac}

\begin{exercice}[Type Bac : compréhension et langue]
Lis le texte suivant, puis réponds aux questions.

\begin{textebox}[Avant le départ]
明天我们全家要去杜阿拉看奶奶。晚上，爸爸把行李放到车上了。妈妈把饭做好了，还把孩子明天要穿的衣服洗干净了。姐姐把火车票放进包里了。弟弟把果汁喝完了，又把玩具收好了。奶奶打电话说：« 你们把门锁好，路上小心！»\\
\emph{Míngtiān wǒmen quánjiā yào qù Dù'ālā kàn nǎinai. Wǎnshang, bàba bǎ xíngli fàng dào chē shang le. Māma bǎ fàn zuò hǎo le, hái bǎ háizi míngtiān yào chuān de yīfu xǐ gānjìng le. Jiějie bǎ huǒchēpiào fàng jìn bāo li le. Dìdi bǎ guǒzhī hē wán le, yòu bǎ wánjù shōu hǎo le. Nǎinai dǎ diànhuà shuō : « Nǐmen bǎ mén suǒ hǎo, lùshang xiǎoxīn ! »}\\
« Demain, toute notre famille va à Douala voir grand-mère. Le soir, papa a mis les bagages dans la voiture. Maman a préparé le repas et a aussi lavé les vêtements que les enfants porteront demain. Ma grande sœur a mis les billets de train dans le sac. Mon petit frère a fini son jus et a rangé ses jouets. Grand-mère a téléphoné : "Fermez bien la porte à clé et faites attention sur la route !" »
\end{textebox}

\textbf{Questions de compréhension} (réponds en français) :
\begin{enumerate}
\item Où la famille va-t-elle demain et pourquoi ?
\item Cite trois choses que les membres de la famille ont préparées le soir.
\item Que demande la grand-mère au téléphone ?
\end{enumerate}

\textbf{Questions de langue} :
\begin{enumerate}
\setcounter{enumi}{3}
\item Releve dans le texte quatre phrases en \zh{把} et recopie-les.
\item Pour chacune, indique le type de complément qui suit le verbe (\zh{在} / \zh{到} / \zh{给} / résultat / \zh{完}).
\item Transforme en phrase en \zh{把} : \zh{弟弟收好了玩具。}
\end{enumerate}
\end{exercice}

\begin{exercice}[Type Bac : thème et version]
\textbf{A. Thème.} Traduis en chinois (caractères + pinyin) les phrases suivantes, en utilisant la structure en \zh{把} chaque fois que c'est nécessaire :
\begin{enumerate}
\item J'ai donné mon dictionnaire à mon ami.
\item Papa a réparé le vélo.
\item N'oublie pas de fermer la porte à clé.
\item Il a mis les livres sur la table.
\end{enumerate}

\textbf{B. Version.} Traduis en français le paragraphe suivant :

\begin{textebox}[Version]
我把行李放到车上，妈妈把饭做好了，弟弟把果汁喝完了。我们可以出发了！\\
\emph{Wǒ bǎ xíngli fàng dào chē shang, māma bǎ fàn zuò hǎo le, dìdi bǎ guǒzhī hē wán le. Wǒmen kěyǐ chūfā le !}
\end{textebox}
\end{exercice}

\begin{exercice}[Type Bac : transformation, correction et expression écrite]
\textbf{A. Transformation.} Réécris les cinq phrases suivantes en phrases en \zh{把} quand c'est possible ; si la transformation est impossible, écris « impossible » et justifie en une phrase.
\begin{enumerate}
\item \zh{他洗干净了衣服。}
\item \zh{我喜欢这部电影。}
\item \zh{她写完了作业。}
\item \zh{他有一个弟弟。}
\item \zh{我们打开了窗户。}
\end{enumerate}

\textbf{B. Correction.} Chacune des phrases suivantes contient une faute. Identifie la faute, corrige la phrase et explique la règle en français.
\begin{enumerate}
\item \zh{我把一杯水喝完了。}
\item \zh{他把门没关好。}
\item \zh{我把书放。}
\end{enumerate}

\textbf{C. Expression écrite.} Rédige un court texte d'environ \textbf{80 caractères chinois} sur le sujet : « Le matin, avant de partir à l'école ». Tu dois employer \textbf{au moins quatre phrases en 把} variées : une avec \zh{在}, une avec \zh{到}, une avec \zh{给} et une avec un complément de résultat (\zh{完}, \zh{好}, \zh{干净}…). Donne le pinyin de ton texte.
\end{exercice}

\newpage

\coursec{Corrigés des exercices}

\begin{corrige}[Exercice 1 — QCM : la bonne structure]
\textbf{1. b} — \zh{我把饭吃完了。} \emph{Wǒ bǎ fàn chī wán le.} « J'ai fini de manger le riz. » La phrase a. est fautive : \zh{吃} (manger) n'est pas un verbe de placement, on ne peut pas dire \zh{*我把饭吃在食堂里} ; la phrase c. place \zh{把} après le verbe, ce qui est impossible : \zh{把} + objet se place toujours \textbf{avant} le verbe.

\textbf{2. b} — \zh{他没把门关好。} \emph{Tā méi bǎ mén guān hǎo.} « Il n'a pas bien fermé la porte. » Règle : la négation \zh{没（有）} / \zh{别} se place \textbf{avant} \zh{把}, jamais entre \zh{把} et le verbe.

\textbf{3. a} — \zh{我把书放在桌子上。} \emph{Wǒ bǎ shū fàng zài zhuōzi shàng.} « J'ai posé le livre sur la table. » Structure : Sujet + \zh{把} + Objet + Verbe + \zh{在} + Lieu. (Note : \zh{上} \emph{shàng} prend le 4\ieme{} ton après \zh{桌子}).

\textbf{4. b} — \zh{请把这杯水喝了。} \emph{Qǐng bǎ zhè bēi shuǐ hē le.} « Bois ce verre d'eau, s'il te plaît. » L'objet après \zh{把} doit être \textbf{déterminé} : \zh{这杯水} (ce verre d'eau) et non \zh{一杯水} (un verre d'eau, indéfini). De plus, le verbe ne peut pas rester « nu » : il faut au minimum un complément (ici \zh{了} marquant l'accomplissement).
\end{corrige}

\begin{corrige}[Exercice 2 — Vrai ou faux ? Justifie]
\begin{enumerate}
\item \textbf{Vrai.} Dans la phrase en \zh{把}, l'objet est déplacé avant le verbe : Sujet + \zh{把} + Objet + Verbe + Complément.
\item \textbf{Faux.} Le verbe d'une phrase en \zh{把} ne peut pas être « nu » : il doit être suivi d'un complément (résultat, \zh{了}, \zh{在}/\zh{到}/\zh{给}…). On dira : \zh{我把书看完了。}
\item \textbf{Vrai.} La négation se place avant \zh{把} : \zh{他没把门关好。} (et non \zh{*他把门没关好}).
\item \textbf{Faux.} La structure en \zh{把} exprime la \emph{disposition} d'un objet déterminé (on agit sur lui avec un résultat). Les verbes d'état ou de sentiment (\zh{喜欢}, \zh{有}, \zh{是}, \zh{知道}…) ne s'y transforment pas.
\item \textbf{Vrai.} L'objet de \zh{把} doit être connu ou précisé : \zh{这本书}, \zh{那杯水}, \zh{我的词典} ; jamais un objet indéfini comme \zh{一本书}.
\end{enumerate}
\end{corrige}

\begin{corrige}[Exercice 3 — Vocabulaire : complète le tableau]
\begin{center}
\begin{tabularx}{\linewidth}{|X|X|X|X|}
\hline
\rowcolor{popblueL} Caractères & Pinyin & Nature & Traduction \\
\hline
关 & guān & verbe & fermer \\
\hline
洗 & xǐ & verbe & laver \\
\hline
放 & fàng & verbe & poser, mettre \\
\hline
锁 & suǒ & verbe / nom & fermer à clé ; serrure \\
\hline
修 & xiū & verbe & réparer \\
\hline
行李 & xíngli & nom & bagages \\
\hline
\end{tabularx}
\end{center}
Points de vigilance : \zh{行} se lit ici \emph{xíng} (2\ieme{} ton, « marcher/aller ») et non \emph{háng} (2\ieme{} ton, « rang/ligne ») ; \zh{锁} a la clé du métal \zh{钅} (radical « métal »), comme beaucoup de mots liés au métal (\zh{针}, \zh{钉}, \zh{铁}…).
\end{corrige}

\begin{corrige}[Exercice 4 — Associe caractères et pinyin]
\begin{enumerate}
\item \zh{把} → \textbf{b. bǎ}
\item \zh{完} → \textbf{c. wán}
\item \zh{放} → \textbf{a. fàng}
\item \zh{干净} → \textbf{f. gānjìng}
\end{enumerate}
Les intrus sont \textbf{d. gǎnjìng} (mauvais ton sur la première syllabe : \zh{干} « sec/propre » est au 1\ier{} ton \emph{gān}, pas \emph{gǎn}) ; \textbf{e. guān} (c'est le pinyin de \zh{关}, absent de la colonne de gauche) ; \textbf{g. wǎn} (mauvais ton : \zh{完} est au 2\ieme{} ton \emph{wán}, pas \emph{wǎn}).
\end{corrige}

\begin{corrige}[Exercice 5 — Traduis en français]
\begin{enumerate}
\item \zh{妈妈把饭做好了。} — « Maman a préparé le repas. » (\zh{做好} : complément de résultat « faire → prêt ».)
\item \zh{弟弟把果汁喝完了。} — « Mon petit frère a fini son jus de fruit. » (\zh{喝完} : boire jusqu'au bout.)
\item \zh{老师把书放在桌子上了。} — « Le professeur a posé le livre sur la table. » (Verbe + \zh{在} + lieu.)
\item \zh{我把作业交给老师了。} — « J'ai remis mes devoirs au professeur. » (Verbe + \zh{给} + personne.)
\item \zh{请把窗户打开。} — « Ouvre la fenêtre, s'il te plaît. » (\zh{打开} : ouvrir ; phrase impérative, \zh{了} final non obligatoire.)
\end{enumerate}
\end{corrige}

\begin{corrige}[Exercice 6 — Transforme en phrase en 把]
\begin{enumerate}
\item \zh{他把衣服洗干净了。} \emph{Tā bǎ yīfu xǐ gānjìng le.} « Il a lavé les vêtements (ils sont propres). »
\item \zh{我把这本书看完了。} \emph{Wǒ bǎ zhè běn shū kàn wán le.} « J'ai fini de lire ce livre. »
\item \zh{她把门关上了。} \emph{Tā bǎ mén guān shàng le.} « Elle a fermé la porte. » (On ajoute le complément directionnel \zh{上} : le verbe ne peut pas rester suivi seulement de \zh{了} dans ce cas ; \zh{她把门关好了} « bien fermé » est aussi correct. \zh{*她把门关了} est à éviter à l'écrit formel.)
\item \zh{我们把饭吃完了。} \emph{Wǒmen bǎ fàn chī wán le.} « Nous avons fini de manger. »
\item \zh{他把自行车修好了。} \emph{Tā bǎ zìxíngchē xiū hǎo le.} « Il a réparé le vélo. »
\end{enumerate}
Méthode : Sujet + \zh{把} + Objet (déterminé) + Verbe + Complément (résultat / direction / lieu / \zh{给}) + \zh{了} (si action accomplie).
\end{corrige}

\begin{corrige}[Exercice 7 — Texte à trous]
Texte complété :

\zh{早上，小明把书包收拾好，然后把早饭吃完了。他把钥匙放在口袋里，把自行车骑到学校。到了教室，他把作业交给老师了。}

\emph{Zǎoshang, Xiǎo Míng bǎ shūbāo shōushi hǎo, ránhòu bǎ zǎofàn chī wán le. Tā bǎ yàoshi fàng zài kǒudài lǐ, bǎ zìxíngchē qí dào xuéxiào. Dào le jiàoshì, tā bǎ zuòyè jiāo gěi lǎoshī le.}

Justifications : \zh{收拾好} (résultat \zh{好}) ; \zh{吃完了} (résultat \zh{完} + \zh{了}) ; \zh{放在口袋里} (verbe + \zh{在} + lieu) ; \zh{骑到学校} (verbe + \zh{到} + destination) ; \zh{交给老师了} (verbe + \zh{给} + personne + \zh{了}). Chaque trou correspond à l'un des quatre grands types de compléments de la phrase en \zh{把} (résultat, lieu \zh{在}, direction \zh{到}, destinataire \zh{给}).
\end{corrige}

\begin{corrige}[Exercice 8 — Remets les mots dans l'ordre]
\begin{enumerate}
\item \zh{我把水喝完了。} \emph{Wǒ bǎ shuǐ hē wán le.} « J'ai fini l'eau. »
\item \zh{请你把门关好。} \emph{Qǐng nǐ bǎ mén guān hǎo.} « Ferme bien la porte, s'il te plaît. »
\item \zh{他把书放在桌子上。} \emph{Tā bǎ shū fàng zài zhuōzi shàng.} « Il a posé le livre sur la table. »
\item \zh{我把词典借给朋友了。} \emph{Wǒ bǎ cídiǎn jiè gěi péngyou le.} « J'ai prêté le dictionnaire à mon ami. »
\item \zh{她没把窗户打开。} \emph{Tā méi bǎ chuānghu dǎkāi.} « Elle n'a pas ouvert la fenêtre. » (La négation \zh{没} précède \zh{把} ; pas de \zh{了} à la forme négative passée.)
\end{enumerate}
\end{corrige}

\begin{corrige}[Exercice 9 — Type Bac : compréhension et langue]
\textbf{Barème indicatif (10 pts)} : compréhension 3 pts (1 pt par question) ; langue 7 pts (questions 4 et 5 : 4 pts ; question 6 : 3 pts).

\textbf{Compréhension}
\begin{enumerate}
\item La famille va à \textbf{Douala} (\zh{杜阿拉} \emph{Dù'ālā}) demain pour \textbf{rendre visite à la grand-mère} (\zh{看奶奶} \emph{kàn nǎinai}).
\item Trois préparatifs possibles parmi : papa a mis les bagages dans la voiture (\zh{爸爸把行李放到车上了}) ; maman a préparé le repas (\zh{妈妈把饭做好了}) ; maman a lavé les vêtements des enfants (\zh{妈妈把衣服洗干净了}) ; la grande sœur a mis les billets de train dans le sac (\zh{姐姐把火车票放进包里了}) ; le petit frère a rangé ses jouets (\zh{弟弟把玩具收好了}).
\item Elle demande de \textbf{bien fermer la porte à clé} (\zh{把门锁好} \emph{bǎ mén suǒ hǎo}) et de \textbf{faire attention sur la route} (\zh{路上小心} \emph{lùshang xiǎoxīn}).
\end{enumerate}

\textbf{Langue}
\begin{enumerate}
\setcounter{enumi}{3}
\item Quatre phrases en \zh{把} (au choix) : \zh{爸爸把行李放到车上了。} / \zh{妈妈把饭做好了。} / \zh{姐姐把火车票放进包里了。} / \zh{弟弟把果汁喝完了。} / \zh{你们把门锁好。}
\item Types de compléments : \zh{放到车上} → verbe + \zh{到} + lieu (direction) ; \zh{做好} → complément de résultat (\zh{好}) ; \zh{放进包里} → verbe + directionnel \zh{进} + lieu ; \zh{喝完} → complément de résultat (\zh{完}) ; \zh{锁好} → complément de résultat (\zh{好}).
\item \zh{弟弟把玩具收好了。} \emph{Dìdi bǎ wánjù shōu hǎo le.} « Mon petit frère a rangé les jouets. »
\end{enumerate}
Points de vigilance : ne pas oublier \zh{了} en fin de phrase pour marquer l'accomplissement ; l'objet après \zh{把} est toujours déterminé (\zh{行李}, \zh{火车票}, \zh{门} sont connus des interlocuteurs).
\end{corrige}

\begin{corrige}[Exercice 10 — Type Bac : thème et version]
\textbf{Barème indicatif (10 pts)} : thème 6 pts (1,5 pt par phrase) ; version 4 pts.

\textbf{A. Thème (Français → Chinois)}
\begin{enumerate}
\item \zh{我把词典给朋友了。} \emph{Wǒ bǎ cídiǎn gěi péngyou le.} (Ou : \zh{我把词典送给朋友了。} si « donner » = offrir.)
\item \zh{爸爸把自行车修好了。} \emph{Bàba bǎ zìxíngchē xiū hǎo le.}
\item \zh{别忘了把门锁好。} \emph{Bié wàng le bǎ mén suǒ hǎo.} (La négation \zh{别} précède tout le groupe verbal \zh{忘了} ; la proposition en \zh{把} suit naturellement.)
\item \zh{他把书放在桌子上了。} \emph{Tā bǎ shū fàng zài zhuōzi shàng le.}
\end{enumerate}

\textbf{B. Version (Chinois → Français)}
Texte source : \zh{我把行李放在车上了，妈妈把饭做好了，弟弟把果汁喝完了。我们可以出发了！}
\emph{Wǒ bǎ xíngli fàng zài chē shàng le, māma bǎ fàn zuò hǎo le, dìdi bǎ guǒzhī hē wán le. Wǒmen kěyǐ chūfā le!}

Traduction : « J'ai mis les bagages dans la voiture, maman a préparé le repas, et mon petit frère a fini son jus de fruit. Nous pouvons partir ! »

Points de vigilance : en A1, le verbe \zh{给} seul suffit (pas de double \zh{给} : on ne dit pas \zh{*把词典给给朋友}) ; en A3, ne pas traduire « n'oublie pas » par \zh{*不要忘} au début sans \zh{了} : la forme naturelle est \zh{别忘了} ; dans la version, \zh{出发} = « partir, se mettre en route ».
\end{corrige}

\begin{corrige}[Exercice 11 — Type Bac : transformation, correction et expression écrite]
\textbf{Barème indicatif (20 pts)} : A. 5 pts (1 pt par phrase) ; B. 6 pts (2 pts par phrase : correction 1 pt + explication 1 pt) ; C. 9 pts (respect des quatre structures 4 pts, correction grammaticale 3 pts, longueur et cohérence 2 pts).

\textbf{A. Transformation}
\begin{enumerate}
\item \zh{他把衣服洗干净了。} \emph{Tā bǎ yīfu xǐ gānjìng le.}
\item \textbf{Impossible} : \zh{喜欢} est un verbe de sentiment, sans action ni résultat sur l'objet ; la phrase en \zh{把} exige une disposition concrète de l'objet.
\item \zh{她把作业写完了。} \emph{Tā bǎ zuòyè xiě wán le.}
\item \textbf{Impossible} : \zh{有} est un verbe de possession (état), incompatible avec la structure en \zh{把}.
\item \zh{我们把窗户打开了。} \emph{Wǒmen bǎ chuānghu dǎkāi le.}
\end{enumerate}

\textbf{B. Correction}
\begin{enumerate}
\item Faute : objet indéfini. Correction : \zh{我把这杯水喝完了。} \emph{Wǒ bǎ zhè bēi shuǐ hē wán le.} Règle : l'objet de \zh{把} doit être déterminé (démonstratif, possessif, ou connu du contexte) ; \zh{一杯水} est indéfini.
\item Faute : place de la négation. Correction : \zh{他没把门关好。} \emph{Tā méi bǎ mén guān hǎo.} Règle : la négation (\zh{没}, \zh{别}) se place \textbf{avant} \zh{把}, jamais entre \zh{把} et le verbe.
\item Faute : verbe « nu ». Correction : \zh{我把书放在桌子上。} ou \zh{我把书放好了。} Règle : dans une phrase en \zh{把}, le verbe doit être accompagné d'un complément (résultat, lieu avec \zh{在}, \zh{了}…).
\end{enumerate}

\textbf{C. Expression écrite — proposition de rédaction modèle}

\zh{早上，我七点起床。我先把被子叠好，再把书包收拾好。妈妈把早饭做好了，我把牛奶喝完，把面包吃完。然后我把钥匙放在口袋里，把自行车骑到学校。到了教室，我把作业交给老师。}

\emph{Zǎoshang, wǒ qī diǎn qǐchuáng. Wǒ xiān bǎ bèizi dié hǎo, zài bǎ shūbāo shōushi hǎo. Māma bǎ zǎofàn zuò hǎo le, wǒ bǎ niúnǎi hē wán, bǎ miànbāo chī wán. Ránhòu wǒ bǎ yàoshi fàng zài kǒudài lǐ, bǎ zìxíngchē qí dào xuéxiào. Dào le jiàoshì, wǒ bǎ zuòyè jiāo gěi lǎoshī.}

« Le matin, je me lève à sept heures. Je plie d'abord ma couette, puis je range mon cartable. Maman a préparé le petit-déjeuner ; je finis mon lait et mon pain. Ensuite je mets mes clés dans ma poche et je vais à l'école à vélo. Arrivé en classe, je remets mes devoirs au professeur. »

Vérification des quatre structures exigées : \zh{放在口袋里} (verbe + \zh{在}) ; \zh{骑到学校} (verbe + \zh{到}) ; \zh{交给老师} (verbe + \zh{给}) ; \zh{叠好} / \zh{收拾好} / \zh{喝完} / \zh{吃完} (compléments de résultat). Environ 80 caractères.

Points de vigilance : objets toujours déterminés ; \zh{了} final pour les actions accomplies ; ne pas placer d'adverbe de temps entre \zh{把} et le verbe (\zh{先} et \zh{再} se mettent \textbf{avant} \zh{把}).
\end{corrige}

\newpage

\coursec{Fiche bilan}

\begin{popfigure}
\begin{tikzpicture}[
    mindmap,
    grow cyclic,
    every node/.style={concept, execute at begin node=\hskip0pt},
    concept color=popblue,
    level 1/.append style={level distance=4.5cm, sibling angle=360/7},
    level 2/.append style={level distance=3cm, sibling angle=360/6},
    texte/.style={concept, font=\small\bfseries, text width=2.8cm, align=center, inner sep=4pt},
    branche1/.style={concept color=poporange},
    branche2/.style={concept color=popgreen},
    branche3/.style={concept color=poppurple},
    branche4/.style={concept color=poppink},
    branche5/.style={concept color=popgold},
    branche6/.style={concept color=popblue},
    branche7/.style={concept color=popdark},
]
\node[texte, font=\large\bfseries, text width=3cm] {把字句\\ (Bǎzìjù)}
    child[branche1] node[texte, align=center]{Structure\\ \cle{S + 把 + O + V + Compl.}}
    child[branche2] node[texte, align=center]{Emplois\\ \cle{Action \emph{aboutie} / \emph{résultat} / \emph{déplacement}}}
    child[branche3] node[texte, align=center]{Compléments\\ \cle{Obligatoires : 了 / 完 / 好 / 走 / 到...}}
    child[branche4] node[texte, align=center]{Négation\\ \cle{没(有) + V (+ 了) / 不 + V}}
    child[branche5] node[texte, align=center]{Comparaison FR\\ \cle{Passif / Causatif / COD antéposé}}
    child[branche6] node[texte, align=center]{Pièges FR\\ \cle{Ordre / 了 vs 过 / Compl. manquant}}
    child[branche7] node[texte, align=center]{Exemples clés\\ \cle{把书放好 / 把门关上}};
\end{tikzpicture}
\legende{Carte mentale récapitulative : la structure \emph{bǎ} (disposition / préposition d'objet).}
\end{popfigure}

\begin{bilan}
\end{bilan}

\courssub{Lexique clé (Mots-outils \& Verbes fréquents)}
\begin{center}
\begin{tabularx}{\linewidth}{|X|X|X|X|}\hline
\rowcolor{popblueL} \textbf{Caractères} & \textbf{Pinyin} & \textbf{Nature} & \textbf{Traduction} \\ \hline
把 & bǎ & Préposition / Verbe & Prendre, saisir / (marque l'objet disposé) \\ \hline
了 & le & Particule & Aspect accompli / changement d'état \\ \hline
完 & wán & Compl. de résultat & Finir, terminer (action) \\ \hline
好 & hǎo & Compl. de résultat & Bien faire, achever correctement \\ \hline
上 & shàng & Compl. de direction/résultat & Monter / Fermer / Allumer / Commencer \\ \hline
下 & xià & Compl. de direction/résultat & Descendre / Éteindre / Arrêter \\ \hline
走 & zǒu & Compl. de direction & Partir, emmener \\ \hline
到 & dào & Compl. de direction/résultat & Arriver à / Atteindre \\ \hline
给 & gěi & Préposition / Verbe & Donner / À (bénéficiaire) / Passif \\ \hline
对 & duì & Préposition & Vers, à (direction interaction) \\ \hline
\end{tabularx}
\end{center}

\courssub{Règles de grammaire à connaître par cœur}
\begin{center}
\begin{tabularx}{\linewidth}{|p{3.5cm}|X|X|}\hline
\rowcolor{popblueL} \textbf{Notion} & \textbf{Schéma / Règle} & \textbf{Exemple (Car. + Pin. + Trad.)} \\ \hline
\cle{Structure canonique} & Sujet + \textbf{把} + Objet + Verbe + \textbf{Complément obligatoire} & \zh{我把作业做完了。} (Wǒ bǎ zuòyè zuò wán le.) — J'ai fini mes devoirs. \\ \hline
\cle{Nature de l'objet} & Objet \textbf{défini, connu, spécifique} (souvent avec 的, 这, 那, mon/ton/son). Interdit : objet indéfini (un livre). & \checkmark \zh{把这本书还给他。} \quad \cross \zh{把一本书还给他。} \\ \hline
\cle{Verbes admis} & Verbes d'action \textbf{transitifs} à résultat/déplacement : 放, 关, 开, 吃, 喝, 做, 买, 卖, 还, 借, 送, 带, 拿, 洗, 写... & \zh{请把门关上。} (Qǐng bǎ mén guān shàng.) — Fermez la porte. \\ \hline
\cle{Complément OBLIGATOIRE} & Le verbe ne peut pas être nu. Il faut : \textbf{了 / 完 / 好 / 走 / 到 / 上 / 下 / 给 / 成 / 为...} & \cross \zh{我把苹果吃。} \quad \checkmark \zh{我把苹果吃\cle{了}。} \\ \hline
\cle{Négation : \textbf{没(有)}} & \textbf{没(有)} + V (+ 了) \rightarrow Action \textbf{non réalisée} (passé). \textbf{不} + V \rightarrow Habitude / Volonté / Futur. & \zh{我\cle{没}把票买\cle{了}。} (Pas acheté.) \quad \zh{我\cle{不}把票买\cle{了}。} (Je ne l'achèterai pas / Je refuse.) \\ \hline
\cle{Objet + 给/到 + Destinataire} & S + 把 + O + \textbf{给/到} + Pers./Lieu + V + Compl. & \zh{把钥匙给我。} (Donne-moi la clé.) \quad \zh{把椅子搬到屋里。} (Déplace la chaise dans la pièce.) \\ \hline
\cle{把 vs 被} & \textbf{把} : Sujet = \textbf{Agent} (actif, intentionnel). \quad \textbf{被} : Sujet = \textbf{Patient} (passif, subi). & \zh{他把书卖了。} (Il a vendu le livre.) \quad \zh{书被他卖了。} (Le livre a été vendu par lui.) \\ \hline
\end{tabularx}
\end{center}

\courssub{Méthodes express}
\begin{enumerate}
    \item \textbf{Identifier l'objet manipulé} : Est-il précis ? (mon téléphone, ce devoir, la porte). Si vague (un téléphone) $\rightarrow$ pas de \cle{把}.
    \item \textbf{Vérifier le verbe} : Action concrète, transitive, qui change l'état/la place de l'objet.
    \item \textbf{Ajouter le complément} : \cle{了} (fini), \cle{完} (terminé), \cle{好} (bien fait), \cle{走/到+lieu} (déplacé), \cle{给+pers.} (donné).
    \item \textbf{Placer la négation} : Devant \cle{把} (\cle{不/没} + 把...). \textit{Attention :} \cle{没} + V + \cle{了} (passé non fait) vs \cle{不} + V + \cle{了} (refus/futur non fait).
    \item \textbf{Traduire en français} : Souvent voix active normale (« J'ai mangé la pomme »), parfois causatif (« Fais sortir le chat » $\rightarrow$ \zh{把猫放出去}), ou passif si contexte l'exige.
\end{enumerate}


\begin{aretenir}[Checklist avant le Bac — \cle{把字句}]
\begin{itemize}
    \item $\square$ Je sais écrire l'ordre des traits de \zh{把} (扌 + 巴 : main + s'accrocher) et je connais son radical \cle{扌} (main).
    \item $\square$ Je récite le schéma : \cle{S + 把 + O (défini) + V + Complément (obligatoire)}.
    \item $\square$ Je distingue \cle{把} (actif / je gère l'objet) et \cle{被} (passif / l'objet me subit).
    \item $\square$ Je ne mets \textbf{jamais} d'objet indéfini après \cle{把} (pas de \zh{一本}, \zh{一些} sans déterminant).
    \item $\square$ Je sais que le verbe \textbf{ne peut pas rester seul} : il faut \cle{了, 完, 好, 上, 下, 走, 到, 给...}.
    \item $\square$ Je place la négation \textbf{devant} \cle{把} : \cle{没(有)把...V...了} (passé) / \cle{不把...V...} (habitude/refus).
    \item $\square$ Je ne confonds pas \cle{了} (accompli) et \cle{过} (expérience) : \cross \zh{把饭吃过了} $\rightarrow$ \checkmark \zh{把饭吃了} / \zh{吃过饭}.
    \item $\square$ Je gère la structure \cle{把 + O + 给/到 + Dest. + V + Compl.} (ex: \zh{把作业交给老师}).
    \item $\square$ Je traduis correctement en français : voix active standard, ou impératif, ou « faire qqch à qqch ».
    \item $\square$ Je repère les pièges : \cle{把} + adjectif (impossible), \cle{把} + verbe psychologique (\zh{喜欢}, \zh{知道}) impossible.
    \item $\square$ Je sais produire 3 phrases types : action courante (ménage), déplacement (rangement), transmission (donner).
\end{itemize}
\end{aretenir}

\findecours

\end{document}
